\documentclass[11pt,a4paper]{article}

\usepackage[utf8]{inputenc}
\usepackage[T1]{fontenc}
\usepackage{lmodern}
\usepackage{geometry}
\usepackage{setspace}
\usepackage{amsmath,amssymb,amsthm}
\usepackage{booktabs}
\usepackage{longtable}
\usepackage{pdflscape}
\usepackage{array}
\usepackage{enumitem}
\usepackage{natbib}
\usepackage{hyperref}
\usepackage{xcolor}
\usepackage{caption}
\usepackage{float}
\usepackage{authblk}
\usepackage{url}

\geometry{margin=1in}
\onehalfspacing
\hypersetup{
    colorlinks=true,
    linkcolor=blue!50!black,
    citecolor=blue!50!black,
    urlcolor=blue!50!black
}
\setlist[itemize]{leftmargin=*}
\setlist[enumerate]{leftmargin=*}
\setlength{\parskip}{0.35em}

\newtheorem{definition}{Definition}
\newtheorem{proposition}{Proposition}
\newtheorem{assumption}{Assumption}
\newtheorem{lemma}{Lemma}

\newcommand{\GDP}{Y}
\newcommand{\FS}{FS}
\newcommand{\MFC}{MFC}
\newcommand{\LS}{LS}
\newcommand{\AI}{AI}
\newcommand{\Prob}{\Pr}

\title{AI, Fiscal Space, and Universal Social Protection in Emerging Economies:\\[0.2em]A Robust Accounting-Fiscal Threshold Framework}
\author{Vicenzo Scavino}
\affil{Preliminary methodological draft}
\date{Version 3.0 --- Empirical version --- July 9, 2026}

\begin{document}
\maketitle

\begin{abstract}
This paper does not forecast whether artificial intelligence (AI) will finance a universal basic income (UBI). It develops a calibrated threshold-accounting framework to identify the threshold combinations of AI-induced productivity, productive adoption, fiscal capture, informality reduction, labor-share allocation, and policy cost under which universal or semi-universal social protection instruments become accounting-fiscally feasible in emerging economies. The framework compares universal social pensions, minimum universal transfers, guaranteed minimum income, partial basic income, and a full UBI benchmark. Its methodological contribution is a set of transparent equations for constructing AI shocks, converting those shocks into tax-base channels, adjusting for institutional frictions without double counting informality, and deriving policy-specific thresholds for growth, adoption, fiscal capture, and required persistence horizons. The paper separates observed variables from literature-disciplined scenario parameters and from unobserved structural parameters that require sensitivity analysis or Monte Carlo simulation. It also introduces impossibility diagnostics, temporal placebos, sectoral negative controls, leave-one-source-out ablations, Monte Carlo feasibility distributions, and global sensitivity diagnostics. The value of the framework is therefore not a point forecast but a falsifiable and reusable frontier that can be recalibrated as better evidence on AI productivity, adoption, labor-share effects, and fiscal capture becomes available.
\end{abstract}

\noindent \textbf{Keywords:} artificial intelligence, fiscal space, universal social protection, universal basic income, informality, fiscal capacity, labor share, emerging economies, Latin America.\\
\textbf{JEL codes:} O33, H53, H24, I38, O54, E62.

\vspace{0.75em}
\noindent\textbf{Status of this version.} Empirical version. All results generated by the tagged reproducibility package (tags fase-B and etapa-7.1); primary specification frozen under pre-registration hash prior to any official run; the horizon-inversion statistic H* was added post-baseline as a derived reporting statistic (declared amendment; no parameter or pre-registered result was altered).

\newpage

\section{Introduction}

The rapid diffusion of artificial intelligence (AI), especially generative AI, has reopened a classic question in political economy: if technology raises productivity while weakening the employment-income link for some workers, can social protection be redesigned so that technological gains are broadly shared? In public debate, this question is often reduced to whether AI will finance a universal basic income. That version is too strong. It mixes an uncertain technological forecast, a fiscal conclusion, and a specific policy instrument.

This paper proposes a narrower and more defensible research question: under what combinations of AI-induced growth, adoption, exposure, labor-share pressure, informality, and fiscal capacity do universal or semi-universal social protection instruments become baseline accounting-fiscally feasible in emerging economies? The aim is not to predict that AI will finance a full UBI. The aim is to build a threshold methodology that shows which instruments become feasible first under explicit assumptions.

The distinction matters for three reasons. First, macroeconomic estimates of AI productivity gains remain uncertain. A conservative task-based approach implies modest aggregate effects when the affected task share or cost savings are limited \citep{acemoglu2024simple}. A micro-to-macro approach by the OECD estimates annual aggregate total-factor productivity gains from AI of about 0.25 to 0.6 percentage points over a ten-year horizon under its main assumptions \citep{filippucci2024miracle}. Second, exposure is not adoption. Countries with high occupational exposure may not obtain high productivity gains if they lack digital infrastructure, human capital, firm capabilities, or regulatory readiness. The IMF AI Preparedness Index explicitly combines digital infrastructure, innovation and economic integration, human capital and labor market policies, and regulation and ethics \citep{cazzaniga2024genai}. Third, even if AI raises output, fiscal capacity determines whether governments can convert output gains into usable revenue. High informality, limited tax capacity, and weak tax systems can block the conversion of technological gains into social protection \citep{gaspar2016tax}.

The central object is a country-policy-scenario feasibility condition. For country $i$, policy instrument $p$, fiscal regime $r$, and AI scenario $s$, the framework asks whether effective additional fiscal space is at least as large as the annual cost of the policy:
\begin{equation}
    fs^{eff}_{i,p,s,r,t} \geq c_{i,p,t},
\end{equation}
where all terms are expressed as shares of the no-AI counterfactual GDP $Y^0_{i,t}$ defined in Section \ref{sec:policy_costs}. Throughout, the unsuperscripted cost $c_{i,p,t}$ denotes the gross-cost ratio $c^{gross}_{i,p,t}$ of Section \ref{sec:policy_costs}, consistent with the convention that $V$ without superscript denotes $V^{gross}$. The associated feasibility index is
\begin{equation}
    V_{i,p,s,r,t}=\frac{fs^{eff}_{i,p,s,r,t}}{c_{i,p,t}}.
\end{equation}
The policy crosses the structural accounting threshold when $V_{i,p,s,r,t}\geq 1$.

The contribution is methodological. The paper develops: (i) a disciplined construction of AI shocks from exposure, preparedness, adoption, and base productivity scenarios; (ii) a fiscal conversion block that separates labor-, consumption-, and capital/profit-base capture, AI-rent capture, and administrative and transition frictions; (iii) a labor-share block that distinguishes productivity gains from their distribution between labor and capital; (iv) explicit thresholds for growth, adoption, and fiscal capture; and (v) a robustness protocol that prevents the results from depending on arbitrary single-number assumptions.

\section{Related Literature and Positioning}\label{sec:literature}

This paper connects five literatures that are usually studied separately.

\subsection{AI, productivity, and macroeconomic growth}

A first literature studies the macroeconomic productivity effects of AI \citep{oecd2024impactai}. \citet{acemoglu2024simple} uses a task-based accounting logic in which aggregate gains depend on the share of tasks affected by AI and the cost savings achieved on those tasks. This provides a conservative benchmark for the lower tail of AI-induced growth scenarios. \citet{filippucci2024miracle} build a micro-to-macro framework that combines micro-level performance gains, AI exposure, adoption rates, and input-output linkages; their main estimates imply annual aggregate TFP gains from AI in the range of 0.25 to 0.6 percentage points over a ten-year horizon.

This paper does not choose one point estimate. Instead, it treats those studies as anchors for scenario design. Low scenarios should be close to conservative task-based estimates; central scenarios should be disciplined by micro-to-macro productivity evidence; high scenarios should be interpreted as stress tests, not forecasts.

\subsection{Exposure, adoption, and preparedness}

A second literature emphasizes that AI exposure is not equal to realized productivity. The IMF's work on AI and the future of work proposes an AI Preparedness Index covering four dimensions: digital infrastructure, innovation and economic integration, human capital and labor market policies, and regulation and ethics \citep{cazzaniga2024genai}. This matters because a country may have a large share of workers in exposed occupations but still fail to capture productivity gains if digital access, skills, or firm capabilities are weak.

For Latin America, the World Bank and ILO estimate that 26--38 percent of jobs could be exposed to generative AI, 8--14 percent could experience productivity-enhancing transformation, and 2--5 percent are at risk of full automation under current capabilities \citep{worldbankilo2024lacgenai}. They also emphasize that digital gaps may block a large fraction of the potential productivity dividend. This motivates the explicit adoption function introduced below.

\subsection{Automation, labor share, and distribution}

A third literature studies how automation affects the distribution of income between labor and capital. \citet{acemoglu2018race} and \citet{acemoglu2019automation} distinguish displacement effects, which move tasks away from labor, from reinstatement effects, which create new labor-intensive tasks. This distinction is crucial for AI. AI can raise productivity and simultaneously reduce the labor share if task displacement dominates task reinstatement.

The fiscal implication is direct. If AI gains accrue mainly to profits, intangible assets, or platform rents, existing labor-tax and payroll-tax systems may capture less revenue than a simple tax-to-GDP ratio suggests. Conversely, stronger capital/profit taxation or AI-rent capture can raise the marginal fiscal conversion of AI growth. Therefore, the model includes labor-share dynamics rather than treating GDP growth as distributionally neutral.

\subsection{AI and fiscal policy}

A fourth literature studies how fiscal policy should adapt to generative AI. \citet{brollo2024fiscal} argue that fiscal policy is central for broadening AI gains, including through social protection, education, labor-market policies, capital-income taxation, and corporate taxation. This paper uses that logic to define fiscal regimes rather than assuming a single tax system. Under the current regime, AI growth is converted into fiscal space through existing tax buoyancy. Under reform regimes, governments may capture a larger share of AI rents or capital income.

\subsection{Universal basic income and social protection}

A fifth literature studies UBI and social protection design. \citet{gentilini2020ubi} emphasize that UBI must be analyzed as a specific transfer architecture, not as a generic promise of technological abundance. \citet{oecd2017basicincome} show that replacing benefit systems with a basic income involves difficult fiscal and distributional tradeoffs. This paper therefore compares a menu of instruments instead of treating full UBI as the only relevant policy.

\section{Conceptual Mechanism}\label{sec:conceptual}

The core mechanism is
\begin{equation}
AI \rightarrow adoption \rightarrow productivity \rightarrow income distribution \rightarrow fiscal space \rightarrow social protection.
\end{equation}
The deliberately one-directional chain omits feedback from protection to adoption or productivity through human capital, entrepreneurship, or labor supply; this is a maintained static-frontier assumption, not an empirical claim. Five frictions intervene: limited digital preparedness, uneven exposure, labor-share displacement, informality, and weak fiscal capacity.

The paper's key conceptual distinction is between \emph{technical exposure} and \emph{fiscal feasibility}. A country can be highly exposed to AI but fiscally unable to finance broad social protection. Conversely, a country with moderate exposure but strong tax capacity and low informality may cross policy thresholds sooner.

\begin{definition}[Baseline accounting-fiscal feasibility]
For country $i$, policy $p$, scenario $s$, fiscal regime $r$, and year $t$, policy $p$ satisfies \emph{baseline accounting-fiscal feasibility} if effective AI-induced fiscal space as a share of no-AI counterfactual GDP $Y^0_{i,t}$ is greater than or equal to the annual recurring policy cost as a share of the same denominator:
\begin{equation}
    fs^{eff}_{i,p,s,r,t}\geq c_{i,p,t}.
\end{equation}
Here $fs^{eff}$ is the fiscal space available to policy $p$: AI-generated revenue after the modeled deductions, including the policy's own net administrative burden. That burden may be negative when the programs replaced by $p$ are administratively more costly than the new instrument; in that case, $fs^{eff}$ exceeds revenue induced purely by AI.
This condition means only that the modeled AI-related fiscal-space channel covers the annual accounting cost of the policy. It does not imply welfare optimality, political feasibility, administrative feasibility, the absence of general-equilibrium effects, or long-run fiscal sustainability by itself. This narrower interpretation is deliberate because the IMF defines fiscal space as a multidimensional concept: the ability to raise spending or lower taxes without endangering market access and debt sustainability \citep{imf2016fiscalspace}. The framework therefore treats accounting coverage as a first-stage screen, not as a full fiscal-space assessment.
\end{definition}

\begin{definition}[Feasibility index and structural threshold]
The baseline feasibility index is
\begin{equation}
    V_{i,p,s,r,t}=\frac{fs^{eff}_{i,p,s,r,t}}{c_{i,p,t}}.
\end{equation}
Because $c_{i,p,t}>0$, the only structural accounting threshold is
\begin{equation}
    V_{i,p,s,r,t}=1 \quad \Longleftrightarrow \quad fs^{eff}_{i,p,s,r,t}=c_{i,p,t}.
\end{equation}
Hence $V<1$ denotes an uncovered fiscal gap relative to the policy cost, $V=1$ denotes exact accounting coverage, and $V>1$ denotes a positive fiscal buffer over the annual policy cost. All other cutoffs used later are reporting bands or prudential buffers around this structural identity, not structural parameters.
\end{definition}


\section{Policy Instruments and Cost Construction}\label{sec:policy_costs}

Let $Y^0_{i,t}$ denote the no-AI counterfactual GDP level used as the fiscal-space denominator, $N_{i,p,t}$ the eligible population for policy $p$, and $B_{i,p,t}$ the annual benefit per eligible person. The framework distinguishes gross cost, net incremental cost, and GDP-normalized cost. The gross nominal cost is
\begin{equation}
    C^{gross}_{i,p,t}=B_{i,p,t}N_{i,p,t},
\end{equation}
and the baseline gross cost ratio is
\begin{equation}
    c^{gross,0}_{i,p,t}=\frac{C^{gross}_{i,p,t}}{Y^0_{i,t}}.
\end{equation}
Unless stated otherwise, $c^{gross}_{i,p,t}$ denotes $c^{gross,0}_{i,p,t}$. The post-AI-denominator version $c^{gross,AI}_{i,p,s,t}=C^{gross}_{i,p,t}/Y^{AI}_{i,s,t}$, using Section \ref{sec:ai_shocks}, is robustness only. With existing social-protection spending, define
\begin{equation}
    C^{net}_{i,p,t}=\max\{0,C^{gross}_{i,p,t}-C^{exist}_{i,p,t}\},
    \qquad
    c^{net,0}_{i,p,t}=\frac{C^{net}_{i,p,t}}{Y^0_{i,t}}.
\end{equation}
$C^{exist}$ includes only substitutable non-contributory spending; contributory pensions and acquired-right benefits cannot be redirected and are excluded.
The gross cost is not net of first-round fiscal recycling: transfer recipients spend part of the benefit on VAT-taxable consumption, so the net fiscal cost is below $c^{gross}$. The labelled first-round extension is
\begin{equation}\label{eq:first_round_benefit_recycling_e5}
    c^{recyc}_{i,p,r,t}=c^{gross}_{i,p,t}\left[1-mpc^{ben}_i(1-m^M_i)(1-exempt^C_{i,t})\tau^{C,eff}_{i,r,t}\right],
\end{equation}
where $mpc^{ben}$ is recipients' marginal propensity to consume. The baseline omits recycling, which biases feasibility conservatively; if the extension is used, it must be reported as a labelled variant and the same convention applied to all instruments. With economy-wide parameters, it is an approximate upper bound because recipient baskets are more exempt or informal and more domestic than average. Greater domestic content works oppositely through lower import leakage, so exemption and effective-rate margins drive that upper-bound reading. Equation \eqref{eq:first_round_benefit_recycling_e5} uses the base-side exemption convention; if exemptions instead enter $\tau^{C,eff}$, omit $(1-exempt^C_{i,t})$ under the assignment rule of Section \ref{sec:fiscal_conversion}.
For readability, $c^{net}_{i,p,t}$ denotes the no-AI-denominator net-cost ratio unless otherwise stated. The corresponding ratios are
\begin{equation}
    V^{gross}_{i,p,s,r,t}=\frac{fs^{eff}_{i,p,s,r,t}}{c^{gross}_{i,p,t}},
    \qquad
    V^{net}_{i,p,s,r,t}=\frac{fs^{eff}_{i,p,s,r,t}}{c^{net}_{i,p,t}}\quad\text{only if }c^{net}_{i,p,t}>0,
\end{equation}
with $V^{gross}$ used as the default reporting object. If $c^{net}=0$, the instrument is classified as already covered by existing spending on the net-cost margin and no net-cost ratio is computed.

The policy menu is intentionally broader than full UBI.

\begin{table}[H]
\centering
\caption{Social protection instruments and cost logic}
\label{tab:policy_rules}
\small
\begin{tabular}{p{0.24\textwidth}p{0.31\textwidth}p{0.31\textwidth}}
\toprule
Instrument & Eligible unit & Benefit benchmark \\
\midrule
Universal social pension & Residents above age threshold $a^*$ & Poverty-line share or minimum pension \\
Minimum universal transfer & Residents or adults & Low fraction of poverty line \\
Guaranteed minimum income & Households below minimum threshold $z$ & Fraction of household poverty gap \\
Partial basic income & Residents or adults & Intermediate fraction of poverty line \\
Full UBI benchmark & Whole population & Fraction of poverty line; values above the poverty line only as labelled stress \\
\bottomrule
\end{tabular}
\end{table}

The corresponding gross-cost rules are
\begin{align}
    C^{PEN,gross}_{i,t} &= B^{PEN}_{i,t}Pop_{i,t}^{a\geq a^*},\\
    C^{MUT,gross}_{i,t} &= \eta_{MUT}PL_{i,t}N^{MUT}_{i,t},\\
    C^{GMI,gross}_{i,t} &= \chi_{GMI}\sum_h w_{h,i,t}n_{h,i,t}\max\{0,z_{i,t}-y^{pc}_{h,i,t}\},\\
    C^{PBI,gross}_{i,t} &= \eta_{PBI}PL_{i,t}N^{PBI}_{i,t},\\
    C^{UBI,gross}_{i,t} &= \eta_{UBI}PL_{i,t}Pop_{i,t}.
\end{align}
Here $PL_{i,t}$ is the poverty-line benchmark; $z_{i,t}$ the per-capita minimum-income threshold; $y^{pc}_{h,i,t}$ household income or consumption per capita; $n_{h,i,t}$ household size; $w_{h,i,t}$ survey weight; and $\chi_{GMI}\in(0,1]$ the poverty-gap share closed. With household-level objects, the equivalent rule is $\chi_{GMI}\sum_h w_h\max\{0,z^H_h-y^H_h\}$. Units of $z$ and $y$ must match before computing GMI cost; benefit levels and poverty or minimum-income thresholds enter the ratios annually. Full UBI is a fiscal stress test, not the default policy.

The GMI equation prices the aggregate poverty gap under three idealizations: exact identification of poor households (no inclusion or exclusion errors), perfect tapering to each household's gap, and no income response. Full gap-filling implies a 100 percent marginal tax rate below $z$; households may reduce or underreport income, widening the gap. The labelled loading $\theta^{target}\geq1$ defines $C^{GMI,load}=\theta^{target}C^{GMI,gross}$ and must be reported alongside ideal targeting, because universal costs lack this idealization and the GMI-before-UBI ranking is partly mechanical.
The cost side holds recipient behavior fixed: labor-supply responses to transfers (small in the cash-transfer evidence, in either direction) and local price effects of large transfers in poor areas are outside the accounting frontier \citep{gentilini2020ubi,oecd2017basicincome}.

\section{Constructing AI Growth Shocks}\label{sec:ai_shocks}

The calibrated $g^{AI}_{i,s,t}$ is a scenario-based \emph{level} increment in the fiscal base relative to a no-AI counterfactual, not a one-year flow compared mechanically with a recurring annual policy cost. The framework separates an externally disciplined AI productivity scenario, bridged into $\phi^{Y,nom}_s$ and accumulated when needed, from its country-specific translation through readiness, exposure, adoption, informality, and digital constraints in $T_{i,s,t}$.

\subsection{Numerically closed AI-productivity scenario ladder}\label{subsec:phi_ladder_main}

Let $\phi^{raw}_{s,u}$ denote the raw AI shock reported by a source under unit $u\in\{TFP,LP,Y\}$, where $TFP$ is total factor productivity, $LP$ is labour productivity, and $Y$ is GDP-equivalent output. The model-ready real GDP-equivalent shock is obtained only after an explicit bridge:
\begin{equation}
    \phi^Y_s=\kappa^{u\to Y}_s\phi^{raw}_{s,u},\qquad 0<\kappa^{u\to Y}_s\leq 1\ \text{(baseline convention)}.
\end{equation}
For $u=Y$, the bridge is the identity, $\kappa^{Y\to Y}=1$. For $u=LP$, $\kappa^{LP\to Y}\leq1$ reflects that GDP gains cannot exceed labour-productivity gains unless employment rises. For $u=TFP$, standard growth accounting implies amplification through capital deepening ($\Delta\ln Y=\Delta\ln TFP/(1-\alpha_K)>\Delta\ln TFP$); the baseline restriction $\kappa^{TFP\to Y}\leq1$ is therefore a deliberate conservative convention that excludes induced capital accumulation from the fiscal-base shock. A capital-deepening robustness bridge $\kappa^{TFP\to Y}\in[1,1/(1-\alpha_K)]$ may be reported separately and must be labelled as such.
The nominal fiscal-base-compatible shock is then
\begin{equation}
    \phi^{Y,nom}_s=(1+\phi^Y_s)(1+\pi^{AI,Y}_s)-1,
\end{equation}
or, for small rates, $\phi^{Y,nom}_s\approx \phi^Y_s+\pi^{AI,Y}_s$. The exact expression is multiplicative, not log-additive, and $\pi^{AI,Y}_s$ may be negative. The nominal bridge is admissible only when fiscal-capture anchors and policy costs are nominal and share the endpoint horizon. For price-indexed benefits, poverty lines, or eligibility costs, the baseline instead pairs a real GDP-equivalent shock with real capture anchors or projects costs to that nominal horizon; price effects alone are never real fiscal capacity. Non-market or zero-price AI output lies outside GDP and the fiscal base, distinct from informality, biasing the frontier conservatively. Although source-unit supports are positive, draws with $\phi^{Y,nom}_s\leq0$ remain economic failure draws in all-draw probabilities and are excluded only from inversions conditional on $\phi^{Y,nom}_s>0$ (Section \ref{subsec:impossibility_score}). A persistently negative aggregate effect would reduce $fs^{eff}$ further, so the reported frontiers are upper bounds on that margin.

When the literature reports annual growth increments over a horizon, the recurring-cost comparison uses the induced level difference relative to a no-AI counterfactual. In the static one-period notation this object is written compactly as
\begin{equation}
    g^{AI,level}_{i,s,t}=\frac{Y^{AI}_{i,s,t}-Y^{0}_{i,t}}{Y^{0}_{i,t}}.
\end{equation}
If annual model-ready shocks are accumulated over horizon $H$, the level object is
\begin{equation}
    g^{AI,level}_{i,s,H}=\prod_{\tau=1}^{H}\left(1+\phi^{Y,nom}_{s,\tau}T_{i,s,\tau}\right)-1.
\end{equation}
Later equations write $g^{AI}_{i,s,t}$ for this fiscal-base-compatible level object unless discussing the annual increment. Official tables use accumulated $g^{AI,level}$ under the frozen 2024 anchor; any static one-period approximation is explicitly labelled. The recurring-cost test compares the endpoint annual fiscal-base increment with policy costs measured or projected for that endpoint, assuming the gain persists at least as long as the obligation. A shift from taxable rents toward lower prices and consumer surplus could reduce long-run capture even at constant $g^{AI}$, making the baseline optimistic about composition; innovation-driven compounding instead makes the level convention conservative. Present-value comparisons are robustness only and are never mixed with the baseline annual-flow threshold. Accumulated-horizon inversion uses
\begin{equation}\label{eq:accumulated_threshold_condition_v22}
    \widetilde{MFC}^{gross}_{i,s,r,t}\left[\prod_{\tau=1}^{H}\left(1+\phi^{Y,nom}_{s,\tau}T_{i,s,\tau}\right)-1\right]
    \geq R^{req,gross}_{i,p,s,r,t}(\xi).
\end{equation}
The debt-consistent accumulated-horizon version replaces $R^{req,gross}_{i,p,s,r,t}(\xi)$ with $R^{req,gross,DC}_{i,p,s,r,t}(\xi)$. Therefore the closed-form $T^*=R/(\phi^{Y,nom}\widetilde{MFC}^{gross})$ used later is a one-period static approximation. Under constant $\phi^{Y,nom}$ and constant $T$ over $H$, the corresponding accumulated-horizon threshold is
\begin{equation}\label{eq:Tstar_accumulated_constant_v22}
    T^{*,H}_{i,p,s,r,t}(\xi)=
    \frac{\left(1+R^{req,gross}_{i,p,s,r,t}(\xi)/\widetilde{MFC}^{gross}_{i,s,r,t}\right)^{1/H}-1}{\phi^{Y,nom}_{s}},
\end{equation}
with non-constant paths solved numerically from Equation \eqref{eq:accumulated_threshold_condition_v22}; the debt-consistent analogue again replaces $R^{req,gross}$ with $R^{req,gross,DC}$.
\begin{table}[H]
\centering
\caption{Mandatory AI-shock objects and admissible use}
\label{tab:ai_shock_objects_v17}
\small
\begin{tabular}{p{0.25\textwidth}p{0.48\textwidth}p{0.18\textwidth}}
\toprule
Object & Meaning & Admissible use \\
\midrule
$\phi^{raw}_{s,u}$ & Literature shock in source unit $u\in\{TFP,LP,Y\}$ & Input only; never used directly in fiscal-space equations \\
$\phi^{Y,nom}_{s,t}$ & Annual nominal GDP-equivalent, fiscal-base-compatible shock after unit and nominal bridging & Annual increment to be accumulated or used in a labelled static approximation \\
$g^{AI,level}_{i,s,H}$ & AI-induced level gain relative to the no-AI counterfactual over horizon $H$ & Object compared with recurring annual policy costs \\
\bottomrule
\end{tabular}
\end{table}
\noindent The implementation must never set $g^{AI}_{i,s,t}=\phi^{Y,nom}_{s,t}$ directly: the translation factor is always applied, and $g^{AI}_{i,s,t}=\phi^{Y,nom}_{s,t}T_{i,s,t}$ is admissible only as an explicitly labelled one-period static approximation. Baseline recurring-cost comparisons use the level object, not the raw annual increment. In particular, a high labour-productivity input is reported after applying $\kappa^{LP\to Y}$ and the nominal bridge, never as the source-unit value.

Only after bridging does the model assign the low, mid, high, or stress GDP-equivalent support listed in Table \ref{tab:phi_numerical_closure}.
Because historical fiscal-capture anchors are usually constructed from nominal revenue and nominal GDP, the baseline fiscal-space implementation must use $\phi^{Y,nom}_s$ and the induced $g^{AI,level}$, or reconstruct the fiscal-capture anchor in real terms with explicit revenue and GDP deflators. Results that multiply a raw real/TFP/labour-productivity AI shock by a nominal historical capture distribution are reported only as heuristic plausibility checks, not as baseline fiscal-space estimates. For compactness, later formulas write $\phi_s$ only as shorthand for the baseline-compatible object $\phi^{Y,nom}_s$.

Micro-level productivity experiments discipline plausibility in exposed tasks but do not directly identify aggregate TFP, GDP, or nominal fiscal-base effects. For example, \citet{noy2023experimental} estimate large effects for professional writing tasks, and \citet{brynjolfsson2023generative} estimate productivity gains in customer support. Those studies justify allowing positive task-level productivity gains, but they are not used as direct macro-growth estimates.

\begin{longtable}{p{0.13\textwidth}p{0.18\textwidth}p{0.13\textwidth}p{0.25\textwidth}p{0.20\textwidth}}
\caption{Numerical closure of the base AI shock $\phi_s$}\label{tab:phi_numerical_closure}\\
\toprule
Scenario & Annual support & Unit & Source discipline & Use in the framework \\
\midrule
\endfirsthead
\toprule
Scenario & Annual support & Unit & Source discipline & Use in the framework \\
\midrule
\endhead
Conservative, $\phi^{low}$ & $[0.0003,0.0007]$ & TFP-equivalent & Annualized from the task-based macro accounting benchmark in \citet{acemoglu2024simple}, which implies a TFP gain of no more than about $0.66\%$ over ten years & Lower-bound macro scenario; used when affected tasks, cost savings, or adoption are limited \\
Moderate, $\phi^{mid}$ & $[0.0025,0.0060]$ & TFP-equivalent & OECD micro-to-macro estimates of annual aggregate TFP gains from AI of about $0.25$--$0.60$ percentage points over a ten-year horizon \citep{filippucci2024miracle} & Main literature-disciplined scenario for the baseline analysis \\
High frontier input, $\phi^{high,LP}$ & $[0.0040,0.0130]$ before bridge & Labour-productivity equivalent & OECD G7 estimates report annual labour-productivity gains in the approximate range $0.4$--$1.3$ percentage points in high-exposure, high-adoption economies \citep{oecd2025g7} & Frontier-adoption case; converted by $\kappa^{LP\to Y}$ and not treated as the baseline for emerging economies \\
Disruptive stress, $\phi^{stress}$ & $[0.0100,0.0200]$ & GDP/TFP-equivalent stress & Deliberate stress support beyond the baseline macro literature & Diagnostic only; excluded from expected-feasibility claims \\
\bottomrule
\end{longtable}

The preferred stochastic assignment applies bounded PERT to the source-unit shock:
\begin{equation}
    \phi^{raw}_{s,u}\sim \mathrm{PERT}(a_s,m_s,b_s).
\end{equation}
Conversion then follows the bridges above; $a_s$ and $b_s$ are the table's source-unit bounds, and $m_s$ is their midpoint unless a source-specific central estimate is available.

For transparent non-stochastic reporting, the deterministic scenario values are
\begin{equation}
    \bar\phi^{low}=0.0005,
    \qquad
    \bar\phi^{mid}=0.00425,
    \qquad
    \bar\phi^{high,LP}=0.0085,
    \qquad
    \bar\phi^{stress}=0.015.
\end{equation}
\subsection{Country translation factor}\label{subsec:translation_factor_main}

The static country-specific level shock is
\begin{equation}
    g^{AI,level}_{i,s,t}=\phi^{Y,nom}_s T_{i,s,t},
\end{equation}
where $T_{i,s,t}\in[0,1]$ is the share of the baseline-compatible nominal GDP-equivalent scenario translated into a domestic fiscal-base level gain. Unless a dynamic horizon is explicit, $g^{AI}_{i,s,t}$ below denotes this level object. Translation uses three distinct inputs:
\begin{itemize}
    \item $E^{prod}_{i,t}$: productive AI exposure, or the share of tasks, occupations, or sectors where AI can plausibly raise output or efficiency.
    \item $A_{i,t}$: AI preparedness or absorptive capacity, including digital infrastructure, human capital and labour-market policies, innovation and economic integration, and regulation and ethics.
    \item $q^{prod}_{i,s,t}$: realized productive adoption, meaning AI use that has been integrated into production processes rather than merely accessed or tested.
\end{itemize}
The IMF AI Preparedness Index (AIPI) supplies $A_{i,t}$ consistently across countries \citep{imf2024aipi,cazzaniga2024genai}. World Bank--ILO ranges show that productive exposure in Latin America and the Caribbean is well below total exposure \citep{worldbankilo2024lacgenai}; the baseline therefore assigns the productivity-enhancing component to $E^{prod}$ and automation exposure to the labour-share/displacement block (Section \ref{subsec:prod_disp_exposure_main}).

Where occupational and sectoral maps exist, $E^{prod}$ is their convex combination (weight $\varpi$), with occupational exposure weighted by employment and sectoral exposure by value added (Appendix \ref{app:ai_shock_module}); the pilot uses the regional fallback below.
When the application is restricted to Latin America and lacks country-specific task maps, the fallback support is
\begin{equation}
    E^{prod}_{i,t}\in[0.08,0.14],
\end{equation}
interpreted as a regional support for productivity-enhancing GenAI exposure, not as total exposure. Total exposure $[0.26,0.38]$ is reported only as a broad-exposure robustness check, and full-automation exposure $[0.02,0.05]$ is used only in the labour-displacement channel.

The implementation adopts published AIPI values consistently across countries; Appendix \ref{app:ai_shock_module} documents reconstruction from its four dimensions.


\subsection{Calibrable logistic adoption module}\label{subsec:logistic_adoption_main}

The framework distinguishes technical access, task-level use, and productive adoption integrated into output processes; micro productivity gains need not become macro gains.

The baseline use equation is
\begin{equation}\label{eq:q_use_logistic_main_v10}
    q^{use}_{i,s,t}=\bar q_{i,t}\Lambda(\ell^q_{i,s,t}),
    \qquad
    \Lambda(x)=\frac{1}{1+e^{-x}}.
\end{equation}
To avoid double counting productive exposure, the baseline adoption index does not include $E^{prod}$ in the logistic adoption equation. Productive exposure enters the translation score directly, while adoption is restricted to preparedness and frictions:
\begin{equation}\label{eq:eta_noE_main_v10}
    \ell^q_{i,s,t}=\mu_{i,s,t}+\nu_A A_{i,t}-\nu_I I_{i,t}-\nu_G Gap_{i,t}.
\end{equation}
Here $A$ is AI preparedness, $I$ is the relevant adoption-side informality or organizational-friction index, and $Gap$ captures digital-access, skills, or infrastructure gaps not already embedded in $A$. In the baseline, $Gap$ is constructed only from indicators excluded from the AIPI preparedness index. If overlapping indicators are unavoidable, the implementation uses a residualized gap, $Gap^{\perp}_{i,t}=Gap_{i,t}-\operatorname{Proj}_{A}(Gap_{i,t})$, or moves the overlapping gap measure to robustness. A robustness specification may include $E^{prod}$ in $\ell^q$, but then exposure must not also enter the translation score as an independent multiplicative factor.

The country adoption ceiling is
\begin{equation}\label{eq:qbar_main_v10}
    \bar q_{i,t}=\max\left\{0,\min\left[1,1-\omega_I I_{i,t}-\omega_G Gap_{i,t}\right]\right\}.
\end{equation}
The economic mechanism behind the informality term is that informality proxies small firm scale, credit constraints, low managerial capital, and weak complementary investments, which limit productive integration of AI even when consumer-side use is high \citep{worldbank2021informality}. Because informal firms do adopt some digital technologies partially, $\omega_I<1$ is the economically plausible range; a unit penalty is an extreme convention. For high-informality countries such as the pilot's primary case, the sensitivity of results to $\omega_I$ is reported prominently: in the official runs, the $p10\to p90$ tornado range of $\omega_I$ on $V^{gross}$ is exactly zero in every headline cell, and its Spearman rank correlation with $V^{gross}$ is within sampling noise ($|\rho^S|\leq0.015$, against tenth-ranked driver magnitudes of $0.03$--$0.05$), because the anchored-target calibration places productive adoption well below the country ceiling $\bar q_i$, leaving the ceiling weight inactive at the official calibration. The full ranking is a post-baseline reporting extension computed exclusively from the stored official draws and recorded in the amendment registry of the reproducibility package.
Informality and $Gap$ may enter both short-run diffusion and long-run adoption capacity only if their empirical content is partitioned across the diffusion and ceiling equations or one margin is residualized against the other. This is an anchored calibration, not an identified cross-country forecasting equation; unless a time trend or scenario path is explicitly calibrated, the anchored intercept is held fixed after $t_0$, $\mu_{i,s,t}=\mu_{i,s,t_0}$ for $t>t_0$. The logit inversion of the intercept, interior projection of the target, initial condition, and support restrictions are specified in Appendix \ref{app:adoption_module}.
Productive adoption follows a lagged diffusion law with an explicit ceiling projection:
\begin{equation}\label{eq:q_prod_diffusion_main_v10}
    q^{prod}_{i,s,t}=\min\left\{\bar q_{i,t},
    (1-\lambda_q)q^{prod}_{i,s,t-1}+\lambda_q q^{use}_{i,s,t}\right\},
    \qquad \lambda_q\in\{0.15,0.30,0.50\}\quad\text{in baseline, with }\lambda_q=1\text{ only as an instant-adoption stress test}.
\end{equation}
The projection is needed because $q^{prod}$ is stock-like; if the ceiling $\bar q_{i,t}$ falls after a shock to infrastructure, informality, or access, the unprojected stock could otherwise exceed the current feasible ceiling.
\subsection[No-double-counting construction of T]{No-double-counting construction of $T_{i,s,t}$}\label{subsec:ndc_translation_main}

The preferred country translation factor separates productive exposure from productive adoption without reusing the same information twice. Under the baseline convention, $E^{prod}$ enters the translation score, while $q^{prod}$ is calibrated from preparedness and frictions, not from $E^{prod}$. Define
\begin{equation}\label{eq:S_ndc_main_v10}
    S^{NDC}_{i,s,t}=\mathbf{1}\{q^{prod}_{i,s,t}>0,\ E^{prod}_{i,t}>0\}(\tilde E^{prod}_{i,t})^{\beta}(\tilde q^{prod,num}_{i,s,t})^{\gamma},
\end{equation}
where $NDC$ denotes no-double-counting and $\tilde E^{prod}_{i,t}$ and $\tilde q^{prod,num}_{i,s,t}$ are the floor-regularized inputs $\tilde X=\epsilon+(1-\epsilon)X$ defined in Appendix \ref{app:ai_shock_module}. The full score with direct preparedness is a labelled robustness specification, valid only when $q^{prod}$ is independently observed from adoption data rather than generated from $A$ or $E^{prod}$ (Appendix \ref{app:ai_shock_module}).

The baseline translation factor is
\begin{equation}\label{eq:T_ndc_main_v10}
    T_{i,s,t}=\min\left\{1,\frac{S^{NDC}_{i,s,t}}{S^{frontier}_{s,t}}\right\}.
\end{equation}
Because $S^{NDC}$ contains the substantive-adoption indicator, $q^{prod}_{i,s,t}\leq0$ implies $S^{NDC}_{i,s,t}=0$, $T_{i,s,t}=0$, and therefore $g^{AI}_{i,s,t}=0$ in the static implementation.
$S^{frontier}_{s,t}$ is an external scenario-year benchmark from high-preparedness, high-exposure, high-adoption economies, never an internal pilot percentile that could assign $T\approx1$ within a weak sample. It remains fixed when inversions solve for $q^{prod,*}$, $E^{prod,*}$, or $T^*$, avoiding circular normalization. Any within-sample percentile is labelled only as a relative-index robustness check. Because $T=1$ denotes benchmark-like translation, population consistency requires the benchmark to match the population underlying $\phi_s$: an average-advanced-economy estimate uses that average set, not its top performer; a stricter benchmark double-discounts adoption and must be labelled as an additional conservative wedge.

\subsection{Productive exposure versus displacement exposure}\label{subsec:prod_disp_exposure_main}

Productivity-enhancing, labor-displacement, and total exposure are distinct, not interchangeable:
\begin{equation}
    E^{prod}_{i,t}\leq E^{total}_{i,t},\qquad E^{disp}_{i,t}\leq E^{total}_{i,t}.
\end{equation}
For source-defined non-overlapping categories, $E^{prod}_{i,t}+E^{disp}_{i,t}\leq E^{total}_{i,t}$ is imposed; otherwise overlap is documented, not forced to zero. World Bank--ILO evidence anchors Latin American total exposure at approximately $26$--$38$ percent of employment, productivity-transforming exposure at $8$--$14$ percent, and automation risk at $2$--$5$ percent \citep{worldbankilo2024lacgenai}. The baseline initializes
\begin{equation}
    E^{prod}_{i,t}\in[0.08,0.14],\qquad
    E^{disp}_{i,t}\in[0.02,0.05],\qquad
    E^{total}_{i,t}\in[0.26,0.38],
\end{equation}
unless better country-sector data exist. In the labor-share module, $E^{auto}$ maps to $E^{disp}$ with aggregate automation-risk data and $E^{aug}$ to $E^{prod}$ with productivity-enhancing data; finer task maps estimate them separately.

\section{Labor-Share Channel}\label{sec:labor_share}

Let
\begin{equation}
    LS_{i,t}=\frac{W_{i,t}}{Y^0_{i,t}}
\end{equation}
denote the labor share of income, with baseline total labor compensation $W_{i,t}$ and no-AI GDP or value added $Y^0_{i,t}$. Legacy calibration notation $Y_{i,t}$ is identified with $Y^0_{i,t}$ before computing ratios. At the anchor $t_0$, $Y^0_{i,t_0}=Y_{i,t_0}$, treating the GenAI-induced component as negligible at $t_0$; $LS_{i,t}\in(0,1)$ is required for the logit transformation.
The complement is not labelled mechanically as corporate profits. It is a non-labor residual share,
\begin{equation}
    NLS_{i,t}=1-LS_{i,t},
\end{equation}
which may include mixed income, depreciation, net taxes on production, gross operating surplus, rents, and other non-labor components depending on the national-accounts concept used. In economies with large self-employment or informal employment, the implementation distinguishes raw and mixed-income-adjusted labor shares; Appendix \ref{app:labor_share_module} defines both measures, and the baseline notation $LS_{i,t}$ denotes the preferred adjusted measure \citep{gollin2002getting}.

The labor-share block does not forecast general-equilibrium wage effects. Its narrower purpose is to allocate the AI-induced output increment between labor and capital/profit income, whose taxability and fiscal capture differ. Automation can shift this distribution even when aggregate output rises \citep{acemoglu2018race,acemoglu2019automation,karabarbounis2014global}.

The model separates exposure to tasks AI can automate or substitute, $E^{auto}_{i,t}$, from those it can augment or complement, $E^{aug}_{i,t}$. Reskilling, task reinstatement, or human--AI complementarity capacity $RST_{i,t}$ may be proxied by active labour-market policy spending, ILOSTAT training participation, or an AIPI human-capital component residualized against $A$ to avoid double counting. The latent labor-share pressure is
\begin{equation}
    \Delta \ell_{i,s,t}
    =
    -\left(\delta_s-\zeta_s RST_{i,t}\right)
    E^{auto}_{i,t}q^{prod}_{i,s,t}
    +
    \psi_s E^{aug}_{i,t}q^{prod}_{i,s,t},
\end{equation}
where $\delta_s\geq0$ is automation pressure, $\psi_s\geq0$ augmentation pressure, and $\zeta_s\geq0$ mitigation through reskilling or task reinstatement. The convention $\psi_s\geq0$ is scenario-specific, not theoretical: the labor-share effect of augmentation depends on substitution elasticities between augmented labor and capital. Low-$\psi_s$ cells approach neutrality, reached exactly only in the no-labor-augmentation ablation ($\psi_s=0$). Baseline $RST$ mitigates displacement rather than independently raising labor share. A conservative restriction is
\begin{equation}
    0\leq \zeta_s RST_{i,t}\leq \delta_s,
\end{equation}
with possible overcompensation allowed only in an explicitly augmentation-dominant robustness scenario.
A single productive-adoption stock $q^{prod}$ scales both pressures, although automation responds to labor costs and augmentation to complementary skills. Using one diffusion object is a deliberate simplification with ambiguous direction, partially spanned by the $\delta_s/\psi_s$ scenario grid. If $E^{auto}$ captures only full-automation risk, displacement pressure is a lower bound because partial task automation also displaces labor income; the automation-heavy scenario partially compensates through higher $\delta_s$.

To keep the post-shock labor share inside the unit interval, the target labor share is defined on the logit scale:
\begin{equation}
    LS^{target}_{i,s,t}
    =
    \Lambda\left(
    \log\frac{LS_{i,t}}{1-LS_{i,t}}+
    \Delta \ell_{i,s,t}
    \right),
\end{equation}
where $\Lambda(x)=1/(1+e^{-x})$. Observed factor shares need not jump immediately to this target. A partial-adjustment rule gives
\begin{equation}
    LS'_{i,s,t}
    =
    (1-\lambda_{LS})LS_{i,t}
    +
    \lambda_{LS}LS^{target}_{i,s,t},
    \qquad 0\leq\lambda_{LS}\leq1.
\end{equation}
The AI-induced labor-share change is therefore
\begin{equation}
    \Delta LS^{AI}_{i,s,t}=LS'_{i,s,t}-LS_{i,t}.
\end{equation}

In accumulated-horizon implementations, the counterfactual labor share $LS_{i,t}$ is held at its anchor-year value unless a documented counterfactual path is available, and the latent pressure $\Delta\ell_{i,s,t}$ is evaluated once at the endpoint year using the endpoint productive-adoption stock; because $q^{prod}$ is a stock object, applying $\Delta\ell$ cumulatively each year would double-count adoption.
The resulting labor and non-labor income flows follow exact accounting identities that exhaust the AI-induced increment; Appendix \ref{app:labor_share_module} gives the full derivation. When $\Delta NLS^{AI}<0$ in augmentation-dominant scenarios, fiscal capture of the decrement is treated symmetrically for accounting simplicity. This abstracts from real tax asymmetries such as loss carry-forwards, an empirical risk listed in Table \ref{tab:preferred_fiscal_channels}. A decline in labor share is not mechanically a decline in total labor income; the framework reports $\Delta LS^{AI}$, $\Delta W^{AI}/Y^0$, $\Delta NLS^{AI}/Y^0$, and $\Delta\Pi^{AI,gross}/Y^0$ to distinguish distribution, labor-income flows, residual income, and the taxable capital/profit base.




\section{Fiscal Conversion Block}\label{sec:fiscal_conversion}

Additional output does not automatically become public revenue, so fiscal conversion is a tax-base-channel problem rather than an average tax-to-GDP ratio. The object is effective marginal capture after base composition, observed effective rates, compliance, exemptions, leakage, administrative costs, and transition costs are assigned without double counting. Under the statutory-accrual convention, revenue belongs to the channel where income legally accrues. Economic incidence shifting---including corporate-tax burdens on labor under capital mobility and incomplete VAT pass-through---is behavioral, lies outside the accounting baseline, and is covered only by the erosion extension in Appendix \ref{app:fiscal_conversion_module}. This implements the policy logic of \citet{brollo2024fiscal} on labor disruption, social protection, capital and corporate taxation, market power, and technological uncertainty.

Let the AI-induced output increment be
\begin{equation}
    \Delta Y^{AI}_{i,s,t}=g^{AI}_{i,s,t}Y^0_{i,t}.
\end{equation}
Output-side income is disciplined by the functional-income identity
\begin{equation}
    \Delta Y^{AI}_{i,s,t}=\Delta W^{AI}_{i,s,t}+\Delta NLS^{AI}_{i,s,t}+\Delta O^{AI}_{i,s,t}.
\end{equation}
The baseline labor-share block sets $\Delta O^{AI}_{i,s,t}=0$ because labor income and the non-labor residual exhaust the AI-induced output increment. A residual $\Delta O^{AI}$ may be activated only as an explicitly labelled robustness channel for output not mapped to the two main income flows. The taxable capital/profit base is a component of the non-labor residual, not the entire residual. The accounting chain is
\begin{equation}\label{eq:capital_base_chain_v17}
\begin{split}
    \Delta NLS^{AI}_{i,s,t}
    &\rightarrow \Delta\Pi^{AI,gross}_{i,s,t}
    =\chi^{Kbase}_{i,t}\Delta NLS^{AI}_{i,s,t},\\
    \Delta\Pi^{AI,gross}_{i,s,t}
    &\rightarrow \Delta\Pi^{AI,dom}_{i,s,t}
    =\chi^{dom}_{i,s,t}(1-\psi^{shift}_{i,s,t})\Delta\Pi^{AI,gross}_{i,s,t},\\
    \Delta\Pi^{AI,dom}_{i,s,t}
    &\rightarrow \Delta\Pi^{AI,disp}_{i,s,t}
    =(1-\tau^{K,disp}_{i,s,r,t})\Delta\Pi^{AI,dom}_{i,s,t}.
\end{split}
\end{equation}
The capital fiscal-capture rate $\tau^{K,fisc,eff}_{i,s,r,t}$ and disposable-income rate $\tau^{K,disp}_{i,s,r,t}$ may differ because compliance, corporate taxation, withholding, and related institutions need not map one-to-one into household income. Consumption is an induced base from disposable income and import leakages, not a third functional output component. Thus the fiscal weights are tax-base exposure ratios, not exhaustive shares of $\Delta Y^{AI}$, and need not sum to one.

\subsection{Gross and effective fiscal conversion}\label{subsec:gross_effective_fiscal_conversion}

For each fiscal channel $j\in\{L,K,C,AI\}$, define the taxable-base exposure weight per unit of AI-induced output as
\begin{equation}
    \omega^j_{i,s,t}
    =\frac{\Delta B^{j,tax}_{i,s,t}/Y^0_{i,t}}{g^{AI}_{i,s,t}},
\end{equation}
where $\Delta B^{j,tax}$ is the additional taxable base associated with channel $j$. For the AI-rent channel this exposure weight is written $\lambda^{AI}_{i,s,t}$ throughout the paper; the sectorally indexed $\omega^{AI}_{i,k,s,t}$ of Appendix \ref{app:informality_module} denotes sectoral shares of the AI shock and is an unrelated object. In particular, the capital/profit exposure weight is constructed on the domestic, profit-shifting-adjusted taxable-profit base:
\begin{equation}\label{eq:omega_k_dom_base_v25}
\begin{split}
    \omega^K_{i,s,t}
    &=\frac{\Delta\Pi^{AI,dom}_{i,s,t}/Y^0_{i,t}}{g^{AI}_{i,s,t}} \\
    &=\chi^{dom}_{i,s,t}(1-\psi^{shift}_{i,s,t})\chi^{Kbase}_{i,t}
      \frac{\Delta NLS^{AI}_{i,s,t}/Y^0_{i,t}}{g^{AI}_{i,s,t}},
\end{split}
\end{equation}
whenever $g^{AI}_{i,s,t}>0$. If $g^{AI}_{i,s,t}=0$, the draw is treated as an economic non-crossing draw and ratio-based exposure weights are not computed. Equation \eqref{eq:omega_k_dom_base_v25} is the baseline convention: foreign ownership and ordinary profit shifting live in the base-side construction of $\Delta\Pi^{AI,dom}$ and are not multiplied again in $\tau^{K,fisc,eff}$ or $\chi^K$. A rate-side profit-shifting robustness convention is allowed only if $\psi^{shift}$ is set to zero in Equation \eqref{eq:capital_base_chain_v17}.

The preferred gross marginal fiscal conversion is
\begin{equation}\label{eq:omega_k_net_v16}
    \omega^{K,net}_{i,s,t}=\begin{cases}
    \max\{0,\omega^K_{i,s,t}-\lambda^{AI,Kbase}_{i,s,t}\}, & \text{if AI rents are treated as a subset of the capital/profit base},\\
    \omega^K_{i,s,t}, & \text{if AI rents are additional to the measured capital/profit base}.
    \end{cases}
\end{equation}
Under the subset convention, the implementation verifies $\lambda^{AI,Kbase}_{i,s,t}\leq\omega^K_{i,s,t}$. If this condition is violated, the draw is labelled convention-inconsistent and reported; it is not silently truncated. The $\max\{0,\cdot\}$ operator is retained only as a numerical safeguard.
\begin{equation}\label{eq:mfc_gross_channel_v10}
    MFC^{gross}_{i,s,r,t}
    =
    \omega^L_{i,s,t}\tau^{L,fisc,eff}_{i,s,r,t}
    +\omega^{K,net}_{i,s,t}\tau^{K,fisc,eff}_{i,s,r,t}
    +\omega^C_{i,s,t}\tau^{C,eff}_{i,s,r,t}
    +\lambda^{AI}_{i,s,t}\tau^{AI,eff}_{i,s,r,t}.
\end{equation}
Effective fiscal rates $\tau^{j,fisc,eff}$ already include rate-side compliance, informality, exemption, and capturability adjustments; base-side adjustments such as consumption exemptions inside $\Delta C^{AI,taxable}$ are not embedded again. Equation \eqref{eq:mfc_gross_channel_v10} therefore applies no additional $\chi^j$; with statutory or normative rates, the appendix first constructs $\tau^{j,fisc,eff}=\tau^{j,base}\chi^j$. The $\omega^L$, $\omega^K$, and $\omega^C$ objects are tax-base response ratios, not probability weights. A negative robustness-channel response is reported as a negative base contribution, never hidden by renormalization.
Labor informality is assigned to exactly one side of the labor channel. The rate-side convention is the baseline, consistent with the convention that $\tau^{eff}$ already includes informality: $\omega^L$ is constructed from gross $\Delta W^{AI}$ and $\chi^L=(1-I^{L,marg})^{\rho_L}$ enters once inside $\tau^{L,fisc,eff}$. Under the alternative base-side convention, $\omega^L$ is constructed from $\Delta W^{taxable}$ in Equation \eqref{eq:labor_informality_base_side_v25}, and $\tau^{L,fisc,eff}$ must not include $\chi^L$. The base-side and rate-side penalties are never applied together.

A key implementation convention is required here. Equation \eqref{eq:mfc_gross_channel_v10} is a realized-scenario marginal-conversion object. In the historical reduced-form mode, $MFC^{gross}$ is drawn directly and is fixed when solving threshold inversions. In the channel-based mode, however, some exposure ratios $\omega^j$ may themselves be functions of the AI shock because $\omega^j=(\Delta B^{j,tax}/Y^0)/g^{AI}$. Therefore the compact product $MFC^{gross}g^{AI}$ is exact at the evaluated scenario, but closed-form inversions that vary $g^{AI}$ are valid only under a frozen-weight convention. When channel weights are recomputed as the shock varies, the implementation solves thresholds from the channel-revenue function in Appendix \ref{app:fiscal_conversion_module}, Equation \eqref{eq:channel_revenue_function_v24}.

Each simulation must declare one AI-rent convention. If AI rents are a subset of the capital/profit base, $\omega^{K,net}=\max\{0,\omega^K-\lambda^{AI,Kbase}\}$ with the consistency verification of Equation \eqref{eq:omega_k_net_v16}; if they are additional, $\omega^{K,net}=\omega^K$. The unit-compatible rescaling rule is in Appendix \ref{app:fiscal_conversion_module}.

Fixed costs and marginal-equivalent deductions are separate margins, and each cost component enters exactly one margin; administrative netting, transition annualization, and the one-margin rule are detailed in Appendix \ref{app:fiscal_conversion_module}. Define the net marginal gross-capture object
\begin{equation}\label{eq:mfc_tilde_v17}
    \widetilde{MFC}^{gross}_{i,s,r,t}
    =MFC^{gross}_{i,s,r,t}-ac^{ME}_{i,s,r,t}-tr^{ME}_{i,s,r,t}-leak^{ME}_{i,s,r,t}.
\end{equation}
The preferred implementation works with effective fiscal space directly:
\begin{equation}\label{eq:fs_gross_eff_v21}
    fs^{gross}_{i,s,r,t}=MFC^{gross}_{i,s,r,t}g^{AI}_{i,s,t},
    \qquad
    fs^{eff}_{i,p,s,r,t}=\widetilde{MFC}^{gross}_{i,s,r,t}g^{AI}_{i,s,t}
    -ac^{net,fix}_{i,p,s,r,t}-tr^{ann}_{i,s,r,t}-leak^{fix}_{i,s,r,t}.
\end{equation}
The leakage term is residual leakage not already embedded in the exposure weights $\omega^j$, effective rates $\tau^{j,fisc,eff}$, domestic profit base $\Delta\Pi^{AI,dom}$, or taxable consumption base $\Delta C^{AI,taxable}$. Import content, profit shifting, or foreign ownership must be placed either inside the relevant base/rate construction or inside residual leakage, but not both.

An optional residual-friction penalty $\Omega^{res}_{i,s,r,t}\geq0$ is used only in stress tests for frictions not already included in $\omega^j$, $\tau^{j,fisc,eff}$, $\Delta\Pi^{AI,dom}$, $\Delta C^{AI,taxable}$, fixed costs, marginal-equivalent deductions, or leakage. It is inactive in the baseline. When activated, the stress-test fiscal-space object is
\begin{equation}\label{eq:omega_res_stress_v17}
    fs^{eff,res}_{i,p,s,r,t}=fs^{eff}_{i,p,s,r,t}-\Omega^{res}_{i,s,r,t}.
\end{equation}
This term must never be used as a second leakage penalty for the same mechanism. The equivalent threshold rule is given in Appendix \ref{app:threshold_module}.

When a per-unit effective conversion rate is useful for reporting, it is derived as
\begin{equation}\label{eq:mfc_eff_v17}
    MFC^{eff}_{i,p,s,r,t}=\frac{fs^{eff}_{i,p,s,r,t}}{g^{AI}_{i,s,t}}
\end{equation}
provided $g^{AI}_{i,s,t}>0$. The implementation must not draw $MFC^{eff}$ independently or subtract the same costs inside both $MFC^{eff}$ and $fs^{eff}$. In simulations, $fs^{eff}$ is the primary object because it avoids unstable division by very small $g^{AI}$ values.

\begin{table}[H]
\centering
\caption{Preferred fiscal-conversion channels}
\label{tab:preferred_fiscal_channels}
\small
\begin{tabular}{p{0.15\textwidth}p{0.33\textwidth}p{0.36\textwidth}}
\toprule
Channel & Tax base created by AI shock & Main empirical risk \\
\midrule
$L$ labor & Formal wages, payroll-taxable employment income, personal income tax base & Informality, exemptions, low personal income tax base, displacement of labor income \\
$K$ capital/profits & Taxable subset of the non-labor residual: corporate profits, capital income, business income, and rents from complementary capital & Mixed-income measurement, depreciation, profit shifting, foreign ownership, deductions, loss carry-forwards, tax incentives \\
$C$ consumption & VAT- or sales-taxable consumption induced by additional disposable income & VAT exemptions, imported digital services, lower taxable domestic consumption multiplier \\
$AI$ AI rents & Excess profits, digital-service rents, AI-platform rents, data/cloud rents & No existing instrument, cross-border location of rents, political/legal limits \\
Costs/leakages & Enforcement cost, transitional compensation, imported cloud/software/hardware, shifted profits & Double counting if already embedded in effective rates \\
\bottomrule
\end{tabular}
\end{table}

\subsection{Consumption as an induced tax base}

Consumption is an induced tax base rather than a third functional component of output. The disposable-income rates $\tau^{j,disp}$ are not fiscal-capture rates, and import content and consumption exemptions enter exactly once, on either the base side or the rate side. The consumption channel is first-round only: induced consumption is not recycled into further income and consumption rounds. This truncation is consistent with the absence of general-equilibrium multipliers elsewhere in the framework and biases the frontier conservatively. Appendix \ref{app:fiscal_conversion_module} gives the full construction.

\subsection{From channel formula to regimes}

Fiscal regimes are indexed by $r$:
\begin{itemize}
    \item $r=0$: status quo - current effective tax structure and no new AI-rent capture.
    \item $r=1$: improved compliance - stronger administration and reporting, with unchanged broad legal structure.
    \item $r=2$: base broadening - lower exemptions and wider effective labor, capital, and consumption bases.
    \item $r=3$: AI-rent capture - explicit taxation of digital or AI-specific rents.
    \item $r=4$: combined reform - compliance gains, base broadening, and AI-rent capture.
\end{itemize}
The regime index changes $\tau^{j,fisc,eff}$, $\tau^{AI,eff}$, $ac^{net,fix}$, $tr^{ann}$, $leak^{fix}$, $ac^{ME}$, $tr^{ME}$, and $leak^{ME}$, but the scenario rent share $\lambda^{AI}_{i,s,t}$ has no $r$ index and regimes do not alter the AI shock. Base broadening acts on legal base-side margins: because exemptions and ordinary shifting enter $\Delta C^{AI,taxable}$ through $(1-exempt^C)$ and $\Delta\Pi^{AI,dom}$ through $\psi^{shift}$, regimes $r\geq2$ index $exempt^C_{i,r,t}$ and $\psi^{shift}_{i,s,r,t}$, inherited by $\omega^C$ and $\omega^K$; alternatively, a labelled rate-side convention acts through $\chi^j$. The same margin must never operate on both sides. Gains for $r\geq1$ are mechanical upper bounds because the baseline omits behavioral responses of evasion, avoidance, and relocation, especially for mobile AI rents. Every headline result relying on $r\geq1$ therefore requires the behavioral-erosion companion for $B^{j,taxable}$ in Appendix \ref{app:fiscal_conversion_module}; omitting it overstates reform payoffs relative to the conservative discipline used elsewhere. Fiscal-reform optimism must remain separate from technology optimism.

\subsection{Observed fiscal anchors}

The preferred fiscal data source for cross-country revenue anchors is the UNU-WIDER Government Revenue Dataset, which provides central and general government revenue series and, where possible, figures with and without natural-resource revenues \citep{unu2025grd}. OECD Revenue Statistics in Latin America and the Caribbean is used as a harmonized tax-to-GDP cross-check for Peru, Chile, Colombia, and Mexico \citep{oecd2025revenue}. The baseline empirical implementation records three fiscal anchors:
\begin{equation}
    \theta^{tax}_{i,t}=\frac{TaxRev_{i,t}}{Y_{i,t}},\qquad
    \theta^{rev}_{i,t}=\frac{TotalRev_{i,t}}{Y_{i,t}},\qquad
    \theta^{nonres}_{i,t}=\frac{Rev^{nonres}_{i,t}}{Y_{i,t}}.
\end{equation}
The non-resource anchor is preferred for countries where commodity cycles can make fiscal capacity look temporarily high. These anchors discipline $MFC^{gross}$; they are not automatically equal to $MFC^{eff}$.

\subsection{Country-specific calibration rule}

For each country $i$, fiscal regime $r$, and scenario $s$, fiscal-conversion draws are constructed in two steps. First draw or compute gross marginal capture. For a positive-capture conditional feasibility benchmark,
\begin{equation}
    MFC^{gross}_{i,s,r,t}
    \sim \text{PERT}\left(P10(MFC^{hist,gross,+}_i),P50(MFC^{hist,gross,+}_i),P90(MFC^{hist,gross,+}_i)\right),
\end{equation}
with country-regime bounds. A tighter $\text{PERT}(P25,P50,P75)$ is a low-dispersion robustness variant, not the baseline, because quartile bounds truncate support more than decile bounds. Unconditional assessment uses the two-component mixture in Section \ref{sec:historical_plausibility}. Unit matching follows the canonical nominal/real rule in Section \ref{subsec:phi_ladder_main}; in particular, the baseline never combines real/TFP-equivalent $g^{AI}$ with nominal historical $MFC^{gross}$.

Then compute $fs^{eff}=\widetilde{MFC}^{gross}g^{AI,level}-ac^{net,fix}-tr^{ann}-leak^{fix}$ and derive $MFC^{eff}=fs^{eff}/g^{AI,level}$ only when needed for reporting and $g^{AI,level}>0$. This ensures that the historical plausibility check compares gross objects with gross objects, while effective feasibility uses net fiscal space.


\section{Historical Fiscal-Capture Plausibility Check}\label{sec:historical_plausibility}

The fiscal block should be checked against country-specific historical experience. The historical object is a gross revenue-capture measure, not an effective net-of-cost measure. For each country-year, define
\begin{equation}\label{eq:mfc_hist_gross_v10}
    MFC^{hist,gross}_{i,t}
    =\frac{(R_{i,t}-R_{i,t-1})/Y_{i,t-1}}{(Y_{i,t}-Y_{i,t-1})/Y_{i,t-1}},
\end{equation}
where $R$ is the chosen revenue aggregate and $Y$ is nominal GDP. The revenue aggregate must be fixed within each country-window: central-government tax revenue, general-government tax revenue, total revenue, or non-resource revenue. A log-change version is
\begin{equation}\label{eq:buoyancy_hist_v10}
    Buoyancy^{hist}_{i,t}=\frac{\Delta\ln R_{i,t}}{\Delta\ln Y_{i,t}}.
\end{equation}
Because $\Delta R$ may be negative even when nominal GDP growth is positive, the implementation reports both the raw historical distribution and a non-negative-capture distribution:
\begin{equation}\label{eq:mfc_hist_positive_v10}
    H_i^+=\{MFC^{hist,gross}_{i,t}:MFC^{hist,gross}_{i,t}\geq0\}.
\end{equation}
Negative-capture years remain reported as fiscal-stress years. The positive-capture calibration $MFC^{gross}\sim MFC^{hist,gross}\mid MFC^{hist,gross}\geq0$ is explicitly conditional and labelled as a positive-capture scenario, not unconditional baseline capacity.
The empirical table must also report
\begin{equation}
    p_i^- = \Pr\left(MFC^{hist,gross}_{i,t}<0\right)
\end{equation}
inside the historical window. When the goal is unconditional uncertainty propagation, the preferred reduced-form draw is a two-component mixture,
\begin{equation}
    MFC^{gross}\sim
    \begin{cases}
    MFC^{hist,gross,-}, & \text{with probability }p_i^-,\\
    MFC^{hist,gross,+}, & \text{with probability }1-p_i^- ,
    \end{cases}
\end{equation}
Here $MFC^{hist,gross,+}$ and $MFC^{hist,gross,-}$ are the positive and fiscal-stress supports. PERT draws from $H_i^+$ are positive-capture conditional and may overstate unconditional capacity in crisis-prone economies. Nor does this imply $MFC^{eff}\geq0$: net fiscal space can be negative after costs and leakage, and negative $fs^{eff}$ draws are economic failures, not invalid draws.
The mixture is unconditional over the sign of capture but conditional on positive nominal GDP growth (cleaning rule 2), the relevant class for a positive AI-induced increment. Zero or negative increments remain economic non-crossings or failure draws without a capture ratio.

The positive-capture PERT is the baseline conditional-feasibility benchmark: it asks whether a policy could cross under years in which marginal revenue capture is positive. The two-component mixture is the unconditional fiscal-risk benchmark: it preserves the empirical probability of negative-capture years. The positive-capture conditional benchmark is reported in the Monte Carlo table of Section \ref{sec:empirical_results}; the unconditional mixture is reported, for both tax and total-revenue anchors, in the historical reduced-form output of the reproducibility package. Claims of expected feasibility must rely on the unconditional mixture or explicitly state that they are conditional on positive fiscal-capture years.

The historical plausibility comparison is
\begin{equation}\label{eq:historical_comparison_v10}
    MFC^{gross}_{i,s,r,t}\quad \text{versus}\quad
    \{P25,P50,P75,P90\}(MFC^{hist,gross,+}_i).
\end{equation}
A scenario requiring $MFC^{gross}>P90(MFC^{hist,gross,+}_i)$ is flagged as historically demanding. A scenario requiring $MFC^{eff}>P90(MFC^{hist,gross,+}_i)$ is not used as the primary comparison because it compares a net object to a gross historical object. Effective feasibility is assessed through $V^{gross}$ and $V^{net}$, while plausibility of fiscal capture is assessed using the gross-to-gross comparison.
The historical prior is composition-weighted: $MFC^{hist}$ reflects sectoral and factor composition, commodity cycles, and formalization, while a more formal and intangible-intensive AI increment pushes capture in opposite directions. Divergence between channel and historical modes therefore signals composition change, not specification error; the $P90$ flag applies under historical composition. The prior is also policy-inclusive: $MFC^{hist,gross}$ and $Buoyancy^{hist}$ embed discretionary rate changes, new taxes, and administration, rather than pure elasticities. It is thus generous to status quo $r=0$, since reform helped produce the historical upper tail, while reform regimes are compared with a history already containing reforms. The $P90$ flag retains this buoyancy-not-elasticity caveat.

The cleaning rules are:
\begin{enumerate}
    \item Use one government level consistently within each country.
    \item Exclude years with non-positive nominal GDP growth when computing increment ratios.
    \item Report raw and positive-capture distributions separately.
    \item Report results with and without pandemic years 2020--2021.
    \item Report results with and without resource revenue where GRD provides the distinction.
    \item Winsorize raw ratios at $P1/P99$ or report robust medians to avoid denominator artifacts.
\end{enumerate}
The comparison labels whether feasibility requires ordinary, demanding, or historically extreme capture rather than mechanically rejecting reform. Unit matching follows the canonical rule in Section \ref{subsec:phi_ladder_main}. The official implementation adopts its nominal convention: raw supports become nominal GDP-equivalent shocks accumulated into $g^{AI,level}$ under the frozen-anchor trajectory, with post-bridge supports reported in the primary-specification table of Section \ref{sec:empirical_results}. Products of real/TFP-equivalent shocks and nominal historical capture remain heuristic plausibility checks, never baseline fiscal-space estimates.


\section{Feasibility Thresholds and Fiscal Guardrails}\label{sec:thresholds}

\subsection{Structural accounting threshold}

The model's primary conservative feasibility object is
\begin{equation}\label{eq:Vgross_main_v10}
    V^{gross}_{i,p,s,r,t}=\frac{fs^{eff}_{i,p,s,r,t}}{c^{gross}_{i,p,t}}.
\end{equation}
The structural accounting threshold is $V^{gross}=1$. A net-cost version,
\begin{equation}\label{eq:Vnet_main_v10}
    V^{net}_{i,p,s,r,t}=\frac{fs^{eff}_{i,p,s,r,t}}{c^{net}_{i,p,t}},\qquad c^{net}_{i,p,t}>0,
\end{equation}
is robustness only when measured spending offsets leave a strictly positive net incremental cost. If $c^{net}=0$, existing spending already covers the policy and $V^{net}$ is not computed. Without a superscript, $V$ denotes $V^{gross}$.

Define the absolute and relative gaps as
\begin{equation}
    GAP^{gross}_{i,p,s,r,t}=fs^{eff}_{i,p,s,r,t}-c^{gross}_{i,p,t},
    \qquad
    RGAP^{gross}_{i,p,s,r,t}=V^{gross}_{i,p,s,r,t}-1.
\end{equation}

\subsection{Prudential buffers}

A prudential buffer $\xi\geq0$ requires
\begin{equation}\label{eq:prudential_main_v10}
    fs^{eff}_{i,p,s,r,t}\geq(1+\xi)c^{gross}_{i,p,t}.
\end{equation}
The standard grid $\xi\in\{0,0.05,0.10,0.25,0.50\}$ defines reporting margins under deep uncertainty, not structural constants.

\subsection{Debt-stabilizing primary-balance guardrail}

Let $PB^0_{i,t}$ and $PB^{stab}_{i,t}$ be baseline and debt-stabilizing primary balances as GDP shares. Alongside $Y^0$-denominated objects, the observed- versus counterfactual-GDP denominator difference is second-order, so $s^{PB}$ may be treated as a share of $Y^0$. The required correction is
\begin{equation}\label{eq:spb_main_v10}
    s^{PB}_{i,t}=\max\{0,PB^{stab}_{i,t}-PB^0_{i,t}\}.
\end{equation}
The debt-consistent prudential guardrail is
\begin{equation}\label{eq:debt_requirement_main_v10}
    fs^{eff}_{i,p,s,r,t}\geq (1+\xi)c^{gross}_{i,p,t}+s^{PB}_{i,t}.
\end{equation}
This is not the primary marginal self-financing benchmark, which asks whether $fs^{eff}\geq(1+\xi)c^{gross}$. The debt-consistent layer asks whether that space also covers the first-order primary-balance correction and must be reported as a strong/debt-consistent test, not as marginal autofinancing failure. Buffering both terms is an ultra-conservative labelled robustness check, never the baseline guardrail.
Because $PB^{stab}$ uses baseline growth, it ignores improved debt dynamics from AI-driven nominal growth and is conservative by construction. The timing caveat remains: feasibility is an endpoint condition, so a crossing year is not a launch year. Earlier launch creates transition deficits while accumulated gains remain below their endpoint, outside both the accounting threshold and first-order guardrail. Early launch requires transition financing outside the framework, and launch scenarios must state this explicitly. The required persistence horizon $H^{*}$ in Section~\ref{subsec:h_star_inversion} exposes this timing without becoming a forecast.

Define
\begin{equation}\label{eq:Rreq_main_v10}
    R^{req}_{i,p,t}(\xi)=(1+\xi)c^{gross}_{i,p,t},
    \qquad
    R^{req,DC}_{i,p,t}(\xi)=(1+\xi)c^{gross}_{i,p,t}+s^{PB}_{i,t}.
\end{equation}
The threshold inequalities are written most cleanly in fiscal-space form:
\begin{equation}\label{eq:effective_thresholds_main_v10}
    fs^{eff}_{i,p,s,r,t}\geq R^{req}_{i,p,t}(\xi),
    \qquad
    fs^{eff}_{i,p,s,r,t}\geq R^{req,DC}_{i,p,t}(\xi).
\end{equation}
Under the generalized gross-capture convention, the corresponding required gross revenue includes the fixed-cost subtractions only:
\begin{equation}\label{eq:gross_req_main_v17}
    R^{req,gross}_{i,p,s,r,t}(\xi)=R^{req}_{i,p,t}(\xi)+ac^{net,fix}_{i,p,s,r,t}+tr^{ann}_{i,s,r,t}+leak^{fix}_{i,s,r,t},
\end{equation}
\begin{equation}\label{eq:gross_req_dc_main_v17}
    R^{req,gross,DC}_{i,p,s,r,t}(\xi)=R^{req,DC}_{i,p,t}(\xi)+ac^{net,fix}_{i,p,s,r,t}+tr^{ann}_{i,s,r,t}+leak^{fix}_{i,s,r,t}.
\end{equation}

\subsection{Thresholds as required values}

The minimum AI-induced level-gain threshold exists only when net marginal gross capture is positive. Define
\begin{equation}\label{eq:valid_growth_threshold_v20}
    Valid^{g*,gross}_{i,p,s,r,t}(\xi)
    =\mathbf{1}\left\{\widetilde{MFC}^{gross}_{i,s,r,t}>0,\ R^{req,gross}_{i,p,s,r,t}(\xi)>0\right\}.
\end{equation}
For valid cells, Proposition \ref{prop:unified_inversion} gives the closed-form $g^{AI,*}$, $MFC^{gross,*}$, and $T^{*,static}$. With endogenous channels, the growth threshold is instead the smallest feasible $g\geq0$ satisfying
\begin{equation}\label{eq:gstar_channel_root_v24}
    fs^{eff,ch}_{i,p,s,r,t}(g)\geq R^{req}_{i,p,t}(\xi),
\end{equation}
using $R^{req,DC}$ for debt consistency and $fs^{eff,ch}(g)-\Omega^{res}$ when the residual stress penalty is active.
If $g^{AI}\leq0$ and the required revenue is positive, no finite gross-capture threshold exists. For reporting only, the derived effective per-unit threshold is
\begin{equation}\label{eq:mfcstar_eff_main_v17}
    MFC^{eff,*}_{i,p,s,r,t}(\xi)=\frac{R^{req}_{i,p,t}(\xi)}{g^{AI}_{i,s,t}},
\end{equation}
provided the user clearly states that fixed costs have already been netted out of $fs^{eff}$.
The maximum financeable gross policy cost is reported as two separate non-negative policy envelopes. The marginal accounting envelope is
\begin{equation}\label{eq:cmax_main_v17}
    c^{gross,max,+}_{i,p,s,r,t}(\xi)=\max\left\{0,\frac{fs^{eff}_{i,p,s,r,t}}{1+\xi}\right\}.
\end{equation}
The stronger debt-consistent envelope is
\begin{equation}\label{eq:cmax_dc_main_v22}
    c^{gross,max,+,DC}_{i,p,s,r,t}(\xi)=\max\left\{0,\frac{fs^{eff}_{i,p,s,r,t}-s^{PB}_{i,t}}{1+\xi}\right\}.
\end{equation}
If the relevant numerator is negative, no positive policy cost is financeable under that scenario and regime.

\subsection{Feasibility frontier and threshold distances}

The baseline frontier is written in terms of the net marginal gross-capture object actually used in the fiscal-space equation:
\begin{equation}\label{eq:frontier_main_v10}
    \mathcal{F}^{\widetilde{MFC}}_{i,p,s,r,t}(\xi)=
    \left\{(g,\widetilde M):g\geq0,\ \widetilde M g\geq R^{req,gross}_{i,p,s,r,t}(\xi)\right\}.
\end{equation}
Here $\widetilde M$ denotes $\widetilde{MFC}^{gross}=MFC^{gross}-ac^{ME}-tr^{ME}-leak^{ME}$. This is a reduced-form or frozen-weight frontier; with weights recomputed as $g$ varies, it becomes the nonlinear set $\{g:fs^{eff,ch}(g)\geq R^{req}\}$ rather than a rectangular hyperbola in $(g,\widetilde M)$. The debt-consistent version substitutes $R^{req,gross,DC}$ for $R^{req,gross}$. Distance diagnostics comprise realized-coverage and support-headroom ratio families, which are never pooled; Appendix \ref{app:threshold_module} gives their complete definitions and use rules.
For adoption, $0<q^{prod,*}\leq\bar q_i$ means support is sufficient, whereas $q^{prod,*}>\bar q_i$ means the country ceiling fails; an absolute excess is a gap, not a ratio. Zero-threshold and floor-driven cases follow the unique operational definition in Section \ref{subsec:impossibility_score}.

\subsection{Monte Carlo feasibility probability}

For simulated draws indexed by $m$, the crossing probability is
\begin{equation}\label{eq:mc_prob_main_v10}
    \widehat{\pi}(\xi)=\frac{1}{M}\sum_{m=1}^{M}\mathbf{1}\{V^{gross,(m)}\geq1+\xi\}.
\end{equation}
The support-valid version conditions on support-valid and no-impossibility draws; the basic-valid version conditions only on arithmetic validity. The debt-consistent version additionally imposes Equation \eqref{eq:debt_requirement_main_v10}.


\section{Theoretical Results: A Single Feasibility Backbone}\label{sec:theory}

All feasibility thresholds transform the accounting backbone in Equation \eqref{eq:fs_gross_eff_v21}, under the preferred generalized gross-capture convention and with $g^{AI}_{i,s,t}=g^{AI,level}_{i,s,t}$. The baseline gross-cost condition is
\begin{equation}\label{eq:backbone_inequality_v17}
    g^{AI}_{i,s,t}\widetilde{MFC}^{gross}_{i,s,r,t}\geq R^{req,gross}_{i,p,s,r,t}(\xi).
\end{equation}

\begin{proposition}[Feasibility-frontier equivalence]\label{prop:frontier_equivalence}
Assume $c^{gross}_{i,p,t}>0$, $\widetilde{MFC}^{gross}_{i,s,r,t}>0$, and $g^{AI}_{i,s,t}>0$. For any buffer $\xi\geq0$, the following statements are equivalent under the non-debt baseline:
\begin{align}
    V^{gross}_{i,p,s,r,t}&\geq1+\xi,\\
    fs^{eff}_{i,p,s,r,t}&\geq(1+\xi)c^{gross}_{i,p,t},\\
    \widetilde{MFC}^{gross}_{i,s,r,t}g^{AI}_{i,s,t}&\geq R^{req,gross}_{i,p,s,r,t}(\xi),\\
    (g^{AI}_{i,s,t},\widetilde{MFC}^{gross}_{i,s,r,t})&\in\mathcal{F}^{\widetilde{MFC}}_{i,p,s,r,t}(\xi).
\end{align}
The debt-consistent version is obtained by replacing $R^{req,gross}$ with $R^{req,gross,DC}$.
\end{proposition}

\begin{proof}
Substitute $V^{gross}=fs^{eff}/c^{gross}$ and $fs^{eff}=\widetilde{MFC}^{gross}g^{AI}-ac^{net,fix}-tr^{ann}-leak^{fix}$, then move the fixed-cost terms to the right-hand side.
\end{proof}

\begin{proposition}[Unified threshold inversion]\label{prop:unified_inversion}
Define $R_{i,p,s,r,t}(\xi)\in\{R^{req,gross}_{i,p,s,r,t}(\xi),R^{req,gross,DC}_{i,p,s,r,t}(\xi)\}$. The following closed-form inversions assume that $\widetilde{MFC}^{gross}$ is fixed with respect to the variable being inverted; this holds in the historical reduced-form mode and in channel-based frozen-weight sensitivity exercises. If channel weights are recomputed endogenously as $g^{AI}$, $T$, or $q^{prod}$ changes, the implementation solves the corresponding root of Equation \eqref{eq:channel_revenue_function_v24} and reports the closed-form number only as a frozen-weight diagnostic. If $R>0$, $\widetilde{MFC}^{gross}>0$, $g^{AI}>0$, and $\phi^{Y,nom}_s>0$, then the required level-gain and gross-capture thresholds are
\begin{align}
    g^{AI,*}_{i,p,s,r,t}(\xi)&=\frac{R_{i,p,s,r,t}(\xi)}{\widetilde{MFC}^{gross}_{i,s,r,t}},\\
    MFC^{gross,*}_{i,p,s,r,t}(\xi)&=\frac{R_{i,p,s,r,t}(\xi)}{g^{AI}_{i,s,t}}+ac^{ME}_{i,s,r,t}+tr^{ME}_{i,s,r,t}+leak^{ME}_{i,s,r,t}.
\end{align}
Under the one-period static approximation $g^{AI}=\phi^{Y,nom}T$, the translation threshold is
\begin{equation}
    T^{*,static}_{i,p,s,r,t}(\xi)=\frac{R_{i,p,s,r,t}(\xi)}{\phi^{Y,nom}_s\widetilde{MFC}^{gross}_{i,s,r,t}}.
\end{equation}
For accumulated horizons, $T^*$ is obtained from Equation \eqref{eq:accumulated_threshold_condition_v22}; the static formula must not be used unless the run is explicitly labelled as one-period static. If $\widetilde{MFC}^{gross}\leq0$ with $R>0$, the $g^{AI,*}$ and $T^*$ inversions are undefined as finite positive thresholds and the cell is labelled non-positive net marginal fiscal capture.
\end{proposition}

\begin{proposition}[Adoption threshold and scale conversion]\label{prop:q_threshold_v10}
Suppose the baseline translation score is $S=\mathbf{1}\{q^{prod}>0,E^{prod}>0\}(\tilde E^{prod})^{\beta}(\tilde q^{prod,num})^{\gamma}$, $T=\min\{1,S/S^{frontier}\}$, $S^{frontier}$ is held fixed for the country-policy cell being inverted, and $\gamma>0$. First compute the required translation factor only under the existence condition
\begin{equation}\label{eq:valid_treq_theory_v20}
\begin{split}
    Valid^{T*}_{i,p,s,r,t}(\xi)=\mathbf{1}\{&R_{i,p,s,r,t}(\xi)>0,\ \phi^{Y,nom}_s>0,\ \widetilde{MFC}^{gross}_{i,s,r,t}>0,\\
    &\gamma>0,\ E^{prod}_{i,t}>0,\ \tilde E^{prod}_{i,t}>0,\ S^{frontier}_{s,t}>0\}.
\end{split}
\end{equation}
This list is identical to the existence conditions in Section \ref{subsec:adoption_threshold_existence} of Appendix \ref{app:threshold_module}.
For valid one-period static cells,
\begin{equation}\label{eq:T_req_theory_v10}
    T^{req,static}_{i,p,s,r,t}(\xi)=\frac{R_{i,p,s,r,t}(\xi)}{\phi^{Y,nom}_s \widetilde{MFC}^{gross}_{i,s,r,t}}.
\end{equation}
This object coincides with $T^{*,static}_{i,p,s,r,t}(\xi)$ of Proposition \ref{prop:unified_inversion}; the $req$ superscript emphasizes that it is a required value prior to the interiority check. If $Valid^{T*}_{i,p,s,r,t}(\xi)=0$ with positive required revenue, no finite adoption threshold is reported. For accumulated horizons, replace $T^{req,static}$ by the numerical solution of Equation \eqref{eq:accumulated_threshold_condition_v22}. If $T^{req,static}>1$, the cell is technology/fiscal-support insufficient: even the saturated translation factor $T=1$ cannot cross the threshold. If $0<T^{req,static}\leq1$, an interior threshold for normalized productive adoption is
\begin{equation}\label{eq:qtilde_star_theory_v10}
    \tilde q^{prod,*}_{i,p,s,r,t}(\xi)=
    \left[
    \frac{T^{req,static}_{i,p,s,r,t}(\xi)S^{frontier}_{s,t}}
    {(\tilde E^{prod}_{i,t})^{\beta}}
    \right]^{1/\gamma}.
\end{equation}
If $\tilde q^{prod,num}=\epsilon+(1-\epsilon)q^{prod}$, first compute the raw threshold on the original adoption scale,
\begin{equation}\label{eq:q_star_original_scale_v10}
    q^{prod,*,raw}_{i,p,s,r,t}(\xi)=\frac{\tilde q^{prod,*}_{i,p,s,r,t}(\xi)-\epsilon}{1-\epsilon},
\end{equation}
and then truncate at the lower support bound:
\begin{equation}\label{eq:q_star_truncated_v10}
    q^{prod,*}_{i,p,s,r,t}(\xi)=\max\{0,q^{prod,*,raw}_{i,p,s,r,t}(\xi)\}.
\end{equation}
If $q^{prod,*}=0$, the cell is not automatically classified as substantively feasible. It is labelled adoption-margin already satisfied only if the crossing also holds without relying solely on the numerical floor. If the crossing is created only by $\epsilon>0$, it is labelled floor-driven (unique operational definition in Section \ref{subsec:impossibility_score}). Impossibility diagnostics apply to the truncated original-scale object $q^{prod,*}$, not only to $\tilde q^{prod,*}$.
\end{proposition}

\begin{proof}
Solve the feasibility inequality for $\tilde q^{prod,num}$ holding $S^{frontier}$, $\phi^{Y,nom}_s$, $\widetilde{MFC}^{gross}$, and $\tilde E^{prod}$ fixed. The second equation is the inverse of the floor transformation used to construct $\tilde q^{prod,num}$.
\end{proof}
\begin{proposition}[No-double-counting convention]\label{prop:ndc_warning}
If $q^{prod}$ is simulated from preparedness and frictions,
\begin{equation}
    q^{prod}_{i,s,t}=h(A_{i,t},I_{i,t},Gap_{i,t}),
\end{equation}
then the baseline translation score may include $E^{prod}$ separately:
\begin{equation}
    S^{NDC}_{i,s,t}=\mathbf{1}\{q^{prod}_{i,s,t}>0,\ E^{prod}_{i,t}>0\}(\tilde E^{prod}_{i,t})^{\beta}(\tilde q^{prod,num}_{i,s,t})^{\gamma}.
\end{equation}
If, instead, the adoption equation also includes $E^{prod}$, then the baseline score must remove the direct exposure term and use $S^{q\text{-}only}=\mathbf{1}\{q^{prod}>0,\ E^{prod}>0\}(\tilde q^{prod,num})^{\gamma}$, retaining the zero-adoption/zero-exposure gate, or the specification must be labeled a partial double-counting correction rather than a no-double-counting specification.
\end{proposition}

\begin{proof}
If exposure enters both $h(\cdot)$ and $S$ directly, exposure affects $S$ through two channels. The convention above assigns exposure to the shock-translation margin and adoption to the diffusion margin, preventing the same exposure information from being multiplied twice.
\end{proof}

\begin{proposition}[Monotonicity and impossibility diagnostics]\label{prop:monotonicity_impossibility}
Holding fixed-cost terms, marginal-equivalent deductions, $S^{frontier}$, and all other arguments fixed as appropriate, and restricting attention to draws with $\phi^{Y,nom}_s\geq0$ so that $g^{AI}\geq0$ because $T_{i,s,t}\in[0,1]$, feasibility is weakly increasing in $\phi$, $T$, and $g^{AI}$ provided $\widetilde{MFC}^{gross}\geq0$. With $\widetilde{MFC}^{gross}<0$, feasibility is weakly decreasing in the AI shock because marginal-equivalent deductions exceed marginal revenue. Feasibility is weakly increasing in $MFC^{gross}$ because $g^{AI}\geq0$ on this support, and weakly decreasing in $c^{gross}$ and $\xi$. Draws with $\phi^{Y,nom}_s<0$ are retained economic failure draws for which the derivatives with respect to $T$ and $MFC^{gross}$ reverse sign; in particular, a draw with both $g^{AI}<0$ and $\widetilde{MFC}^{gross}<0$ would mechanically yield $\widetilde{MFC}^{gross}g^{AI}>0$ and is classified as an economic failure draw, not as a crossing. For components entering through $T=\min\{1,S/S^{frontier}\}$, monotonicity is weak and becomes flat after saturation at $T=1$. Hard-support violations---$q^{prod,*}>\bar q_i$, $q^{use,*}>\bar q_i$, $\tau^{AI,eff,*}>\bar\tau^{AI,eff}$, $\tau^{AI,*}>\bar\tau^{AI}$, and $g^{AI,*}>g^{AI,max}$---make the policy-scenario pair infeasible through that margin. By contrast, $MFC^{gross,*}>P90(MFC^{hist,gross,+})$ is historically extreme, a plausibility flag rather than a support violation.
\end{proposition}

\begin{proof}
The signs follow from differentiating Equation \eqref{eq:backbone_inequality_v17} under the stated ceteris-paribus conditions. The sign of the derivative with respect to $g^{AI}$ is the sign of $\widetilde{MFC}^{gross}$, while the derivative with respect to $MFC^{gross}$ is proportional to $g^{AI}$, which is non-negative on the restricted support $\phi^{Y,nom}_s\geq0$. Because $T$ is capped at one, changes in exposure or adoption that operate through $S$ have zero marginal effect after saturation. Hard-support diagnostics follow from the maintained supports of adoption, tax rates, and scenario-bounded AI growth; the historical fiscal-capture comparison is instead a plausibility diagnostic.
\end{proof}

The substantive implication is precise: disagreement about the size of the AI shock changes $\phi_s$; disagreement about diffusion changes $T$; disagreement about tax institutions changes $MFC^{gross}$ and the net fiscal-space deductions; disagreement about gross versus net program cost changes $c^{gross}$ or $c^{net}$; and debt prudence changes $s^{PB}$. The paper reports the frontier and the required thresholds rather than a single forecast.

\section{Parameter Assignment Protocol}\label{sec:parameters}

The main credibility threat is arbitrary calibration. The protocol therefore treats parameter assignment as a reproducible audit trail rather than as a set of discretionary point assumptions. Every input must be classified, sourced, transformed, bounded, and linked to an uncertainty treatment before it enters the feasibility frontier.

\subsection{Parameter classes and admissible uses}

Let the full parameter vector be
\begin{equation}
    \Theta=
    \Theta^{obs}\cup\Theta^{constructed}\cup\Theta^{lit}\cup\Theta^{calib}\cup\Theta^{stress}.
\end{equation}
The classes are mutually interpretable even when they overlap empirically. Observed inputs are country-year data. Constructed inputs are indices built from documented indicators. Literature-disciplined inputs are externally anchored ranges. Calibrated inputs are not directly observed and must be varied over grids or probability distributions. Stress inputs are deliberately extreme values used only to test the frontier.

\begin{longtable}{p{0.18\textwidth}p{0.30\textwidth}p{0.34\textwidth}}
\caption{Parameter hierarchy and assignment rules}\label{tab:parameter_hierarchy}\\
\toprule
Class & Examples & Assignment rule \\
\midrule
\endfirsthead
\toprule
Class & Examples & Assignment rule \\
\midrule
\endhead
Observed country data & GDP, population, eligible population, tax revenue/GDP, primary balance, informality, poverty lines, social protection expenditure & Use public data directly after unit harmonization; no scenario assignment unless the data are missing or stale \\
Constructed indices & AI preparedness, productive exposure, digital gap, fiscal capturability, policy cost net of existing transfers & Report raw indicators, normalization, weights, and aggregation formula; no hidden transformations \\
Literature-disciplined scenarios & Base AI shock $\phi_s$, exposure ranges, micro productivity relevance, diffusion speeds & Anchor low/central/high values in published studies; report the source-specific unit before conversion \\
Calibrated structural parameters & Translation elasticities $\alpha,\beta,\gamma$, informality elasticities $\rho_L,\rho_K,\rho_C$, leakage, AI-rent share, rent-capture rate & Use grids and Monte Carlo; no single preferred value unless estimated from data \\
Stress-test parameters & Disruptive AI shock, instant adoption, zero leakage, high rent taxation, high formalization effect & Use only to identify frontier requirements; do not present as baseline forecasts \\
\bottomrule
\end{longtable}

A result is baseline only if it uses no $\Theta^{stress}$ and is robust over the declared $\Theta^{calib}$ supports. Because calibrated inputs enter every run, the label means robustness over their supports, not their absence. Results requiring $\Theta^{stress}$ are stress feasibility.

\subsection{Unit discipline}

All fiscal variables are expressed as shares of the no-AI counterfactual GDP denominator $Y^0$ unless stated otherwise. Policy cost, effective fiscal space, fixed administrative costs, fixed transition costs, fixed leakage, and revenue gains must therefore be dimensionless ratios. Marginal-equivalent administrative, transition, or leakage terms are rates per unit of the AI shock and enter through $\widetilde{MFC}^{gross}$. The AI shock from the literature must be reported in its source unit and then converted through the mandatory chain
\begin{equation}
    \phi^{raw}_{s,u}\rightarrow \phi^Y_s=\kappa^{u\to Y}_s\phi^{raw}_{s,u}\rightarrow \phi^{Y,nom}_s\rightarrow g^{AI,level}_{i,s,H}.
\end{equation}
Micro studies such as \citet{noy2023experimental} and \citet{brynjolfsson2023generative} discipline the plausibility of productivity effects in exposed tasks, but they do not directly identify aggregate GDP growth or a recurring fiscal-base level increment.

\subsection{AI-productivity scenario ladder}

The externally disciplined scenario supports and deterministic closures are reported in Table \ref{tab:phi_numerical_closure}. These objects are reported across the paper and the reproducibility package: the source-unit raw supports in Table \ref{tab:phi_numerical_closure}; the scenario bridges $\kappa^{u\to Y}$ and nominal adjustments $\pi^{AI,Y}$ in the calibration registry; the post-bridge nominal supports in the primary-specification table of Section \ref{sec:empirical_results}; and the per-cell accumulated level objects $g^{AI,level}$ in the run traces.
A result may be called ``baseline feasible'' only if it holds after converting conservative or moderate TFP-equivalent raw supports into model-ready $\phi^{Y,nom}$ and then into $g^{AI,level}$. A result that requires the high G7 frontier case must be labelled ``frontier feasible.'' A result that requires the disruptive support must be labelled ``stress feasible'' and not interpreted as a forecast.

\subsection{Master calibration matrix}

The required parameter-reporting standard is in Table \ref{tab:master_calibration_matrix} (Appendix \ref{app:calibration_ledger}). The machine-readable country-parameter registry, including values, supports, sources, and approval status, is provided in the reproducibility package's \texttt{appendix\_i} calibration-registry files.

\subsection{Fiscal-conversion modes}

Two mutually exclusive modes obtain $MFC^{gross}$: channel mode derives it from effective rates and base-exposure weights; historical reduced-form mode draws it from the historical gross-capture distribution. A run must not treat both channel rates and independent $MFC^{gross}$ as primitives. If both are reported, one is the operative specification and the other a diagnostic decomposition.

\subsection{Distributional assignment rules}

When a parameter has a lower value $a$, a modal or central value $m$, and an upper value $b$, the baseline uncertainty distribution is a bounded PERT distribution. If $a=b$, the parameter is treated as fixed at that value and no beta draw is attempted:
\begin{equation}
    X=a+(b-a)Z,
    \qquad
    Z\sim Beta\left(1+4\frac{m-a}{b-a},\ 1+4\frac{b-m}{b-a}\right).
\end{equation}
With only a defensible support interval, use a uniform distribution and state that no central value is imposed. A lognormal is allowed only for strictly positive multiplicative parameters, such as leakage or administrative cost, after reporting its median and 90 percent interval. Bounded shares are truncated to their natural support.

Both independence and scenario-consistent dependence specifications are reported. Without an empirical correlation matrix, correlations are scenario assumptions, not estimates; Appendix \ref{app:uncertainty_module} gives the dependence structure.

\subsection{Data construction rules}

Observed variables are constructed from declared primary sources: GRD, ILOSTAT, ASPIRE, PIP or national poverty lines, AIPI, ITU and World Bank indicators excluded from AIPI, and WEO/DSA. Appendix \ref{app:calibration_ledger} gives the complete construction rules. Executed choices are recorded in the calibration registry and the Results data section.

\subsection{Calibration audit trail}

Each empirical table should be accompanied by a machine-readable calibration file with the following fields:
\begin{equation}
\begin{gathered}
    \{\text{parameter, country, year, raw source, raw value, transformation, unit,}\}\\
    \{\text{support, distribution, dependence block, scenario class}\}.
\end{gathered}
\end{equation}
This audit trail is not optional. It is what makes the framework reusable when future AI productivity estimates, adoption measures, or fiscal datasets change.

\subsection{Threshold-based rather than point-estimate reporting}

The empirical section reports continuous threshold objects before any categorical labels:
$V^{gross}$, the baseline accounting-fiscal index; $GAP^{gross}=fs^{eff}-c^{gross}$, the absolute GDP-share gap; $RGAP^{gross}=V^{gross}-1$, the relative structural gap; $g^*$, required AI-induced growth; $MFC^*$, required marginal fiscal conversion; $\phi^*$ and $T^*$, required base shock and translation; $\tau^{AI,*}$, required AI-rent capture; $q^*$, required productive adoption; $PB^0-PB^{stab}$, distance from the debt guardrail; and $\pi(\xi)$ and $\pi^{DC}(\xi)$, accounting and debt-consistent crossing probabilities.

This makes the framework robust even if AI forecasts change. Future evidence changes parameter values, not the logic of the frontier. Visual classification may use the baseline 10 and 25 percent bands, but conclusions are checked over the $\xi$ grid declared in Section \ref{sec:thresholds} and $\bar{\pi}\in\{0.50,0.75,0.90\}$.



\section{Impossibility, Placebo, and Falsification Diagnostics}\label{sec:impossibility_falsification}

\subsection{Infeasibility versus impossibility}\label{subsec:infeasibility_impossibility}

A policy-scenario pair can fail in two ways. It is \emph{below threshold} if feasible-support values exist but the current draw or scenario does not cross the frontier. It is \emph{impossible through a margin} if the required value of a bounded component lies outside the admissible support. The second classification is stronger and must be reported explicitly.

The maximum scenario-supported AI level shock is defined as
\begin{equation}
    g^{AI,max}_{s,t}=\phi^{Y,nom,max}_sD_{s,t}T^{max},\qquad T^{max}\equiv1,
\end{equation}
with $D_{s,t}=1$ in the static model. This notation separates growth support ($\phi^{Y,nom,max}$), dynamic translation support ($D_{s,t}$), and adoption/exposure translation support ($T^{max}$).

\begin{longtable}{p{0.27\textwidth}p{0.37\textwidth}p{0.25\textwidth}}
\caption{Impossibility diagnostics}\label{tab:impossibility_diagnostics_v10}\\
\toprule
Condition & Diagnosis & Reporting action \\
\midrule
\endfirsthead
\toprule
Condition & Diagnosis & Reporting action \\
\midrule
\endhead
$q^{prod,*}_{i,p,s,r,t}>\bar q_{i,t}$ & Required productive adoption exceeds the country ceiling & Adoption-impossible \\
$q^{use,*}_{i,p,s,r,t}>\bar q_{i,t}$ & Required use exceeds country adoption ceiling & Use-ceiling-impossible \\
$MFC^{gross,*}_{i,p,s,r,t}>P90(MFC^{hist,gross,+}_i)$ & Required gross fiscal capture exceeds historical high-capture benchmark & Historically demanding/extreme \\
$\tau^{AI,*}_{i,p,s,r,t}>\bar\tau^{AI}_{i,r,t}$ & Required AI-rent tax exceeds admissible policy bound & AI-rent-tax-impossible \\
$\tau^{AI,eff,*}_{i,p,s,r,t}>\bar\tau^{AI,eff}_{i,s,r,t}$ & Required effective AI-rent capture exceeds the effective bound $\bar\tau^{AI}\chi^{AI}$ & AI-rent-tax-impossible (effective bound) \\
$g^{AI,*}_{i,p,s,r,t}>g^{AI,max}_{s,t}$ & Required AI-induced growth exceeds scenario support & Technology-scenario-insufficient \\
$g^{AI,*}_{i,p,s,r,t}>g^{AI,max}_{s,H}$ & Required accumulated-horizon AI-induced growth exceeds the scenario-supported maximum defined in Appendix \ref{app:threshold_module} & Technology-scenario-insufficient (accumulated horizon) \\
$\widetilde{MFC}^{gross}_{i,s,r,t}\leq0$ and $R^{req,gross}_{i,p,s,r,t}>0$ & Net marginal fiscal capture is non-positive after marginal-equivalent deductions & Net-capture-impossible \\
$c^{gross}_{i,p,t}>c^{gross,max,+}_{i,p,s,r,t}$ & Program cost exceeds the financeable envelope & Policy-cost-too-large (threshold failure, not a support violation) \\
$(1+\xi)c^{gross}_{i,p,t}+s^{PB}_{i,t}>fs^{eff}_{i,p,s,r,t}$ & Debt-consistent guardrail fails & Fiscal-sustainability failure (threshold failure, not a support violation) \\
\bottomrule
\end{longtable}

The last two rows are threshold failures; only exogenous adoption, tax-rate, and scenario supports define impossibility through a margin under Section \ref{subsec:infeasibility_impossibility}. These diagnostics must not be hidden by a continuous score: a full UBI may fail because the shock is too small, adoption is too low, capture is historically extreme, the tax instrument is implausible, or gross cost is too large. Each failure has a distinct policy meaning.

\subsection{Impossibility score}\label{subsec:impossibility_score}

Because productive adoption is stock-like while use is a flow, a use threshold must be derived before a use-ceiling diagnostic is applied. For $\lambda_q>0$,
\begin{equation}\label{eq:q_use_star_v11}
    q^{use,*}_{i,p,s,r,t}(\xi)=\max\left\{0,
    \frac{q^{prod,*}_{i,p,s,r,t}(\xi)-(1-\lambda_q)q^{prod}_{i,s,t-1}}{\lambda_q}
    \right\}.
\end{equation}
If the inherited productive-adoption stock is already sufficient, the required contemporaneous use is set to zero and the use margin is classified as already satisfied. If $\lambda_q=0$, contemporaneous use cannot move productive adoption and the use-threshold inversion is not defined.

\subsection{Required persistence horizon}\label{subsec:h_star_inversion}

The inversions above hold the accumulation horizon fixed at $H$ and ask how large each
structural margin would have to be. The dual question is temporal: holding the scenario and
the frozen 2024 structure fixed, for how many years would the annual increment
$\phi^{Y,nom}_{s}T_{i,s}$ have to persist before the accumulated level gain covers the
buffered requirement? The required persistence horizon is
\begin{equation}\label{eq:h_star_inversion}
    H^{*}_{i,p,s,r}(\xi)=\min\left\{H\in\mathbb{N}:\;
    \widetilde{MFC}^{gross}_{i,s,r}\left[(1+\phi^{Y,nom}_{s}T_{i,s})^{H}-1\right]
    -F^{fix}_{i,p,s,r}\;\geq\;(1+\xi)\,c^{gross}_{i,p}(2024+H)\right\},
\end{equation}
where $c^{gross}_{i,p}(2024+H)$ carries the projected demographic eligibility ratio of the
evaluation year for age-conditioned instruments and is constant for gap-based instruments,
consistent with the endpoint cost convention. The debt-consistent analogue adds
$s^{PB+}_{i}$ to the requirement. When $c^{gross}$ is constant in $H$, the closed form is
$H^{*}=\lceil \ln(1+R^{req}/\widetilde{MFC}^{gross})/\ln(1+\phi^{Y,nom}T)\rceil$.

Three conventions discipline this statistic. First, $H^{*}$ is a requirement, not a
forecast: it states the number of years of scenario persistence that would be needed,
exactly as $MFC^{gross,*}$ states the capture that would be needed; it does not assert that
any calendar year will deliver it. Second, $H^{*}$ inherits the frozen-structure convention,
so it extends the maintained persistence assumption from the baseline $H=10$ to a variable
horizon; because long horizons stretch that assumption, reported values are censored at a
declared cap of $25$ years, beyond which the statistic is reported only as ``not reachable
within a defensible persistence horizon.'' Third, $H^{*}$ was added to the reporting set
after the pre-registered baseline was computed; it is a derived statistic of previously
frozen quantities, introduces no new calibrated parameter, and alters no pre-registered
result. The noncomputability flags of the preceding subsections apply unchanged, and the
Monte Carlo layer reports $\Pr(H^{*}\leq 10)$ and $\Pr(H^{*}\leq 25)$ under the same draw
discipline.

For each draw $m$, define basic denominator validity and support validity separately. Negative effective fiscal space and zero AI-induced level gains are not invalid draws; they are failure outcomes that remain in the denominator of all-draw probabilities:
\begin{equation}
    Valid^{basic,(m)}=\mathbf{1}\{c^{gross,(m)}>0\}.
\end{equation}
For threshold inversions, define the stricter computability indicator
\begin{equation}\label{eq:valid_threshold_operational_v25}
\begin{split}
    Valid^{threshold,(m)}=\mathbf{1}\{&c^{gross,(m)}>0,
    \ R^{req,gross,(m)}>0,
    \ g^{AI,level,(m)}>0,
    \ \widetilde{MFC}^{gross,(m)}>0,
    \ \phi^{Y,nom,(m)}>0,\\
    &\gamma^{(m)}>0,
    \ S^{frontier,(m)}>0,
    \ E^{prod,(m)}>0,
    \ \tilde E^{prod,(m)}>0\}.
\end{split}
\end{equation}
Draws with $R^{req,gross,(m)}\leq0$ are requirement-already-covered on the fixed-cost margin, not inverted. Their separately reported share is trivially feasible there, so exclusion from threshold computability is not fragility in $\pi^{threshold|basic}$. AI-rent thresholds distinguish effective and statutory computability:
\begin{align}
    Valid^{AIrent,eff,(m)}
    &=\mathbf{1}\{g^{AI,level,(m)}>0,\ \lambda^{AI,(m)}>0\},\\
    Valid^{AIrent,stat,(m)}
    &=\mathbf{1}\{g^{AI,level,(m)}>0,\ \lambda^{AI,(m)}>0,\ \chi^{AI,(m)}>0\}.
\end{align}
To prevent the numerical floor from creating substantive feasibility, define
\begin{equation}\label{eq:floor_driven_operational_v25}
    FloorDriven^{(m)}=\mathbf{1}\{q^{prod,*,(m)}=0,\ \epsilon^{(m)}>0,\ \text{crossing fails when }\epsilon=0\}.
\end{equation}
Then define support validity as
\begin{equation}
\begin{split}
    Valid^{support,(m)}=\mathbf{1}\{&Valid^{threshold,(m)}=1,
    \ q^{prod,*,(m)}\leq\bar q^{(m)},
    \ q^{use,*,(m)}\leq\bar q^{(m)},\\
    &\left(\lambda^{AI,(m)}=0\ \text{or}\ \tau^{AI,eff,*,(m)}\leq\bar\tau^{AI,eff,(m)}\right),
    \ g^{AI,*,(m)}\leq g^{AI,max,(m)},
    \ FloorDriven^{(m)}=0\}.
\end{split}
\end{equation}
A non-floor-driven $q^{prod,*}=0$ means the adoption margin is already satisfied and support-valid. The effective AI-rent bound $\bar\tau^{AI,eff}=\bar\tau^{AI}\chi^{AI}$ subsumes the statutory check: $\tau^{AI,eff,*}\leq\bar\tau^{AI}\chi^{AI}$ implies $\tau^{AI,*}\leq\bar\tau^{AI}$ when $\chi^{AI}>0$, while $\chi^{AI}=0$ with $\tau^{AI,eff,*}>0$ is a support violation.
The validity shares describe the chain from all draws to basic arithmetic validity, threshold computability, and bounded-support validity:
\begin{align}
    \pi^{basic}&=\frac{1}{M}\sum_{m=1}^{M}Valid^{basic,(m)},\\
    \pi^{threshold|basic}&=\frac{\sum_m Valid^{threshold,(m)}}{\sum_m Valid^{basic,(m)}}
    \quad\text{if }\sum_m Valid^{basic,(m)}>0,\\
    \pi^{support|basic}&=\frac{\sum_m Valid^{support,(m)}}{\sum_m Valid^{basic,(m)}}
    \quad\text{if }\sum_m Valid^{basic,(m)}>0.
\end{align}
The bounded-support impossibility share is reported conditionally on threshold computability,
\begin{equation}
    IS^{support|threshold}_{i,p,s,r,t}
    =1-\frac{\sum_m Valid^{support,(m)}}{\sum_m Valid^{threshold,(m)}}
    \quad\text{if }\sum_m Valid^{threshold,(m)}>0,
\end{equation}
while $1-\pi^{basic}$ is the denominator-invalid share and $\pi^{support|basic}$ remains aggregate reporting. Draws with $g^{AI}=0$ and $c^{gross}>0$ are basic-valid non-crossings; $c^{gross,(m)}\leq0$ draws are denominator-invalid, excluded from $\widehat{\pi}^{basic}$ but retained in all-draw probability with zero crossing. Report all-, basic-valid-, and support-valid-draw probabilities. High all-draw but low support-valid probability signals non-computability or support violations, separated by the intermediate layer.

\subsection{Temporal placebo}

The placebo protocol has two layers. The first is a mechanical accounting-null check. In pre-GenAI years, the GenAI-specific productivity shock is set to zero,
\begin{equation}
    \phi^{GenAI}_{s,t}=0\qquad \text{for accounting-null placebo years.}
\end{equation}
With the zero-adoption gate active, this implies $T^{GenAI}=0$ and $g^{AI,level}=0$ for the GenAI channel. This null check verifies that the code does not manufacture AI fiscal space mechanically, but it is not interpreted as an informative historical placebo because $0$ minus positive deployment costs is arithmetically negative.

The second layer is the informative pre-GenAI ICT placebo. For years such as 2010--2018, the implementation feeds the translation block with historical digitalization or ICT-adoption proxies rather than GenAI adoption:
\begin{equation}
    S^{ICT}_{i,t}=\mathbf{1}\{q^{ICT}_{i,t}>0,\ E^{ICT}_{i,t}>0\}(\tilde E^{ICT}_{i,t})^{\beta_{ICT}}(\tilde q^{ICT,num}_{i,t})^{\gamma_{ICT}},
    \qquad
    T^{ICT}_{i,t}=\min\left\{1,\frac{S^{ICT}_{i,t}}{S^{ICT,frontier}_{t}}\right\}.
\end{equation}
The placebo shock is then
\begin{equation}
    g^{ICT,placebo}_{i,t}=\phi^{ICT,placebo}_{t}T^{ICT}_{i,t},
\end{equation}
where $\phi^{ICT,placebo}$ is a documented growth-accounting ICT/digitalization analogue, not a GenAI shock, with comparability defaults $\beta_{ICT}=\beta$ and $\gamma_{ICT}=\gamma$. A strong flag marks broad pre-GenAI placebo feasibility of AI-funded UBI/PEN/GMI under the same thresholds:
\begin{equation}
    PlaceboFail^{ICT}_{i,p,t}=\mathbf{1}\{V^{gross,ICT}_{i,p,t}\geq1\ \text{for broad policy cells in pre-GenAI years}\}.
\end{equation}
The test asks whether ordinary digitalization, growth, or buoyancy is being relabelled as AI fiscal space; it is diagnostic, not causal. Its known-outcome benchmark is that digitalization in 2010--2018 financed universal social protection nowhere in the region, so broad feasibility contradicts history, not merely a prior. Validity requires the full era-consistent pipeline---historical capture distributions, policy costs, and informality---not translation alone.

\subsection{Sectoral negative controls}

For sectors with very low AI exposure, the translation score should be low. Let $\mathcal{K}^{NC}$ be a set of low-exposure sectors. A negative-control diagnostic reports
\begin{equation}
    T^{NC}_{i,s,t}=\sum_{k\in\mathcal{K}^{NC}}\omega^{NC}_{i,k,t}T_{i,k,s,t}.
\end{equation}
The weights $\omega^{NC}_{i,k,t}$ are employment or value-added shares within the negative-control set and satisfy $\sum_{k\in\mathcal{K}^{NC}}\omega^{NC}_{i,k,t}=1$.
The expectation is not that $T^{NC}=0$, but that it remains substantially below the country aggregate in scenarios where the aggregate is driven by high-exposure sectors. If negative-control sectors generate high translation scores, the exposure/adoption construction is too permissive. These placebo and negative-control diagnostics are diagnostic stress tests, not substitutes for a causal treatment-effect design \citep{bertrand2004trust,athey2017state,lipsitch2010negative}.

\subsection{Leave-one-source-out and channel-ablation protocol}

The replication package reports leave-one-source-out and channel-ablation checks:
\begin{itemize}
    \item AI-productivity source: Acemoglu-only, OECD-only, and high/stress-only exclusions.
    \item Exposure source: IMF/OECD exposure only, World Bank--ILO exposure only, and country-sector exposure only where available.
    \item Fiscal channel: no AI-rent tax, no formalization effect, no labor augmentation, high import leakage, and no base-broadening reform.
    \item Cost convention: gross-cost baseline, net-cost robustness, and debt-consistent requirement.
\end{itemize}
If a positive feasibility result depends on a single source or a single fiscal channel, the result is reported as source-dependent rather than robust.


\section{Robustness Strategy}\label{sec:robustness}

\subsection{Robust-decision logic}

The framework reports thresholds, not a single headline forecast. A result is robust only if it survives variation in AI productivity, productive adoption, fiscal capture, informality, leakage, program cost, and debt guardrails.

\subsection{Types of uncertainty}

The implementation separates four uncertainty layers: scenario uncertainty, structural uncertainty, parameter uncertainty, and data uncertainty. Scenario uncertainty is handled by $\phi^{Y,nom}_s$ and fiscal regime $r$. Structural uncertainty is handled by alternative formulas. Parameter uncertainty is handled by Monte Carlo draws. Data uncertainty is handled by reporting observed-source alternatives and harmonized-data alternatives.

\subsection{Specification robustness}

Specification checks include effective versus gross fiscal conversion, gross-cost versus net-cost policy denominators, debt-consistent versus non-debt thresholds, no-AI-rent-tax regimes, high-leakage regimes, and no-formalization regimes.

\subsection{Parameter-grid robustness}

The deterministic grid includes
\begin{align}
    \phi^{Y,nom}_s &\in \{\phi^{low,Y,nom},\phi^{mid,Y,nom},\phi^{high,Y,nom},\phi^{stress,Y,nom}\},\\
    \lambda^{AI} &\in \{\lambda^{AI}_L,\lambda^{AI}_M,\lambda^{AI}_H\},\\
    \lambda_q &\in \{0.15,0.30,0.50\}\quad\text{baseline; }\lambda_q=1\text{ stress},\\
    \xi &\in \{0,0.05,0.10,0.25,0.50\}.
\end{align}
The symbol $\lambda^{AI}$ is reserved for AI-rent share. The symbol $\lambda_q$ is reserved for adoption-diffusion speed. They are not interchangeable.

\subsection{Monte Carlo uncertainty propagation}\label{subsec:mc_main}

Let $\Theta^{(m)}$ denote draw $m$ of the uncertain-input vector. The full vector must state which variables are observed and fixed and which are stochastic. The implementation must choose one fiscal-conversion mode and one adoption mode for a given run.

Observed-adoption mode draws $q^{prod}$ from adoption or measurement-error distributions and treats $A,I,Gap$ as controls. Anchored simulated-adoption mode takes $A,I,Gap,\lambda_q,\nu_A,\nu_I,\nu_G,q^{use,target}_{i,t_0}$ as primitives; logit-inverted $\mu_{i,s,t_0}$ is not an independently identified structural coefficient. Structural-intercept robustness instead draws a common or hierarchical intercept without target inversion. Simulated-adoption modes derive $q^{use}$ and $q^{prod}$ endogenously and never draw $q^{prod}$ independently.

Channel mode derives gross capture from tax bases rather than drawing it. Threshold runs declare frozen versus endogenous weights; only the former use closed forms, while the latter solve Equation \eqref{eq:channel_revenue_function_v24}. Appendix \ref{app:uncertainty_module} gives both mode vectors and the complete AI-rent primitives $\tau^{AI},\chi^{AI,dom},\psi^{shift,tech},\chi^{AI,compliance},\chi^{AI,enforce}$, with derived $\tau^{AI,eff}=\tau^{AI}\chi^{AI}$; historical reduced-form mode instead draws gross capture.
Capturability elasticities $\rho_j$ are outside $\Theta_{hist}$. Labelled alternatives add their parameters to the audit trail: $\psi^{K,rate}$ for rate-side shifting and $\epsilon^{ero,j}$ for behavioral erosion.
Historical-mode channel rates may be descriptive but do not also construct $MFC^{gross}$. The run declares observed or simulated adoption; only observed-adoption mode treats $q^{prod}$ as primitive. Fixed GDP-share costs and marginal-equivalent deductions each occupy one margin. Nominal $AC$, $TR$, and $Leak$ are derived from GDP, never drawn independently alongside normalized objects.

For each draw, effective fiscal space is computed directly from net marginal gross capture and fixed costs; when $g^{AI,level,(m)}>0$, the per-unit object $MFC^{eff,(m)}=fs^{eff,(m)}/g^{AI,level,(m)}$ is derived rather than drawn. The complete draw-level specification is in Appendix \ref{app:uncertainty_module}.
A draw with $fs^{eff,(m)}<0$ is not invalid. It is an economically meaningful failure draw: administrative costs, transition costs, and leakage exceed gross AI-generated fiscal space. Such draws remain in the all-draw feasibility probability and lower the crossing probability.
Arithmetic, computability, support, numerical-floor, and failure semantics follow the canonical conservative reporting chain in Section \ref{subsec:impossibility_score}.
The gross feasibility ratio is $V^{gross,(m)}=fs^{eff,(m)}/c^{gross,(m)}$. The net ratio $V^{net,(m)}=fs^{eff,(m)}/c^{net,(m)}$ is computed only for draws with $c^{net,(m)}>0$.
If $c^{net,(m)}=0$, the draw is classified as already covered by existing spending on the net-cost margin. $V^{gross}$ remains the baseline reporting object.

\subsection{Monte Carlo output layers}

For each country-policy-scenario-regime cell, the percentile and conditional-probability layers are those specified in Appendix \ref{app:uncertainty_module}; the canonical all-draw crossing grid is Equation \eqref{eq:mc_prob_main_v10}.
If the denominator is zero, the corresponding conditional probability is reported as undefined/NA and the relevant valid-share, $\pi^{basic}$ or $\pi^{support}$, is reported as zero. This distinguishes ``no valid support exists'' from ``valid support exists but no draw crosses the threshold.'' The debt-consistent version uses Equation \eqref{eq:debt_requirement_main_v10}.

\subsection{Distributional choices and dependence}

The baseline distribution for low-central-high calibrated scalar parameters is bounded PERT. Triangular and truncated-uniform draws are robustness checks, not the baseline convention. Bounded share variables may also be drawn from beta distributions when moments are available. Dependence between preparedness, adoption, informality, and fiscal capacity may be implemented through correlated latent shocks. If dependence is imposed, the paper reports Spearman and tornado diagnostics as the baseline sensitivity outputs and reserves classical Sobol indices for an independent-input diagnostic run or a dependence-aware sensitivity method; Appendix \ref{app:uncertainty_module} gives the scenario-consistent dependence structure.

\subsection{Global sensitivity and stress tests}

Report Spearman rank correlations, one-at-a-time tornado ranges, and grouped block sensitivity for AI productivity, adoption, fiscal capture, informality, leakage, and cost \citep{mckay1979lhs,morris1991factorial,saltelli2008global,sobol2001global}. Classical Sobol indices are reserved for independent-input diagnostic runs or dependence-aware sensitivity methods; Appendix \ref{app:uncertainty_module} gives the full formulas and protocol.

\subsection{Numerical precision and reproducibility}

The replication package should publish the country-year dataset, parameter matrix, random seed, draw count, distributional supports, cleaning rules, and exact code used to generate tables. Simulations should be rerun with at least two seeds or with convergence diagnostics for the main crossing probabilities.

\subsection{Placebo and negative-control checks}

Placebo and negative-control outputs are required tables, not optional prose. They include the temporal placebo, low-exposure sector negative controls, leave-one-source-out results, and channel ablations.

\section{Empirical Design and Minimum Pilot Application}\label{sec:empirical_design}

The deliberately small application uses Peru as primary case and Chile, Colombia, and Mexico as comparators. It does not claim regional representativeness; it shows that common equations can use observed data across fiscal and institutional settings and that feasibility is driven by transparent rather than hidden inputs.

\subsection{Pilot country sample}

\begin{table}[H]
\centering
\caption{Minimum pilot sample and identification logic}
\label{tab:pilot_country_logic}
\small
\begin{tabular}{p{0.18\textwidth}p{0.64\textwidth}}
\toprule
Country & Reason for inclusion \\
\midrule
Peru & Primary case; high informality, low tax-to-GDP ratio, relevant for social protection expansion and fiscal-capacity constraints. \\
Chile & Higher institutional and fiscal capacity within the same region; useful upper comparator for preparedness and formalization. \\
Colombia & Andean comparator with larger social-policy system and persistent informality. \\
Mexico & Large economy with deeper productive integration, low tax-to-GDP ratio by OECD standards, and distinct digital/productive structure. \\
\bottomrule
\end{tabular}
\end{table}

Three pilot countries are upper-middle-income and Chile is high income. The sample spans informality and fiscal capacity only within this middle-to-high band, not the low-income, low-capacity corner where impossibility may bind most; a low-income case is the natural next application.

\subsection{Observed pilot anchors}

The harmonization year is $t_0=2024$: for a 2024 observation, use the 2024 value; carry earlier values forward with a flag. A newer value such as 2026 requires either all macro anchors updated to 2026 or a non-harmonized source-specific robustness flag. Baseline comparisons must not silently mix years.

The preliminary seed values in the methodological draft have been superseded by the official observed anchors constructed by the reproducibility pipeline and reported in Table \ref{tab:results_anchors} in the results section. Those official anchors draw on GRD, OECD Revenue Statistics, WDI, AIPI, and national informality sources, under the harmonization rules above \citep{imf2024aipi,worldbank2026wdi,oecd2025revenue,oecd2025peru,worldbank2026humancapital,inechile2025ene,dane2026informalidad,inegi2024enoe}.

\subsection{Policy instruments and cost equations}

For each country $i$ and policy $p$, the annual policy-cost ratio $c_{i,p,t}$ divides the nominal equations in Section \ref{sec:policy_costs} by no-AI GDP $Y^0_{i,t}$. Required data are older population and pension benefit for PEN; eligible population and poverty line for MUT/PBI; household-income distribution and minimum line for GMI; and total population for UBI.

ASPIRE supplies existing expenditure and coverage; PIP, SEDLAC, and national surveys supply poverty lines and gaps, with detail in Appendix \ref{app:calibration_ledger}. A full UBI is retained only as a benchmark. The paper should expect universal social pensions, targeted or guaranteed minimum income instruments, and partial basic income instruments to cross feasibility thresholds before a full UBI benchmark, with the GMI ranking subject to the ideal-targeting caveat of Section \ref{sec:policy_costs}.

\subsection{AI exposure, adoption, and fiscal conversion in the pilot}

The pilot uses the World Bank--ILO exposure ranges defined in Section \ref{subsec:prod_disp_exposure_main}. Country differentiation through AIPI, informality, digital gaps, and productive adoption follows the logistic and diffusion module in Section \ref{subsec:logistic_adoption_main}.
The official tables report, for every country-policy-scenario-regime cell,
\begin{equation}
    \left(A_i,E^{prod}_i,q^{prod}_{i,s,t},I_i,MFC^{gross}_{i,s,r,t},fs^{eff}_{i,p,s,r,t},c_{i,p,t},V_{i,p,s,r,t}\right).
\end{equation}
The pilot's fiscal-conversion mode is declared once: all $V$ columns are computed in the channel-based mode, which represents reform regimes natively, while the historical gross-capture distribution enters through the plausibility flag and through a labelled historical reduced-form robustness run; the two modes are never mixed inside one cell. Results from this design are reported in Section \ref{sec:empirical_results}.
The sequence was implemented in full in the tagged reproducibility package described in the data and code availability statement, and its outputs are reported in Section \ref{sec:empirical_results}.
This empirical design is intentionally modular. If a reader disagrees with the AI shock, adoption speed, tax-base channel, or policy-cost benchmark, the framework allows that block to be replaced without changing the threshold logic.

\subsection{Expected qualitative pattern}\label{sec:expected}

Before empirical implementation, the expected pattern is:
\begin{itemize}
    \item Full UBI is likely to require implausibly high AI growth or fiscal capture in many emerging economies.
    \item Universal social pensions, minimum universal transfers, and targeted guaranteed minimum income schemes should cross thresholds earlier, with the GMI comparison subject to the ideal-targeting caveat of Section \ref{sec:policy_costs}: its early crossing partly reflects the poverty-gap costing convention and must be re-examined under the $\theta^{target}$ loading.
    \item Countries with stronger AI preparedness and stronger fiscal capacity should cross thresholds first.
    \item Countries with high exposure, low preparedness, high informality, and low fiscal capacity should score high on the pressure-without-capacity index.
    \item Fiscal reform may matter as much as AI optimism: a moderate increase in marginal fiscal conversion can reduce the required AI growth threshold substantially.
\end{itemize}

These are hypotheses, not results.

\section{Empirical Results}\label{sec:empirical_results}

This section reports the official results of the pilot application to Peru, Chile, Colombia,
and Mexico. All numbers are generated by the tagged reproducibility package described in the
data and code availability statement; no value in this section is transcribed by hand, and
every table is read directly from the verified output of the pipeline. The narrative follows
the pre-registered reporting order: policy costs, AI-induced fiscal space, feasibility
ratios, historical plausibility, debt guardrails, Monte Carlo uncertainty, threshold
inversions with required persistence horizons, diagnostics, and limits.

\subsection{Data, snapshot, and calibration}\label{subsec:results_data}

The empirical base is a frozen, hash-verified snapshot of public sources (World Bank WDI and
PIP, OECD Revenue Statistics, IMF WEO and AIPI, ILOSTAT, PWT, ITU, Ookla, UN WPP, and the
UNU-WIDER GRD as a contrast series), harmonized to the 2024 anchor year. The digital-access
gap is constructed exclusively from indicators verified to be outside the IMF AI Preparedness
Index (mobile-network coverage and median connection speeds), which enforces the
no-double-counting rule of Section~\ref{subsec:ndc_translation_main} by construction. All
calibrated parameters were registered, source-disciplined, and signed by the author before
any official run; the full assignment is reported in the master calibration matrix in the
appendix. Table~\ref{tab:results_anchors} reports the observed country anchors and
Table~\ref{tab:results_primary_spec} the frozen primary specification.

\begingroup
\begin{table}[!htbp]
\centering
\caption{Observed 2024 country anchors. GDP is in USD billions; fiscal ratios, informality, and adoption are percentages. The complete anchor matrix is provided in the replication package.}\label{tab:results_anchors}
\fontsize{7.2}{8.6}\selectfont
\setlength{\tabcolsep}{2.4pt}
\begin{tabular}{@{}lrrrrrrrr@{}}
\toprule
Country & \shortstack{GDP\\(USD bn)} & \shortstack{Tax rev.\\(\% GDP)} & \shortstack{Total rev.\\(\% GDP)} & \shortstack{Debt\\(\% GDP)} & AIPI & \shortstack{Digital\\gap} & \shortstack{Informality\\(\%)} & \shortstack{Adoption\\target (\%)} \\
\midrule
Chile & 329.3 & 20.5 & 23.9 & 41.7 & 0.586 & 0.089 & 27.5 & 13.2 \\
Colombia & 420.5 & 19.9 & 28.3 & 61.0 & 0.489 & 0.239 & 56.1 & 11.0 \\
Mexico & 1830.5 & 18.3 & 25.0 & 59.1 & 0.532 & 0.351 & 50.6 & 12.0 \\
Peru & 291.8 & 16.3 & 18.7 & 32.1 & 0.491 & 0.265 & 67.5 & 11.1 \\
\bottomrule
\end{tabular}
\end{table}
\endgroup

\begingroup
\begin{table}[!htbp]
\centering
\caption{Primary AI scenarios and fiscal regimes.}\label{tab:results_primary_spec}
\small
\setlength{\tabcolsep}{7pt}
\begin{tabular}{@{}lr@{\hspace{3em}}lr@{}}
\toprule
Scenario & \shortstack{Annual AI shock\\(\% of GDP)} & Regime & \shortstack{Median $MFC^{gross}$\\(\%)} \\
\midrule
Low & 0.050 & r0 & 8.92 \\
Moderate & 0.425 & r1 & 9.39 \\
High & 0.306--0.421 & r2 & 9.82 \\
Disruptive stress & 1.500 & r3 & 8.72 \\
 &  & r4 & 10.93 \\
\bottomrule
\multicolumn{4}{@{}p{0.88\linewidth}@{}}{\footnotesize The high scenario reports the country range after the source-unit and nominal-output bridges; other scenarios are point supports.} \\
\end{tabular}
\end{table}
\endgroup


\subsection{Policy costs}\label{subsec:results_costs}

Table~\ref{tab:results_policy_costs} reports gross annual costs at the 2024 anchor and at
the 2034 endpoint under the demographic-projection convention. Three features matter for
everything that follows. First, the cost gradient across instruments spans two orders of
magnitude: gap-based guaranteed minimum income is the cheapest instrument where poverty gaps
are small (about $0.15$ percent of GDP in Chile under ideal targeting), while the full UBI
benchmark reaches roughly a fifth of GDP in Peru. Second, the expected cost ordering across
instruments is not universal: Chile's universal pension guarantee is several times more
expensive than its poverty gap, which reverses the ex-ante ordering of
Section~\ref{sec:expected} for that country and illustrates that instrument rankings are a
country-structure outcome rather than a design constant. Third, ageing raises age-conditioned
costs materially by the endpoint, which the demographic-ratio convention captures without
projecting monetary levels. For the primary case, the guaranteed-minimum-income cost was
additionally validated in a microdata-strong extension using the ENAHO 2024 household survey
at the same PIP poverty line: the official grouped-data cost ($2.79\%$ of GDP at ideal
targeting) lies between the income-based ($2.57\%$) and expenditure-based ($3.06\%$)
microdata costs, and no cell classification changes under either variant (amendment registry,
reproducibility package). For Chile, whose ideal-targeting GMI drives the headline conditional
result, the same validation using the CASEN 2024 survey reproduces the official grouped-data
cost to within $0.4\%$ at the same line ($0.1469\%$ versus $0.1464\%$ of GDP), again with no
classification changes.

\begingroup
\singlespacing
\fontsize{8}{9.6}\selectfont
\setlength{\tabcolsep}{3pt}
\begin{longtable}{@{}p{0.12\linewidth}p{0.12\linewidth}p{0.13\linewidth}r p{0.36\linewidth}@{}}
\caption{Gross annual policy costs at the 2024 anchor. GMI rows report ideal and loaded targeting separately; net costs equal gross costs under the pure-layering baseline.}\label{tab:results_policy_costs}\\
\toprule
Country & Instrument & Variant & \shortstack{Gross cost\\(\% GDP)} & Endpoint convention \\
\midrule
\endfirsthead
\caption[]{Gross annual policy costs at the 2024 anchor. GMI rows report ideal and loaded targeting separately; net costs equal gross costs under the pure-layering baseline. (continued)}\\
\toprule
Country & Instrument & Variant & \shortstack{Gross cost\\(\% GDP)} & Endpoint convention \\
\midrule
\endhead
Peru & PEN & -- & 0.5574 & WPP eligible-share ratio \\
Peru & GMI & Ideal & 2.7872 & Fixed aggregate cost \\
Peru & GMI & Loaded & 4.1808 & Fixed aggregate cost \\
Peru & MUT & -- & 5.3395 & WPP eligible-share ratio \\
Peru & PBI & -- & 7.9695 & WPP eligible-share ratio \\
Peru & UBI & -- & 21.3579 & WPP eligible-share ratio \\
Chile & PEN & -- & 3.1044 & WPP eligible-share ratio \\
Chile & GMI & Ideal & 0.1464 & Fixed aggregate cost \\
Chile & GMI & Loaded & 0.2195 & Fixed aggregate cost \\
Chile & MUT & -- & 2.9040 & WPP eligible-share ratio \\
Chile & PBI & -- & 4.8517 & WPP eligible-share ratio \\
Chile & UBI & -- & 11.6161 & WPP eligible-share ratio \\
Colombia & PEN & -- & 0.9046 & WPP sex- and age-specific eligibility \\
Colombia & GMI & Ideal & 2.9859 & Fixed aggregate cost \\
Colombia & GMI & Loaded & 4.4789 & Fixed aggregate cost \\
Colombia & MUT & -- & 4.7456 & WPP eligible-share ratio \\
Colombia & PBI & -- & 7.4321 & WPP eligible-share ratio \\
Colombia & UBI & -- & 18.9824 & WPP eligible-share ratio \\
Mexico & PEN & -- & 1.5996 & WPP eligible-share ratio \\
Mexico & GMI & Ideal & 1.0120 & Fixed aggregate cost \\
Mexico & GMI & Loaded & 1.5180 & Fixed aggregate cost \\
Mexico & MUT & -- & 3.9531 & WPP eligible-share ratio \\
Mexico & PBI & -- & 5.9054 & WPP eligible-share ratio \\
Mexico & UBI & -- & 15.8125 & WPP eligible-share ratio \\
\bottomrule
\end{longtable}
\endgroup


\subsection{Baseline feasibility}\label{subsec:results_baseline}

Table~\ref{tab:results_vgross} reports the gross-cost feasibility ratio $V^{gross}$ for the
primary-case moderate-scenario block and all crossing cells at the prudential buffer
$\xi=0.10$; the full four-country grid (480 cells) is provided in the reproducibility package. The central
result is stark: under the conservative and moderate scenario supports, no instrument in any
country reaches the accounting threshold in any fiscal regime; feasibility ratios remain
below $0.7$ throughout. Crossings exist only under the disruptive stress support, and they
are concentrated and instructive: Chile's guaranteed minimum income crosses in every regime,
including the status quo, and passes the debt guardrail with feasibility ratios between
$2.6$ and $5.8$ across the loaded- and ideal-targeting variants ($3.9$--$5.8$ under ideal
targeting; $2.6$--$3.9$ under the loaded companion); Peru's universal social pension crosses only marginally ($V$ between $1.0$
and $1.4$), fails the debt guardrail in every regime, and its most marginal crossing (status quo regime, $V=1.02$) does not survive the raw labor-share robustness extension of Appendix \ref{app:labor_share_module}, falling to $V=0.98$. All stress-support results carry
the mandated label: they are conditional on a scenario deliberately beyond the baseline
macroeconomic literature and are not forecasts.

\begingroup
\fontsize{8}{9.6}\selectfont
\setlength{\tabcolsep}{1.2pt}
\let\tableunderscore\_
\renewcommand{\_}{\tableunderscore\allowbreak}
\begin{longtable}{@{}
>{\raggedright\arraybackslash}p{0.06\linewidth}
>{\raggedright\arraybackslash}p{0.23\linewidth}
>{\raggedright\arraybackslash}p{0.10\linewidth}
>{\raggedright\arraybackslash}p{0.07\linewidth}
>{\raggedright\arraybackslash}p{0.19\linewidth}
@{\hspace{0.35em}}
>{\raggedright\arraybackslash}p{0.12\linewidth}
>{\raggedright\arraybackslash}p{0.18\linewidth}
@{}}
\caption{Table 4. Baseline $V^{gross}$: four-country moderate-support null and all crossing cells. Headline cells; the full grid is provided in the Online Appendix and the reproducibility package.}
\label{tab:results_vgross}\\
\toprule
country\_id & policy\_variant\_id & scenario\_id & regime\_id & v\_gross & crosses\_v1 & cell\_result\_class \\
\midrule
\endfirsthead
\caption[]{Table 4. Baseline $V^{gross}$: four-country moderate-support null and all crossing cells. Headline cells; the full grid is provided in the Online Appendix and the reproducibility package. (continued)}\\
\toprule
country\_id & policy\_variant\_id & scenario\_id & regime\_id & v\_gross & crosses\_v1 & cell\_result\_class \\
\midrule
\endhead
\multicolumn{7}{l}{\textit{Panel A: all countries, moderate scenario, r0, primary policy variants}} \\
\midrule
PER & PEN & mid & r0 & 0.0781898593701847 & False & not\_feasible \\
PER & GMI:GMI\_ideal\_aggregate & mid & r0 & -0.0525921743988333 & False & not\_feasible \\
PER & MUT & mid & r0 & 0.0191207387806886 & False & not\_feasible \\
PER & PBI & mid & r0 & 0.0164818464107766 & False & not\_feasible \\
PER & UBI & mid & r0 & 0.0068349135400345 & False & not\_feasible \\
CHL & PEN & mid & r0 & 0.0324037073756413 & False & not\_feasible \\
CHL & GMI:GMI\_ideal\_aggregate & mid & r0 & -0.6119963968804701 & False & not\_feasible \\
CHL & MUT & mid & r0 & 0.0547886336536989 & False & not\_feasible \\
CHL & PBI & mid & r0 & 0.0388238152491779 & False & not\_feasible \\
CHL & UBI & mid & r0 & 0.0174750960041515 & False & not\_feasible \\
COL & PEN & mid & r0 & 0.059840611829065 & False & not\_feasible \\
COL & GMI:GMI\_ideal\_aggregate & mid & r0 & -0.045557782985473 & False & not\_feasible \\
COL & MUT & mid & r0 & 0.0237371951087018 & False & not\_feasible \\
COL & PBI & mid & r0 & 0.0190932876417367 & False & not\_feasible \\
COL & UBI & mid & r0 & 0.0082461584314325 & False & not\_feasible \\
MEX & PEN & mid & r0 & 0.0650294871158911 & False & not\_feasible \\
MEX & GMI:GMI\_ideal\_aggregate & mid & r0 & -0.0851234261100203 & False & not\_feasible \\
MEX & MUT & mid & r0 & 0.0411154331581197 & False & not\_feasible \\
MEX & PBI & mid & r0 & 0.0324771874434579 & False & not\_feasible \\
MEX & UBI & mid & r0 & 0.0130541684472778 & False & not\_feasible \\
\midrule
\multicolumn{7}{l}{\textit{Panel B: all crossing cells}} \\
\midrule
PER & PEN & stress & r0 & 1.0152041145854496 & True & accounting\_feasible\_debt\_failed \\
PER & PEN & stress & r1 & 1.0926140795425037 & True & accounting\_feasible\_debt\_failed \\
PER & PEN & stress & r2 & 1.15780271593806 & True & accounting\_feasible\_debt\_failed \\
PER & PEN & stress & r3 & 1.0499800745562868 & True & accounting\_feasible\_debt\_failed \\
PER & PEN & stress & r4 & 1.4222907379520815 & True & accounting\_feasible\_debt\_failed \\
CHL & GMI:GMI\_ideal\_aggregate & stress & r0 & 3.901937871138183 & True & feasible\_cell \\
CHL & GMI:GMI\_loaded\_aggregate & stress & r0 & 2.6012919140921213 & True & feasible\_cell \\
CHL & GMI:GMI\_ideal\_aggregate & stress & r1 & 4.239349420616151 & True & feasible\_cell \\
CHL & GMI:GMI\_loaded\_aggregate & stress & r1 & 2.826232947077433 & True & feasible\_cell \\
CHL & GMI:GMI\_ideal\_aggregate & stress & r2 & 4.550059741080454 & True & feasible\_cell \\
CHL & GMI:GMI\_loaded\_aggregate & stress & r2 & 3.033373160720303 & True & feasible\_cell \\
CHL & GMI:GMI\_ideal\_aggregate & stress & r3 & 4.015460658881374 & True & feasible\_cell \\
CHL & GMI:GMI\_loaded\_aggregate & stress & r3 & 2.676973772587582 & True & feasible\_cell \\
CHL & GMI:GMI\_ideal\_aggregate & stress & r4 & 5.799668762204835 & True & feasible\_cell \\
CHL & GMI:GMI\_loaded\_aggregate & stress & r4 & 3.86644584146989 & True & feasible\_cell \\
\bottomrule
\end{longtable}
\endgroup


\subsection{Historical plausibility and debt guardrails}\label{subsec:results_plausibility}

The channel-based capture rates that generate the baseline results are themselves
historically ordinary: constructed marginal capture lies below each country's median
historical capture, so no baseline result relies on fiscally unusual behavior. The
discipline binds on the \emph{required} side. Table~\ref{tab:results_plausibility} reports
the classification of required captures against each country's historical distribution,
with bootstrap confidence intervals on the percentile boundaries. Only one requirement in
the moderate scenario lies inside historical experience: Chile's guaranteed minimum income,
whose required capture of roughly $0.17$ is classified as fiscally ordinary, with a
borderline flag reflecting proximity to the confidence band. Every other requirement is
outside historical support, in most cases by multiples of the ninetieth percentile. Peru's
universal pension at stress is the instructive marginal case: its required capture of
$0.305$ exceeds the ninetieth percentile of $0.297$ but falls inside that percentile's
bootstrap interval, and is reported as extreme with an explicit borderline flag rather than
as a clean classification. These class assignments are invariant to the historical window:
recomputing the percentiles on post-2010, post-2015, and pandemic-excluded samples changes
no classification, and only adds borderline flags where the shorter samples widen the
confidence bands.

The debt-stabilizing guardrail then separates the pilot countries more sharply than any
other layer. Chile's stabilizing primary-balance gap is approximately zero, so its
debt-consistent requirements coincide with the baseline requirements. Peru requires a fiscal
adjustment of roughly $0.7$ percent of GDP before new spending is debt-consistent, which is
why its stress crossings fail the guardrail throughout. For Colombia and Mexico the
guardrail is decisive, as the required persistence horizons of
Section~\ref{subsec:results_inversions} make explicit.

\begin{landscape}
\fontsize{7}{8.4}\selectfont
\setlength{\tabcolsep}{1.8pt}
\let\tableunderscore\_
\renewcommand{\_}{\tableunderscore\allowbreak}
\begin{longtable}{@{}
>{\raggedright\arraybackslash}p{0.055\linewidth}
>{\raggedright\arraybackslash}p{0.065\linewidth}
>{\raggedright\arraybackslash}p{0.05\linewidth}
>{\raggedright\arraybackslash}p{0.11\linewidth}
>{\raggedright\arraybackslash}p{0.19\linewidth}
>{\raggedright\arraybackslash}p{0.09\linewidth}
>{\raggedright\arraybackslash}p{0.22\linewidth}
>{\raggedright\arraybackslash}p{0.10\linewidth}
@{}}
\caption{Historical plausibility of constructed channel-based captures. Headline rows; constant provenance fields and the full 80-row grid in the Online Appendix and the reproducibility package.}\label{tab:results_plausibility}\\
\toprule
country\_id & scenario\_id & regime\_id & mfc\_gross & tax\_class & tax\_borderline & total\_revenue\_class & total\_revenue\_borderline \\
\midrule
\endfirsthead
\caption[]{Historical plausibility of constructed channel-based captures. Headline rows; constant provenance fields and the full 80-row grid in the Online Appendix and the reproducibility package. (continued)}\\
\toprule
country\_id & scenario\_id & regime\_id & mfc\_gross & tax\_class & tax\_borderline & total\_revenue\_class & total\_revenue\_borderline \\
\midrule
\endhead
\multicolumn{8}{l}{\textit{Panel A: moderate scenario, all countries and regimes}} \\
\midrule
PER & mid & r0 & 0.0847312395305157 & fiscally\_ordinary & False & fiscally\_ordinary & False \\
PER & mid & r1 & 0.0893889881373077 & fiscally\_ordinary & False & fiscally\_ordinary & False \\
PER & mid & r2 & 0.0935151115929293 & fiscally\_ordinary & False & fiscally\_ordinary & False \\
PER & mid & r3 & 0.083061682901925 & fiscally\_ordinary & False & fiscally\_ordinary & False \\
PER & mid & r4 & 0.1040846588277146 & fiscally\_ordinary & False & fiscally\_ordinary & False \\
CHL & mid & r0 & 0.0893304528213086 & fiscally\_ordinary & False & fiscally\_ordinary & False \\
CHL & mid & r1 & 0.0935828857462592 & fiscally\_ordinary & False & fiscally\_ordinary & False \\
CHL & mid & r2 & 0.097890095749623 & fiscally\_ordinary & False & fiscally\_ordinary & False \\
CHL & mid & r3 & 0.0869487663566865 & fiscally\_ordinary & False & fiscally\_ordinary & False \\
CHL & mid & r4 & 0.1088182048463007 & fiscally\_ordinary & False & fiscally\_ordinary & False \\
COL & mid & r0 & 0.0891126224653829 & fiscally\_ordinary & False & fiscally\_ordinary & False \\
COL & mid & r1 & 0.0942179579136647 & fiscally\_ordinary & False & fiscally\_ordinary & False \\
COL & mid & r2 & 0.0986287384548064 & fiscally\_ordinary & False & fiscally\_ordinary & False \\
COL & mid & r3 & 0.0874569590495911 & fiscally\_ordinary & False & fiscally\_ordinary & False \\
COL & mid & r4 & 0.1097821135538837 & fiscally\_ordinary & False & fiscally\_ordinary & False \\
MEX & mid & r0 & 0.0996944510359932 & fiscally\_ordinary & False & fiscally\_ordinary & False \\
MEX & mid & r1 & 0.1014204647253859 & fiscally\_ordinary & False & fiscally\_ordinary & False \\
MEX & mid & r2 & 0.1058634498915898 & fiscally\_ordinary & False & fiscally\_ordinary & False \\
MEX & mid & r3 & 0.094363472129331 & fiscally\_ordinary & False & fiscally\_ordinary & False \\
MEX & mid & r4 & 0.1170075658078403 & fiscally\_ordinary & False & fiscally\_ordinary & False \\
\midrule
\multicolumn{8}{l}{\textit{Panel B: low, high, and stress scenarios at r0}} \\
\midrule
PER & low & r0 & 0.0981653639417149 & fiscally\_ordinary & False & fiscally\_ordinary & False \\
CHL & low & r0 & 0.1030774989002298 & fiscally\_ordinary & False & fiscally\_ordinary & False \\
COL & low & r0 & 0.0976167873194511 & fiscally\_ordinary & False & fiscally\_ordinary & False \\
MEX & low & r0 & 0.1293057836721054 & fiscally\_ordinary & False & fiscally\_ordinary & False \\
PER & high & r0 & 0.0828730385225199 & fiscally\_ordinary & False & fiscally\_ordinary & False \\
CHL & high & r0 & 0.0872636558453767 & fiscally\_ordinary & False & fiscally\_ordinary & False \\
COL & high & r0 & 0.0859243641845317 & fiscally\_ordinary & False & fiscally\_ordinary & False \\
MEX & high & r0 & 0.0993773434305047 & fiscally\_ordinary & False & fiscally\_ordinary & False \\
PER & stress & r0 & 0.0788881354954433 & fiscally\_ordinary & False & fiscally\_ordinary & False \\
CHL & stress & r0 & 0.0829220288081224 & fiscally\_ordinary & False & fiscally\_ordinary & False \\
COL & stress & r0 & 0.0811501638454694 & fiscally\_ordinary & False & fiscally\_ordinary & False \\
MEX & stress & r0 & 0.0929324928024102 & fiscally\_ordinary & False & fiscally\_ordinary & False \\
\bottomrule
\end{longtable}
\end{landscape}


\subsection{Monte Carlo feasibility probabilities}\label{subsec:results_mc}

Table~\ref{tab:results_mc} propagates the registered parameter uncertainty through the headline
cells and reports crossing probabilities on the frozen buffer grid, under the three
pre-registered denominators (all draws, arithmetically valid draws, and support-valid
draws) and with convergence verified on a ladder of draw sizes; the full 480-cell grid is
provided in the reproducibility package.
Under the moderate scenario the probabilistic results confirm the deterministic ones with
near-certainty: crossing probabilities are indistinguishable from zero for every instrument
and country. Under the stress support, Chile's guaranteed minimum income crosses with
probability above $0.97$ in every regime while passing the debt guardrail; Peru's universal
pension crosses with probability between roughly $0.6$ and $0.8$ depending on the regime,
but its debt-consistent crossing probability is below one percent, which is what the
fragile classification records. Three robustness properties of these probabilities deserve
emphasis. First, they are insensitive to the declared institutional-dependence structure:
imposing the signed rank-dependence at intensities between $0.3$ and $0.8$, under either a
one-factor or a pairwise implementation, trims stress probabilities by two to four
percentage points and never moves any result across a decision threshold. Second, the
driver ranking is stable across methods: the scenario shock dominates everywhere, followed
by the adoption anchor and the frontier-exposure calibration, with rent-capture parameters
entering only in the reform regimes that activate them. Third, the labelled historical
reduced-form robustness mode is systematically more optimistic than the channel mode ---
for Chile's mid-scenario minimum income it yields crossing probabilities in the $0.7$ to
$0.9$ range where the channel mode yields effectively zero --- which quantifies how much of
the optimism in reduced-form reasoning the channel accounting disciplines away.

\begin{landscape}
\fontsize{7}{8.4}\selectfont
\setlength{\tabcolsep}{1.5pt}
\begin{longtable}{llllllllll}
\caption{Table 7. Monte Carlo probabilities and percentiles}\label{tab:results_mc}\\
\toprule
country\_id & policy\_variant\_id & scenario\_id & regime\_id & prob\_v\_ge\_1 & prob\_v\_ge\_1\_10 & v\_p05 & v\_p50 & v\_p95 & converged\_flag \\
\midrule
\endfirsthead
\caption[]{Table 7. Monte Carlo probabilities and percentiles (continued)}\\
\toprule
country\_id & policy\_variant\_id & scenario\_id & regime\_id & prob\_v\_ge\_1 & prob\_v\_ge\_1\_10 & v\_p05 & v\_p50 & v\_p95 & converged\_flag \\
\midrule
\endhead
PER & PEN & low & r0 & 0.0 & 0.0 & -0.280804716583315 & -0.2647799071131987 & -0.2439667941816916 & True \\
PER & GMI:GMI\_ideal\_aggregate & low & r0 & 0.0 & 0.0 & -0.1243837331284456 & -0.1211790996123289 & -0.1170169035251048 & True \\
PER & GMI:GMI\_loaded\_aggregate & low & r0 & 0.0 & 0.0 & -0.082922488752297 & -0.0807860664082192 & -0.0780112690167365 & True \\
PER & MUT & low & r0 & 0.0 & 0.0 & -0.0183544548761689 & -0.016681636180756 & -0.014508969823225 & True \\
PER & PBI & low & r0 & 0.0 & 0.0 & -0.0086262833396594 & -0.0075055059732242 & -0.0060498340076888 & True \\
PER & UBI & low & r0 & 0.0 & 0.0 & -0.0025338848741798 & -0.0021156802003266 & -0.0015725136109439 & True \\
PER & PEN & low & r1 & 0.0 & 0.0 & -0.274902912445793 & -0.2578261787682911 & -0.2355709066968619 & True \\
PER & GMI:GMI\_ideal\_aggregate & low & r1 & 0.0 & 0.0 & -0.1222973979505572 & -0.118862804773338 & -0.1143937069247924 & True \\
PER & GMI:GMI\_loaded\_aggregate & low & r1 & 0.0 & 0.0 & -0.0815315986337048 & -0.079241869848892 & -0.0762624712831949 & True \\
PER & MUT & low & r1 & 0.0 & 0.0 & -0.0178763568176213 & -0.016108408511155 & -0.0137975508757999 & True \\
PER & PBI & low & r1 & 0.0 & 0.0 & -0.0083529385026931 & -0.0071727399174334 & -0.0056270182961965 & True \\
PER & UBI & low & r1 & 0.0 & 0.0 & -0.0024409921709079 & -0.0020007741584386 & -0.001422739880091 & True \\
PER & PEN & low & r2 & 0.0 & 0.0 & -0.2729987440801695 & -0.2551085743721981 & -0.2319049592375962 & True \\
PER & GMI:GMI\_ideal\_aggregate & low & r2 & 0.0 & 0.0 & -0.1219723691750035 & -0.118327627604195 & -0.1135814504750464 & True \\
PER & GMI:GMI\_loaded\_aggregate & low & r2 & 0.0 & 0.0 & -0.0813149127833356 & -0.0788850850694633 & -0.0757209669833643 & True \\
PER & MUT & low & r2 & 0.0 & 0.0 & -0.0176925857349113 & -0.0158205519480509 & -0.013394478208276 & True \\
PER & PBI & low & r2 & 0.0 & 0.0 & -0.0082312387983086 & -0.0069802263821024 & -0.0053471631421403 & True \\
PER & UBI & low & r2 & 0.0 & 0.0 & -0.0023960690003928 & -0.0019288954793936 & -0.0013195292128714 & True \\
PER & PEN & low & r3 & 0.0 & 0.0 & -0.2797413894001765 & -0.2640377724105686 & -0.2432863705432179 & True \\
PER & GMI:GMI\_ideal\_aggregate & low & r3 & 0.0 & 0.0 & -0.1237725284562668 & -0.1206100200513087 & -0.1164575139697089 & True \\
PER & GMI:GMI\_loaded\_aggregate & low & r3 & 0.0 & 0.0 & -0.0825150189708445 & -0.0804066800342058 & -0.0776383426464726 & True \\
PER & MUT & low & r3 & 0.0 & 0.0 & -0.0183047417017689 & -0.0166693089912612 & -0.014522115305649 & True \\
PER & PBI & low & r3 & 0.0 & 0.0 & -0.0086119930167172 & -0.0075193541161494 & -0.0060775087319363 & True \\
PER & UBI & low & r3 & 0.0 & 0.0 & -0.0025321736657254 & -0.0021248232147978 & -0.0015882096910363 & True \\
PER & PEN & low & r4 & 0.0 & 0.0 & -0.2655341702231792 & -0.2463613921149541 & -0.221047802230051 & True \\
PER & GMI:GMI\_ideal\_aggregate & low & r4 & 0.0 & 0.0 & -0.119515406425721 & -0.1156334012130659 & -0.1105439063693619 & True \\
PER & GMI:GMI\_loaded\_aggregate & low & r4 & 0.0 & 0.0 & -0.0796769376171473 & -0.0770889341420439 & -0.0736959375795746 & True \\
PER & MUT & low & r4 & 0.0 & 0.0 & -0.0170671277914975 & -0.0150607808841668 & -0.0124345331000499 & True \\
PER & PBI & low & r4 & 0.0 & 0.0 & -0.0078648695844642 & -0.0065208898808358 & -0.0047566692774386 & True \\
PER & UBI & low & r4 & 0.0 & 0.0 & -0.0022698623057646 & -0.0017677902562943 & -0.0011087942671315 & True \\
PER & PEN & mid & r0 & 0.0 & 0.0 & -0.0778164699494851 & 0.0535751585119926 & 0.2357362802872228 & True \\
PER & GMI:GMI\_ideal\_aggregate & mid & r0 & 0.0 & 0.0 & -0.0837902434052994 & -0.0575146101698452 & -0.0210861186322114 & True \\
PER & GMI:GMI\_loaded\_aggregate & mid & r0 & 0.0 & 0.0 & -0.0558601622701996 & -0.0383430734465635 & -0.0140574124214743 & True \\
PER & MUT & mid & r0 & 0.0 & 0.0 & 0.0028353467593133 & 0.0165512273082204 & 0.0355668998908652 & True \\
PER & PBI & mid & r0 & 0.0 & 0.0 & 0.0055707423974533 & 0.0147602908656151 & 0.0275006646410949 & True \\
PER & UBI & mid & r0 & 0.0 & 0.0 & 0.0027635655346906 & 0.0061925356719174 & 0.0109464538175786 & True \\
PER & PEN & mid & r1 & 0.0 & 0.0 & -0.0604731707587133 & 0.0773468093548506 & 0.2683230476658478 & True \\
PER & GMI:GMI\_ideal\_aggregate & mid & r1 & 0.0 & 0.0 & -0.0793684336403056 & -0.0518594179968017 & -0.0136697127351917 & True \\
PER & GMI:GMI\_loaded\_aggregate & mid & r1 & 0.0 & 0.0 & -0.0529122890935371 & -0.0345729453312011 & -0.0091131418234611 & True \\
PER & MUT & mid & r1 & 0.0 & 0.0 & 0.0044822878511361 & 0.0188823263751837 & 0.0388178868598617 & True \\
PER & PBI & mid & r1 & 0.0 & 0.0 & 0.0066146636746602 & 0.0162690962237668 & 0.0296283125295264 & True \\
PER & UBI & mid & r1 & 0.0 & 0.0 & 0.0031420404855384 & 0.0067449850250472 & 0.0117309400023792 & True \\
PER & PEN & mid & r2 & 0.0 & 0.0 & -0.0467389985709556 & 0.0972085294438288 & 0.2999783704330296 & True \\
PER & GMI:GMI\_ideal\_aggregate & mid & r2 & 0.0 & 0.0 & -0.0765936156653832 & -0.0478658868283457 & -0.0072538642205776 & True \\
PER & GMI:GMI\_loaded\_aggregate & mid & r2 & 0.0 & 0.0 & -0.0510624104435888 & -0.0319105912188971 & -0.0048359094803851 & True \\
PER & MUT & mid & r2 & 0.0 & 0.0 & 0.0059448648639813 & 0.0209688356486348 & 0.0421598290193806 & True \\
PER & PBI & mid & r2 & 0.0 & 0.0 & 0.0076066411547585 & 0.0176838274420554 & 0.0318797011409474 & True \\
PER & UBI & mid & r2 & 0.0 & 0.0 & 0.0035138334581562 & 0.0072755071488945 & 0.0125733153453045 & True \\
PER & PEN & mid & r3 & 0.0 & 0.0 & -0.0766220110920464 & 0.0536235975904663 & 0.2353403298434292 & True \\
PER & GMI:GMI\_ideal\_aggregate & mid & r3 & 0.0 & 0.0 & -0.0831305314403983 & -0.0570945820983098 & -0.0208445384822125 & True \\
PER & GMI:GMI\_loaded\_aggregate & mid & r3 & 0.0 & 0.0 & -0.0554203542935988 & -0.0380630547322065 & -0.0138963589881417 & True \\
PER & MUT & mid & r3 & 0.0 & 0.0 & 0.0029038365451377 & 0.0164924192301619 & 0.0354653184952051 & True \\
PER & PBI & mid & r3 & 0.0 & 0.0 & 0.0056014055549311 & 0.0147010498022561 & 0.0273959017717674 & True \\
PER & UBI & mid & r3 & 0.0 & 0.0 & 0.0027693464855577 & 0.0061659335919686 & 0.0108992483308426 & True \\
PER & PEN & mid & r4 & 0.0 & 0.0 & -0.0144663083598352 & 0.1452800686702738 & 0.3700650241572971 & True \\
PER & GMI:GMI\_ideal\_aggregate & mid & r4 & 0.0 & 0.0 & -0.0691967900996816 & -0.0373134582547987 & 0.0076282321629495 & True \\
PER & GMI:GMI\_loaded\_aggregate & mid & r4 & 0.0 & 0.0 & -0.0461311933997877 & -0.0248756388365324 & 0.005085488108633 & True \\
PER & MUT & mid & r4 & 0.0 & 0.0 & 0.0091578114694959 & 0.0258283804473755 & 0.0492876139050117 & True \\
PER & PBI & mid & r4 & 0.0 & 0.0 & 0.0096944087798487 & 0.020882655570372 & 0.036592385503379 & True \\
PER & UBI & mid & r4 & 0.0 & 0.0 & 0.0042808869577225 & 0.008457142435398 & 0.0143203871991269 & True \\
PER & PEN & high & r0 & 0.0 & 0.0 & -0.1197760798261581 & 0.0064055030258902 & 0.1805557822559446 & True \\
PER & GMI:GMI\_ideal\_aggregate & high & r0 & 0.0 & 0.0 & -0.092181305550826 & -0.0669475746738123 & -0.0321210874882726 & True \\
PER & GMI:GMI\_loaded\_aggregate & high & r0 & 0.0 & 0.0 & -0.0614542037005507 & -0.0446317164492082 & -0.0214140583255151 & True \\
PER & MUT & high & r0 & 0.0 & 0.0 & -0.0015447876806515 & 0.0116272198371496 & 0.0298066461480013 & True \\
PER & PBI & high & r0 & 0.0 & 0.0 & 0.0026360815428624 & 0.0114612387083985 & 0.0236413330604348 & True \\
PER & UBI & high & r0 & 0.0 & 0.0 & 0.0016685319246994 & 0.0049615338041497 & 0.0095063903818626 & True \\
PER & PEN & high & r1 & 0.0 & 0.0 & -0.1051444236233894 & 0.0290436154580095 & 0.2123842238634628 & True \\
PER & GMI:GMI\_ideal\_aggregate & high & r1 & 0.0 & 0.0 & -0.0882808277878826 & -0.0615116696556362 & -0.0248586727431852 & True \\
PER & GMI:GMI\_loaded\_aggregate & high & r1 & 0.0 & 0.0 & -0.0588538851919217 & -0.0410077797704241 & -0.0165724484954568 & True \\
PER & MUT & high & r1 & 0.0 & 0.0 & -0.0001792310127901 & 0.013826809533239 & 0.0330351355721435 & True \\
PER & PBI & high & r1 & 0.0 & 0.0 & 0.0034930185574446 & 0.0128823685401276 & 0.0257488060678942 & True \\
PER & UBI & high & r1 & 0.0 & 0.0 & 0.0019771018064567 & 0.0054824630657117 & 0.0102843361376515 & True \\
PER & PEN & high & r2 & 0.0 & 0.0 & -0.0942168280226432 & 0.0469615432480898 & 0.2379162620671473 & True \\
PER & GMI:GMI\_ideal\_aggregate & high & r2 & 0.0 & 0.0 & -0.086145028194222 & -0.0578745192667392 & -0.0195429083683854 & True \\
PER & GMI:GMI\_loaded\_aggregate & high & r2 & 0.0 & 0.0 & -0.057430018796148 & -0.0385830128444928 & -0.0130286055789236 & True \\
PER & MUT & high & r2 & 0.0 & 0.0 & 0.0009685842479098 & 0.0157161505569207 & 0.035662650506037 & True \\
PER & PBI & high & r2 & 0.0 & 0.0 & 0.0042636148662047 & 0.0141520105211384 & 0.0275216767722818 & True \\
PER & UBI & high & r2 & 0.0 & 0.0 & 0.0022655240320006 & 0.0059563023293532 & 0.0109462456289674 & True \\
PER & PEN & high & r3 & 0.0 & 0.0 & -0.1153534323989077 & 0.0136378251875298 & 0.1909656774213741 & True \\
PER & GMI:GMI\_ideal\_aggregate & high & r3 & 0.0 & 0.0 & -0.0908294285156541 & -0.0650189941216629 & -0.0296831765125295 & True \\
\bottomrule
\end{longtable}
\end{landscape}

\subsection{Threshold inversions and required persistence horizons}\label{subsec:results_inversions}

Because the baseline does not cross, the informative objects are the inversions: what any
country would need for crossing. Table~\ref{tab:results_inversions} reports the headline cells'
required AI growth, required capture, and adoption-side requirements with their impossibility
flags; the full 960-row grid across both requirement bases is provided in the reproducibility package.
The structural finding is that the binding margin at moderate scenarios is translation, not
taxation. For every country-policy pair in the moderate support, the required translation
factor exceeds one --- no attainable level of domestic adoption would suffice, because the
translation factor is capped at the frontier by construction --- and the required capture is
therefore not the operative constraint even where it is astronomical. The AI-rent channel
cannot bridge the gap either: required effective rent-capture rates exceed any declared or
statutory bound by an order of magnitude, so reform regimes relax but do not remove the
constraint.

Table~\ref{tab:results_h_star} translates the same requirements into time, using the
required persistence horizon of Section~\ref{subsec:h_star_inversion}. Under the disruptive
support, Chile's guaranteed minimum income would require six years of scenario persistence
in the status quo and five under combined reform, with debt-consistency adding nothing
because Chile's stabilizing gap is approximately zero; under the moderate support it would
require nineteen years, within the reporting cap but beyond any comfortable policy horizon.
Peru's universal pension would require eight to eleven years of disruptive persistence on an
accounting basis, but debt-consistency doubles the requirement to seventeen to twenty-two
years. Colombia's pension and Mexico's minimum income reach the threshold only under
thirteen to twenty-one years of disruptive persistence, and in both cases the
debt-consistent horizon is censored: it exceeds the twenty-five-year cap. Every other
combination is censored on both bases. The Monte Carlo layer sharpens the same message:
the probability that Chile's stress-scenario minimum income crosses within a decade is
$0.98$, while the corresponding probability for the moderate scenario is zero. The
feasibility frontier, the historical-plausibility classes, and the persistence horizons are
three languages for a single set of facts, ordered here from most technical to most
communicable.

\begin{landscape}
\fontsize{6}{7.2}\selectfont
\setlength{\tabcolsep}{1.5pt}
\setlength{\LTleft}{0pt}
\setlength{\LTright}{0pt}
\let\tableunderscore\_
\renewcommand{\_}{\tableunderscore\allowbreak}
\begin{longtable}{@{}
>{\raggedright\arraybackslash}p{0.03\linewidth}
>{\raggedright\arraybackslash}p{0.11\linewidth}
>{\raggedright\arraybackslash}p{0.04\linewidth}
>{\raggedright\arraybackslash}p{0.03\linewidth}
>{\raggedright\arraybackslash}p{0.07\linewidth}
>{\raggedright\arraybackslash}p{0.07\linewidth}
>{\raggedright\arraybackslash}p{0.07\linewidth}
>{\raggedright\arraybackslash}p{0.065\linewidth}
>{\raggedright\arraybackslash}p{0.13\linewidth}
>{\raggedright\arraybackslash}p{0.055\linewidth}
>{\raggedright\arraybackslash}p{0.13\linewidth}
>{\raggedright\arraybackslash}p{0.16\linewidth}
@{}}
\caption{Table 5. Threshold inversion: moderate-support base accounting and cited crossing cells. Headline cells; the full grid is provided in the Online Appendix and the reproducibility package.}\label{tab:results_inversions}\\
\toprule
country\_id & policy\_variant\_id & scenario\_id & regime\_id & requirement\_basis & g\_ai\_required & mfc\_required\_gross & q\_prod\_required & historical\_tax\_class & historical\_tax\_borderline & historical\_total\_revenue\_class & flags \\
\midrule
\endfirsthead
\caption[]{Table 5. Threshold inversion: moderate-support base accounting and cited crossing cells. Headline cells; the full grid is provided in the Online Appendix and the reproducibility package. (continued)}\\
\toprule
country\_id & policy\_variant\_id & scenario\_id & regime\_id & requirement\_basis & g\_ai\_required & mfc\_required\_gross & q\_prod\_required & historical\_tax\_class & historical\_tax\_borderline & historical\_total\_revenue\_class & flags \\
\midrule
\endhead
\multicolumn{12}{l}{\textit{Panel A: all countries, moderate scenario, r0, baseline requirement}} \\
\midrule
PER & PEN & mid & r0 & baseline & 0.09307805457271823 & 0.30496566450665274 &  & extreme\_or\_outside\_historical\_support & True & fiscally\_demanding & frontier\_exceeds\_one;ai\_rent\_capture\_impossible \\
PER & GMI:GMI\_ideal\_aggregate & mid & r0 & baseline & 0.4050023484263854 & 1.3269702604077591 &  & extreme\_or\_outside\_historical\_support & False & extreme\_or\_outside\_historical\_support & frontier\_exceeds\_one;ai\_rent\_capture\_impossible \\
PER & MUT & mid & r0 & baseline & 0.7069946375039964 & 2.316433131513287 &  & extreme\_or\_outside\_historical\_support & False & extreme\_or\_outside\_historical\_support & frontier\_exceeds\_one;ai\_rent\_capture\_impossible \\
PER & PBI & mid & r0 & baseline & 1.0449706775167822 & 3.4237949914380765 &  & extreme\_or\_outside\_historical\_support & False & extreme\_or\_outside\_historical\_support & frontier\_exceeds\_one;ai\_rent\_capture\_impossible \\
PER & UBI & mid & r0 & baseline & 2.7813649797830293 & 9.113005457504299 &  & extreme\_or\_outside\_historical\_support & False & extreme\_or\_outside\_historical\_support & frontier\_exceeds\_one;ai\_rent\_capture\_impossible \\
CHL & PEN & mid & r0 & baseline & 0.40192065576351904 & 1.1615055065554287 &  & extreme\_or\_outside\_historical\_support & False & extreme\_or\_outside\_historical\_support & frontier\_exceeds\_one;ai\_rent\_capture\_impossible \\
CHL & GMI:GMI\_ideal\_aggregate & mid & r0 & baseline & 0.05896136989683802 & 0.17039172987795692 &  & fiscally\_ordinary & False & fiscally\_ordinary & frontier\_exceeds\_one;ai\_rent\_capture\_impossible \\
CHL & MUT & mid & r0 & baseline & 0.3706956372998297 & 1.0712687138756083 &  & extreme\_or\_outside\_historical\_support & False & extreme\_or\_outside\_historical\_support & frontier\_exceeds\_one;ai\_rent\_capture\_impossible \\
CHL & PBI & mid & r0 & baseline & 0.607261427850001 & 1.7549172510842361 &  & extreme\_or\_outside\_historical\_support & False & extreme\_or\_outside\_historical\_support & frontier\_exceeds\_one;ai\_rent\_capture\_impossible \\
CHL & UBI & mid & r0 & baseline & 1.4385688968552328 & 4.157302371241761 &  & extreme\_or\_outside\_historical\_support & False & extreme\_or\_outside\_historical\_support & frontier\_exceeds\_one;ai\_rent\_capture\_impossible \\
COL & PEN & mid & r0 & baseline & 0.13136722647379864 & 0.4542080169969672 &  & extreme\_or\_outside\_historical\_support & True & extreme\_or\_outside\_historical\_support & frontier\_exceeds\_one;ai\_rent\_capture\_impossible \\
COL & GMI:GMI\_ideal\_aggregate & mid & r0 & baseline & 0.40962064991749136 & 1.4162815803768314 &  & extreme\_or\_outside\_historical\_support & False & extreme\_or\_outside\_historical\_support & frontier\_exceeds\_one;ai\_rent\_capture\_impossible \\
COL & MUT & mid & r0 & baseline & 0.5989270419598275 & 2.0708168342786077 &  & extreme\_or\_outside\_historical\_support & False & extreme\_or\_outside\_historical\_support & frontier\_exceeds\_one;ai\_rent\_capture\_impossible \\
COL & PBI & mid & r0 & baseline & 0.9272671166038335 & 3.206067217557692 &  & extreme\_or\_outside\_historical\_support & False & extreme\_or\_outside\_historical\_support & frontier\_exceeds\_one;ai\_rent\_capture\_impossible \\
COL & UBI & mid & r0 & baseline & 2.351386437925959 & 8.130022988472925 &  & extreme\_or\_outside\_historical\_support & False & extreme\_or\_outside\_historical\_support & frontier\_exceeds\_one;ai\_rent\_capture\_impossible \\
MEX & PEN & mid & r0 & baseline & 0.19410435280954638 & 0.6900821949601332 &  & extreme\_or\_outside\_historical\_support & False & extreme\_or\_outside\_historical\_support & frontier\_exceeds\_one;ai\_rent\_capture\_impossible \\
MEX & GMI:GMI\_ideal\_aggregate & mid & r0 & baseline & 0.1483442186317258 & 0.5273952001657148 &  & extreme\_or\_outside\_historical\_support & False & extreme\_or\_outside\_historical\_support & frontier\_exceeds\_one;ai\_rent\_capture\_impossible \\
MEX & MUT & mid & r0 & baseline & 0.44791630397528764 & 1.5924375818041663 &  & extreme\_or\_outside\_historical\_support & False & extreme\_or\_outside\_historical\_support & frontier\_exceeds\_one;ai\_rent\_capture\_impossible \\
MEX & PBI & mid & r0 & baseline & 0.6603909875923281 & 2.3478302039767063 &  & extreme\_or\_outside\_historical\_support & False & extreme\_or\_outside\_historical\_support & frontier\_exceeds\_one;ai\_rent\_capture\_impossible \\
MEX & UBI & mid & r0 & baseline & 1.752047909784934 & 6.228902390695059 &  & extreme\_or\_outside\_historical\_support & False & extreme\_or\_outside\_historical\_support & frontier\_exceeds\_one;ai\_rent\_capture\_impossible \\
\midrule
\multicolumn{12}{l}{\textit{Panel B: cited stress crossing cells}} \\
\midrule
CHL & GMI:GMI\_ideal\_aggregate & stress & r0 & baseline & 0.06351805361681283 & 0.046621787839469436 & 0.07573616388301008 & fiscally\_ordinary & False & fiscally\_ordinary & requirement\_already\_covered \\
CHL & GMI:GMI\_ideal\_aggregate & stress & r0 & debt\_consistent & 0.06351805361681283 & 0.046621787839469436 & 0.07573616388301008 & fiscally\_ordinary & False & fiscally\_ordinary & requirement\_already\_covered \\
CHL & GMI:GMI\_loaded\_aggregate & stress & r0 & baseline & 0.07322587895792605 & 0.05374726079817179 & 0.08695075719003827 & fiscally\_ordinary & False & fiscally\_ordinary & requirement\_already\_covered \\
CHL & GMI:GMI\_loaded\_aggregate & stress & r0 & debt\_consistent & 0.07322587895792605 & 0.05374726079817179 & 0.08695075719003827 & fiscally\_ordinary & False & fiscally\_ordinary & requirement\_already\_covered \\
PER & PEN & stress & r0 & baseline & 0.09997218070252074 & 0.08391723213266877 & 0.11738267782663821 & fiscally\_ordinary & False & fiscally\_ordinary &  \\
PER & PEN & stress & r0 & debt\_consistent & 0.19006141116034964 & 0.1595386581319568 &  & fiscally\_ordinary & False & fiscally\_ordinary & frontier\_exceeds\_one \\
CHL & GMI:GMI\_ideal\_aggregate & stress & r4 & baseline & 0.04846101111926 & 0.045736627752129985 & 0.05815848320663104 & fiscally\_ordinary & False & fiscally\_ordinary & requirement\_already\_covered \\
PER & PEN & stress & r4 & baseline & 0.07645867094528785 & 0.08340649018171022 & 0.09066505763630459 & fiscally\_ordinary & False & fiscally\_ordinary & requirement\_already\_covered \\
\bottomrule
\end{longtable}
\end{landscape}

\begingroup
\let\tableunderscore\_
\renewcommand{\_}{\tableunderscore\allowbreak}
\begin{table}[!htbp]
\centering
\caption{Time to threshold $H^*$ (post-baseline extension). Both status fields are \texttt{crosses\_within\_cap} in all four data rows; the common note is ``$H^*$ censored at 25 years; values are conditional years of frozen scenario persistence, not calendar forecasts.''}\label{tab:results_h_star}
\fontsize{7}{8.4}\selectfont
\setlength{\tabcolsep}{1.5pt}
\begin{tabular}{@{}
>{\raggedright\arraybackslash}p{0.06\linewidth}
>{\raggedright\arraybackslash}p{0.055\linewidth}
>{\raggedright\arraybackslash}p{0.17\linewidth}
>{\raggedright\arraybackslash}p{0.14\linewidth}
>{\raggedright\arraybackslash}p{0.07\linewidth}
>{\raggedright\arraybackslash}p{0.055\linewidth}
>{\raggedright\arraybackslash}p{0.075\linewidth}
>{\raggedright\arraybackslash}p{0.075\linewidth}
>{\raggedright\arraybackslash}p{0.105\linewidth}
>{\raggedright\arraybackslash}p{0.105\linewidth}
@{}}
\toprule
\shortstack{country\\id} & \shortstack{policy\\id} & \shortstack{policy\\variant} & \shortstack{GMI\\version} & scenario & regime & \shortstack{$H^*$\\baseline} & \shortstack{$H^*$\\debt} & \shortstack{$\Pr(H^*\leq10)$} & \shortstack{$\Pr(H^*\leq25)$} \\
\midrule
CHL & GMI & GMI:GMI\_ideal\_aggregate & GMI\_ideal\_aggregate & mid & r0 & 19.0 & 19.0 & 0.0 & 0.797 \\
CHL & GMI & GMI:GMI\_ideal\_aggregate & GMI\_ideal\_aggregate & stress & r0 & 6.0 & 6.0 & 0.978 & 1.0 \\
PER & PEN & PEN &  & stress & r2 & 10.0 & 20.0 & 0.4666 & 0.994 \\
PER & PEN & PEN &  & stress & r4 & 8.0 & 17.0 & 0.7186 & 0.9988 \\
\bottomrule
\end{tabular}
\end{table}
\endgroup


\subsection{Diagnostics and falsification}\label{subsec:results_diagnostics}

The temporal placebo is the central credibility test and it passes: running the complete
pipeline on the 2010--2018 digitalization wave with era-consistent inputs --- an ICT
productivity shock disciplined to $0.3$--$0.6$ percent annually, 2010 connectivity anchors,
and pre-2009 historical percentiles --- produces no crossing anywhere, with a maximum
feasibility ratio of $0.75$. The framework does not fabricate feasibility from a
technological wave that, as a matter of record, did not finance universal social protection.
The sectoral negative control behaves almost as sharply, and its one exception is reported
as a declared failure rather than repaired: substituting displacement-only exposure for
productive exposure collapses feasibility by exactly the exposure ratio, confirming the
mechanism, yet Chile's stress-scenario minimum income still crosses. Channel ablation
corroborates the reading that this cell is overdetermined --- it survives switching off any
single fiscal channel --- so its conditional label carries an explicit caveat: at the
disruptive support, that crossing is driven by the scale of the shock relative to a very
small policy cost, not by the productive-exposure mechanics. The remaining diagnostics
support the baseline reading. Leave-one-source-out reruns leave headline classifications
unchanged where alternative sources exist. The labor channel carries the largest share of
fiscal capture everywhere, and Peru's fragile crossing requires the labor and consumption
channels jointly. Reform-regime crossings are not rent-dependent: removing AI-rent taxation
lowers Peru's stress-reform ratio from $1.40$ to $1.27$ without changing its status.
Finally, the adoption-side informality ablation relaxes the adoption ceiling by up to
$0.34$ without changing any feasibility outcome, because current adoption anchors sit far
below the ceiling: in this design and at these adoption levels, informality binds through
the fiscal side embedded in observed effective rates rather than through the adoption
margin. Full diagnostic tables are reported in the appendix.
\subsection{Hypothesis adjudication}\label{subsec:results_adjudication}

Section~\ref{sec:expected} recorded the qualitative pattern this framework expected before
any official run, and the pre-registration discipline requires adjudicating against that
record rather than around it. Table~\ref{tab:results_adjudication} reports the verdicts.
Two hypotheses hold cleanly. Fiscal regimes matter materially (H4): reform regimes change
feasibility ratios and required horizons throughout, and every regime-dependent result
survives its mandatory behavioral-erosion companion. The fragility discipline also holds
(H5): no official crossing relies on historically extreme capture, and the one borderline
case is flagged as such rather than absorbed. Two hypotheses receive a mixed verdict, and
the mixtures are informative. The cost gradient (H1) governs within countries --- cheaper
instruments cross first everywhere --- but the instrument ordering is not universal across
countries: Chile's guaranteed minimum income is far cheaper than its universal pension
guarantee, reversing the expected ranking and confirming that instrument hierarchies are
country-structure outcomes. The informality hypothesis (H3) is supported on the fiscal
margin, which is embedded in observed effective rates, but the adoption-side channel is
inert at current adoption levels: ablating informality from the adoption ceiling relaxes
the ceiling substantially without changing any feasibility outcome, because adoption
anchors sit far below it. Finally, the exposure hypothesis (H2) is recorded as not testable
across countries by design, because the pilot uses a common regional productive-exposure
anchor; within-design it is a maintained property, not a finding, and testing it requires
the country-specific occupational exposure measures discussed under limitations.

\begin{landscape}
\fontsize{6.3}{7.6}\selectfont
\setlength{\tabcolsep}{1pt}
\setlength{\LTleft}{0pt}
\setlength{\LTright}{0pt}
\let\tableunderscore\_
\renewcommand{\_}{\tableunderscore\allowbreak}
\begin{longtable}{@{}
>{\raggedright\arraybackslash}p{0.045\linewidth}
>{\raggedright\arraybackslash}p{0.04\linewidth}
>{\raggedright\arraybackslash}p{0.115\linewidth}
>{\raggedright\arraybackslash}p{0.275\linewidth}
>{\raggedright\arraybackslash}p{0.065\linewidth}
>{\raggedright\arraybackslash}p{0.055\linewidth}
>{\raggedright\arraybackslash}p{0.185\linewidth}
>{\raggedright\arraybackslash}p{0.17\linewidth}@{}}
\caption{Final hypothesis adjudication}\label{tab:results_adjudication}\\
\toprule
tier & \shortstack{hypothesis\\id} & statement & \shortstack{evidence\\summary} & verdict & \shortstack{preliminary\\flag} & \shortstack{evidence\\file} & notes \\
\midrule
\endfirsthead
\caption[]{Final hypothesis adjudication (continued)}\\
\toprule
tier & \shortstack{hypothesis\\id} & statement & \shortstack{evidence\\summary} & verdict & \shortstack{preliminary\\flag} & \shortstack{evidence\\file} & notes \\
\midrule
\endhead
A+B+D & H1 & Cost/feasibility gradient across policy instruments & Final evidence is mixed: high-cost UBI remains only a stress benchmark, but the lower-cost ordering is not universal. PGU/PEN-style categorical pensions and ideal GMI switch order across countries; ideal GMI cannot alone support an operational ranking. & MIXED & False & fiscal\_space\_result.csv; policy\_parameter\_official\_4c.csv; country\_policy\_classification\_final.csv & H1 is retained as a qualitative gradient claim, not a strict PEN/MUT<GMI<PBI<UBI ordering. \\
A+B+D & H2 & Preparedness/adoption matter more than total exposure across countries & The official four-country baseline uses a common LAC fallback E\_prod, so cross-country exposure variation is not identified. Preparedness, adoption, Gap and informality move T, but the total-exposure margin is not separately testable. & NOT\_TESTABLE & False & value\_assignment\_table\_baseline-official-v3.csv; informality\_adoption\_ablation\_result.csv & Requires country-specific E\_prod before becoming a cross-country test. \\
A+B+D & H3 & High informality reduces feasibility through adoption and fiscal channels & Adoption-only ablation removes I from q\_bar and logit, raising q\_bar but leaving q\_prod fixed by the frozen q\_use\_target/mu inversion; max\_abs\_delta\_v=0, gained\_crossings=0. The fiscal side is not ablated because informality is embedded in observed effective tax rates. & MIXED & False & informality\_adoption\_ablation\_result.csv; diagnostic\_result.csv & The adoption-margin mechanism is inactive in the official frozen baseline because q\_bar is not binding; the fiscal channel remains non-separable without a tax-rate counterfactual. \\
A+B+D & H4 & Fiscal regimes materially affect feasibility & Regime changes materially move V and crossings in stress cells; r2/r4 raise feasible support relative to r0, with erosion companions reported for r>=1. & HOLDS & False & fiscal\_space\_result.csv; robustness\_table.csv; diagnostic\_channel\_ablation\_result.csv & The claim is about regime sensitivity, not about robust feasibility under mid scenarios. \\
A+B+D & H5 & Crossings requiring historically extreme fiscal capture are fragile & Required-MFC classes are computed against country percentiles and stress/fragile labels are retained. No final robust claim is supported by an extreme historical-capture requirement. & HOLDS & False & threshold\_inversion\_result.csv; historical\_window\_plausibility\_changes.csv; country\_policy\_classification\_final.csv & The rule is enforced by schema/tests and the required-MFC class pipeline. \\
\bottomrule
\end{longtable}
\end{landscape}


\subsection{Limitations of the empirical application}\label{subsec:results_limitations}

The empirical results inherit every maintained assumption of the framework, and the most
consequential ones deserve explicit restatement alongside the bias ledger reported in the
appendix. First, the guaranteed-minimum-income results use the ideal-targeting aggregate,
which is a mechanical lower bound on cost; the loaded companion with the declared targeting
factor is reported throughout, and a microdata-based estimate remains the natural upgrade
before journal submission. Second, productive exposure uses a common regional anchor, which
is why cross-country exposure differences are maintained rather than tested. Third, the
frozen-trajectory convention is conservative on the translation margin: the declared
conservative variant with a zero initial adoption stock eliminates the stress crossings,
while quantifying the upside would require a declared adoption path that the baseline
deliberately does not model. Fourth, a crossing year is not a launch year: the persistence
horizons of Table~\ref{tab:results_h_star} state when accumulated gains would cover
steady-state costs, and financing any earlier launch requires transition resources outside
the framework. Fifth, the pure-layering convention sets existing-program offsets to zero,
so the net-cost specification is conservative by construction. Sixth, all headline
crossings are conditional on the disruptive support, which is deliberately beyond the
baseline macroeconomic literature and is reported as a diagnostic, not a forecast; and the
required persistence horizon was added after the pre-registered baseline as a derived
statistic, a disclosure repeated here because transparency about post-baseline additions is
part of the framework's identification discipline.

\section{Interpretation Rules and Non-Claims}\label{sec:claims}

The central contribution should be stated narrowly:
\begin{quote}
This paper does not forecast whether AI will finance UBI. It provides a calibrated threshold-accounting framework to identify the threshold combinations of AI productivity, productive adoption, fiscal capture, informality reduction, labor-share allocation, and policy cost under which universal or semi-universal social protection instruments become accounting-fiscally feasible in emerging economies.
\end{quote}

This wording is not a rhetorical disclaimer. It is the identification discipline of the paper. The object of analysis is a feasibility frontier, not a prediction of future revenue.

\subsection{Allowed claims}

The framework can claim that, under stated assumptions, a country-policy pair crosses or does not cross a baseline accounting-fiscal threshold. It can report the distance from that threshold, the buffer required for robust coverage, whether the result passes a first-order debt-stabilizing primary-balance guardrail, and whether the result violates any impossibility condition. It can rank policies by required AI growth, required adoption, required fiscal capture, or maximum financeable benefit. It can also identify countries with high exposure pressure but low fiscal capacity. The optional pressure-without-capacity diagnostic is defined as
\begin{equation}\label{eq:pwc_index_v22}
    PWC_{i,p,s,r,t}=E^{disp}_{i,t}(1-A_{i,t})I^L_{i,t}\max\{0,1-V^{gross}_{i,p,s,r,t}\},
\end{equation}
where $E^{disp}$ captures displacement exposure, $(1-A)$ captures low preparedness, $I^L$ captures labour informality, and $\max\{0,1-V^{gross}\}$ captures uncovered fiscal distance. The index is a ranking diagnostic only; it is not a feasibility threshold and is not added to the fiscal-space equation.

\subsection{Maintained economic assumptions and bias directions}

A result should be read against this ledger: headline feasibility that survives the optimistic-side assumptions and fails only conservative ones is robust; the reverse is fragile. The consolidated ledger is in Table \ref{tab:maintained_assumptions_bias} (Appendix \ref{app:calibration_ledger}).

\subsection{Claims to avoid}

The paper should not claim that AI will finance a full UBI, that a given transfer is politically optimal, that implementation capacity exists, or that the model estimates actual future revenues. It should also avoid interpreting stress-test scenarios as forecasts. Baseline accounting-fiscal feasibility is not equivalent to welfare optimality, political feasibility, administrative feasibility, or a full debt-sustainability assessment. A result that fails placebo, negative-control, or leave-one-source-out diagnostics should not be presented as robust even if its point estimate satisfies $V\geq1$. Symmetrically, a finding of infeasibility must not be read as evidence that AI gains are small. Because part of the welfare gain from AI accrues as consumer surplus through lower prices and free services, which is outside the fiscal base, the framework measures taxable-income feasibility, not welfare. Infeasibility can coexist with large welfare gains that are simply hard to tax.

\section{Conclusion}\label{sec:conclusion}

This paper developed a calibrated, falsifiable threshold framework for AI-financed social protection and then subjected it to a pre-registered empirical application to Peru, Chile, Colombia, and Mexico, with every reported number generated by a tagged reproducibility package. The framework separates exposure, preparedness, productive adoption, labor-share allocation, tax-base channels, fiscal regimes, AI-rent capture, leakage, administrative costs, and debt-consistency guardrails, so that disagreement about any one link can be located and recalibrated without discarding the method.

No instrument is robustly feasible under the primary specification. The most favourable result is conditional: Chile's ideal-targeting guaranteed minimum income crosses feasibility thresholds under stress scenarios, with $V^{gross}$ ranging from about $3.9$ to $5.8$ across fiscal regimes, and survives the debt-consistent guardrail; the loaded-targeting companion is reported alongside, and the result carries the negative-control overdetermination caveat below. Peru's universal social pension achieves fragile accounting crossings that fail debt consistency. In the headline cells for Chile's ideal GMI and Peru's universal social pension, the binding constraint is translation rather than taxation: under the moderate scenario, the required translation factor exceeds one in those cells, so even saturated adoption cannot close the gap. The required persistence horizon makes this temporal: at the ten percent prudential buffer under the baseline requirement, Chile's ideal GMI needs six years of stress-scenario persistence in the status quo ($r0$) and nineteen under the moderate scenario in the status quo; imposing the debt-consistent requirement roughly doubles Peru's pension horizon, from eight--eleven to seventeen--twenty-two years across regimes. In the $H^*$ extension, Colombia's pension and Mexico's minimum income reach the threshold only under disruptive stress persistence.

The falsification protocol was informative and disciplining. The era-consistent ICT placebo for 2010--2018 does not relabel ordinary digitalization as AI fiscal space (maximum $V^{gross}\approx0.75$). The sectoral negative control produced the expected mechanical collapse when productive exposure is replaced by automation risk; the survival of Chile's GMI under stress reveals overdetermination by the scale of the shock relative to program cost, and it is reported explicitly as a caveat, not as favourable evidence.

The central contribution remains a reusable threshold method. A reader who disputes the AI shock changes $\phi_s$; who disputes diffusion changes $T$; who disputes tax institutions changes $MFC^{gross}$; who disputes program design changes $c^{gross}$ --- and the same equations return revised frontiers. Two readings survive every specification: fiscal reform matters at least as much as technological optimism, and every crossing reported here is a requirement statement about what would have to persist, not a forecast that it will.



\appendix
\section{AI-Shock Construction Module}\label{app:ai_shock_module}

This appendix formalizes the robustness checks behind the AI-shock block. The baseline text uses the frontier-normalized multiplicative translation factor because it captures complementarity and avoids dependence on internal-sample statistics. None of the auxiliary choices below is treated as structural.

\subsection{Normalization of input variables}

The preparedness reconstruction from the AIPI's four dimensions is
\begin{equation}
    A_{i,t}=\left(D_{i,t}^{w_D}H_{i,t}^{w_H}R_{i,t}^{w_R}G_{i,t}^{w_G}\right)^{1/(w_D+w_H+w_R+w_G)},
\end{equation}
Here $D$ is digital infrastructure, $H$ human capital and labour-market policy readiness, $R$ innovation and economic integration, and $G$ regulation and ethics. Geometric aggregation makes a severe bottleneck reduce absorptive capacity even when another dimension is strong; reconstructed and published AIPI values need not coincide. These symbols are local: elsewhere $D_{s,t}$ is diffusion, $H$ an accumulation horizon, and $R$ denotes revenue objects.

All AI-shock inputs must be transformed to a common $[0,1]$ scale before aggregation. The simplest transformation is min--max normalization:
\begin{equation}
    X^{MM}_{i,t}=\frac{X_{i,t}-\min_i X_{i,t}}{\max_i X_{i,t}-\min_i X_{i,t}}.
\end{equation}
This expression is used only when $\max_iX_{i,t}>\min_iX_{i,t}$. If all observations are equal, min--max normalization is not informative; a diagnostic value such as $0.5$ may be assigned only for descriptive plots, while threshold calculations must use external-frontier normalization. Because min--max normalization can depend strongly on sample composition and outliers, percentile-rank and frontier normalizations are available as labelled robustness alternatives:
\begin{align}
    X^{pct}_{i,t}&=\frac{rank(X_{i,t})-1}{N_t-1}\quad (N_t>1),\\
    X^{frontier}_{i,t}&=\min\left\{1,\frac{X_{i,t}}{X^{frontier}_t}\right\}.
\end{align}
The frontier reference is preferred to the maximum because it is less sensitive to a single outlier when constructed from an external benchmark set. For the baseline AI-translation factor, $X^{frontier}_t$ must be external to the pilot emerging-economy sample and should represent a high-preparedness, high-adoption, high-exposure benchmark. An internal pilot-sample percentile is a relative-position diagnostic, not an absolute translation benchmark. Single-country pilots must not use min--max or percentile-rank normalizations for threshold construction; they must use external-frontier normalization. Share-valued inputs that enter the translation score directly, $E^{prod}$ and $q^{prod}$, are already measured on the natural $[0,1]$ scale and enter $S$ in natural units: the only frontier normalization in the baseline translation factor is the score-level division by $S^{frontier}_{s,t}$ in Equation \eqref{eq:T_ndc_main_v10}. The input-level normalizations of this subsection apply to index-type inputs such as the preparedness dimensions and $Gap$. Applying an input-level frontier normalization to $E^{prod}$ or $q^{prod}$ and then dividing the score by $S^{frontier}$ would discount the same distance to the frontier twice and is not admissible.

\subsection{Exposure construction from occupational and sectoral maps}

When both maps are available, productive exposure is constructed as
\begin{align}
    E^{occ,prod}_{i,t}&=\sum_o \omega^{emp}_{i,o,t}e^{prod}_o,\\
    E^{sec,prod}_{i,t}&=\sum_j \omega^{va}_{i,j,t}e^{prod}_j,
\end{align}
and combined by
\begin{equation}
    E^{prod}_{i,t}=\varpi E^{occ,prod}_{i,t}+(1-\varpi)E^{sec,prod}_{i,t},
    \qquad \varpi\in[0,1].
\end{equation}
Employment weights the occupational map and value added the sectoral map; $\varpi$ selects their convex combination.

\subsection{Technical floor}

The regularized input is
\begin{equation}
    \tilde X_{i,t}=\epsilon+(1-\epsilon)X_{i,t},
\end{equation}
with
\begin{equation}
    \epsilon\in\{0,0.01,0.05\}.
\end{equation}
For low but positive interior adoption, the floor also creates a level distortion: with $\epsilon=0.05$, a small $q^{prod}$ is inflated proportionally by the transformation. The grid therefore includes $\epsilon=0$, which is the preferred option whenever threshold inversion is numerically stable. The official implementation freezes the numerical floor at $\epsilon=10^{-9}$ for support checks and de-flooring, the numerically stabilized version of the preferred grid point $\epsilon=0$, with the logit-inversion floor frozen at $\epsilon_\mu=10^{-6}$.
The floor is a numerical regularization device, not a substantive claim that every country has positive AI capability in all dimensions. The implementation therefore separates the numerical object $\tilde q^{prod,num}=\epsilon+(1-\epsilon)q^{prod}$ from the substantive adoption object $q^{prod}$ and the substantive exposure object $E^{prod}$. The floor is never allowed to create translation or fiscal space by itself. In every translation score, multiplicative or additive, zero substantive adoption or zero substantive productive exposure is imposed as a hard gate:
\begin{equation}\label{eq:q_zero_gate_v21}
    q^{prod}_{i,s,t}\leq0\ \text{or}\ E^{prod}_{i,t}\leq0\quad\Longrightarrow\quad
    S_{i,s,t}=0,\quad T_{i,s,t}=0,\quad g^{AI}_{i,s,t}=0.
\end{equation}
Equivalently, the multiplicative score is written with the indicator $\mathbf{1}\{q^{prod}>0,E^{prod}>0\}$. The regularized values $\tilde q^{prod,num}$ and $\tilde E^{prod}$ may stabilize elasticities or threshold inversions conditional on positive substantive adoption and exposure, but they cannot convert zero adoption or zero exposure into positive AI-induced output. The gate removes crossings with substantive adoption exactly equal to zero, but floor-assisted crossings remain possible when $q^{prod}$ is positive but very small. Two complementary checks therefore remain active in the baseline implementation: the $FloorDriven$ diagnostic, under its unique operational definition in Section \ref{subsec:impossibility_score}, Equation \eqref{eq:floor_driven_operational_v25}, flags truncated cells with $q^{prod,*}=0$ whose crossing depends on $\epsilon>0$, while floor sensitivity of crossings with small positive adoption is diagnosed by re-running the threshold inversions at the preferred grid point $\epsilon=0$.


\subsection{Alternative aggregation forms}

The baseline no-double-counting multiplicative index is
\begin{equation}
    S^{NDC}_{i,s,t}=\mathbf{1}\{q^{prod}_{i,s,t}>0,\ E^{prod}_{i,t}>0\}(\tilde E^{prod}_{i,t})^{\beta}(\tilde q^{prod,num}_{i,s,t})^{\gamma}.
\end{equation}
This is the preferred form when productive adoption is simulated from preparedness and frictions, $q^{prod}=h(A,I,Gap)$. Preparedness affects the shock through adoption and therefore is not multiplied again in the baseline translation index.

A full-score robustness specification is
\begin{equation}
    S^{full}_{i,s,t}=\mathbf{1}\{q^{prod}_{i,s,t}>0,\ E^{prod}_{i,t}>0\}\tilde A_{i,t}^{\alpha}(\tilde E^{prod}_{i,t})^{\beta}(\tilde q^{prod,num}_{i,s,t})^{\gamma}.
\end{equation}
This form is valid only when $q^{prod}$ is observed independently from adoption data, or when it is explicitly labeled as a full-score robustness check rather than a no-double-counting baseline. If the adoption equation includes $E^{prod}$ directly, the direct exposure term must be removed from the baseline score and the score becomes
\begin{equation}
    S^{q-only}_{i,s,t}=\mathbf{1}\{q^{prod}_{i,s,t}>0,\ E^{prod}_{i,t}>0\}(\tilde q^{prod,num}_{i,s,t})^{\gamma}.
\end{equation}

An additive robustness specification is
\begin{equation}
    S^{add}_{i,s,t}=\mathbf{1}\{q^{prod}_{i,s,t}>0,\ E^{prod}_{i,t}>0\}\left(w_E E^{prod}_{i,t}+w_q q^{prod}_{i,s,t}\right),
    \qquad w_E+w_q=1,
\end{equation}
with preparedness added only when adoption is independently observed. The zero-adoption/zero-exposure gate applies to the additive form exactly as to the multiplicative forms: without the indicator, a robustness run could report positive translation and fiscal space with zero substantive adoption, which Equation \eqref{eq:q_zero_gate_v21} rules out by construction. The additive form is not the preferred baseline because, conditional on passing the gate, it still permits one strong component to compensate for a nearly absent component.

\subsection{Elasticity grid}

Since the structural elasticities governing the translation of AI exposure into aggregate productivity are not directly observed, the baseline uses unit elasticities for transparency. The named specifications below span the main directions; the full sensitivity grid for active elasticities is $\{0.5,1,1.5,2\}$, as declared in the master calibration matrix. The elasticity menu is:
\begin{table}[H]
\centering
\caption{Elasticity-grid robustness for AI-shock translation}
\label{tab:elasticity_grid}
\small
\begin{tabular}{lccc p{0.34\textwidth}}
\toprule
Specification & $\alpha$ & $\beta$ & $\gamma$ & Interpretation \\
\midrule
No-double-counting baseline & 0 & 1 & 1 & Preparedness affects adoption, not the score directly \\
Exposure-heavy & 0 & 1.5 & 1 & Occupational/sectoral opportunity matters more \\
Adoption-heavy & 0 & 1 & 1.5 & Realized diffusion matters more \\
Low complementarity & 0 & 0.5 & 0.5 & Lower penalty for weak components \\
High complementarity & 0 & 1.5 & 1.5 & Stronger penalty for weak components \\
Full-score robustness & 1 & 1 & 1 & Only when adoption is independently observed \\
\bottomrule
\end{tabular}
\end{table}

\subsection{No-double-counting checks}

The implementation must report which information enters each component. The recommended checks are:
\begin{align}
    S^{NDC}_{i,s,t}&=\mathbf{1}\{q^{prod}_{i,s,t}>0,\ E^{prod}_{i,t}>0\}(\tilde E^{prod}_{i,t})^{\beta}(\tilde q^{prod,num}_{i,s,t})^{\gamma},\qquad q^{prod}=h(A,I,Gap),\\
    S^{q-only}_{i,s,t}&=\mathbf{1}\{q^{prod}_{i,s,t}>0,\ E^{prod}_{i,t}>0\}(\tilde q^{prod,num}_{i,s,t})^{\gamma},\qquad q^{prod}=h(A,E^{prod},I,Gap),\\
    S^{full}_{i,s,t}&=\mathbf{1}\{q^{prod}_{i,s,t}>0,\ E^{prod}_{i,t}>0\}\tilde A_{i,t}^{\alpha}(\tilde E^{prod}_{i,t})^{\beta}(\tilde q^{prod,num}_{i,s,t})^{\gamma},\qquad q^{prod}\ \text{observed independently.}
\end{align}
The first is the baseline. The second is used if exposure enters adoption. The third is a robustness specification, not the no-double-counting baseline.

\subsection{Scenario timing as robustness only}

The baseline paper is not designed to estimate calendar dates for threshold crossing. If a dynamic diffusion path is used, it should be treated only as a robustness extension:
\begin{equation}
    g^{AI}_{i,s,t}=\phi^{Y,nom}_s D_{s,t}T_{i,s,t},
    \qquad g^{AI,max}_{s,t}=\phi^{Y,nom,max}_sD_{s,t}T^{max},\quad T^{max}\equiv1,
\end{equation}
where
\begin{equation}
    D_{s,t}=\frac{1}{1+\exp[-k_s(t-t^0_s)]},
\end{equation}
and $D_{s,t}=1$ in the static implementation.
This extension changes the annual path of the shock but does not alter the central accounting-fiscal threshold logic.

\subsection{Robustness menu}

The framework's robustness menu spans the following dimensions; the officially executed subset and its outputs are recorded in the pre-registered plans and the robustness summary of the reproducibility package:
\begin{longtable}{p{0.25\textwidth}p{0.62\textwidth}}
\caption{AI-shock robustness dimensions}\label{tab:ai_shock_robustness}\\
\toprule
Dimension & Variants \\
\midrule
\endfirsthead
\toprule
Dimension & Variants \\
\midrule
\endhead
Normalization & Min--max, percentile rank, external frontier; internal-sample percentile only as relative-index robustness \\
Functional form & Multiplicative baseline, additive robustness \\
Elasticities & No-double-counting baseline, exposure-heavy, adoption-heavy, low-complementarity, high-complementarity, full-score only when adoption is independently observed \\
Base scenario $\phi^{Y,nom}_s$ & Conservative, moderate, high, disruptive stress test \\
Adoption & Observed adoption, externally calibrated adoption, simulated logistic adoption \\
Exposure & Occupational, sectoral, task-based, productive versus labor-displacement exposure \\
Double counting & $S^{NDC}$ baseline, $S^{q\text{-}only}$ when exposure enters adoption, $S^{full}$ robustness \\
Temporal structure & Static scenario, logistic diffusion robustness only \\
\bottomrule
\end{longtable}

The methodological claim is conditional: given an externally disciplined AI-productivity scenario, the model evaluates how much of that scenario can plausibly translate into domestic fiscal space under alternative assumptions about exposure, preparedness, adoption, and fiscal capture.




\section{Productive AI Adoption Module}\label{app:adoption_module}

This appendix gives the reproducible adoption module. The baseline separates exposure from adoption: exposure determines where AI can affect tasks; adoption determines whether firms and workers actually integrate AI into production.

\subsection{Access, use, and productive adoption}

Let $q^{access}$ denote access, $q^{use}$ denote use, and $q^{prod}$ denote productive adoption. The ordering is
\begin{equation}
    0\leq q^{prod}_{i,s,t}\leq \bar q_{i,t}\leq1,
\end{equation}
with $q^{access}$ and $q^{use}$ reported when observed. The stronger inequality $q^{prod}\leq q^{use}$ is imposed only if the empirical design treats both as same-period flow variables or assumes non-decreasing use. In the baseline, $q^{prod}$ is stock-like organizational integration, so it may exceed current use after a temporary decline in use.

\subsection{Logistic diffusion equation}

The baseline logistic use equation is
\begin{equation}
    q^{use}_{i,s,t}=\bar q_{i,t}\Lambda(\ell^q_{i,s,t}),
    \qquad
    \ell^q_{i,s,t}=\mu_{i,s,t}+\nu_A A_{i,t}-\nu_I I_{i,t}-\nu_G Gap_{i,t}.
\end{equation}
Productive exposure $E^{prod}$ is deliberately excluded from this baseline adoption equation because it already enters the translation score $S^{NDC}$. A robustness version may include $\nu_EE^{prod}$ in $\ell^q$, but then the direct $E^{prod}$ term must be removed from $S$.

\subsection{Country-specific adoption ceiling}

The ceiling is
\begin{equation}
    \bar q_{i,t}=\max\{0,\min[1,1-\omega_I I_{i,t}-\omega_G Gap_{i,t}]\}.
\end{equation}
For the economic rationale behind the informality term, see Section \ref{subsec:logistic_adoption_main}. A convex ceiling variant, for example $\bar q=(1-I)^{a_I}(1-Gap)^{a_G}$, may be reported as a functional-form robustness check; the baseline linear ceiling is a convention, not an estimated form.
The ceiling uses observed or harmonized informality and digital/skills gaps. Informality and $Gap$ may affect both diffusion speed and long-run capacity only if their empirical content is partitioned or residualized; otherwise the same barrier is counted twice. If a target adoption moment falls outside the interior support, the target is adjusted:
\begin{equation}
    q^{use,target,adj}_{i,t_0}
    =\min\left\{\max\left\{q^{use,target}_{i,t_0},\epsilon_\mu\right\},\bar q_{i,t_0}-\epsilon_\mu\right\},
    \qquad \bar q_{i,t_0}>2\epsilon_\mu.
\end{equation}
If $\bar q_{i,t_0}\leq2\epsilon_\mu$, the target moment is not logit-invertible and the cell is labelled ceiling-impossible for that calibration.

\subsection{Initial adoption moment}

The initial intercept is
\begin{equation}
    \mu_{i,s,t_0}=\log\left(\frac{q^{use,target,adj}_{i,t_0}}{\bar q_{i,t_0}-q^{use,target,adj}_{i,t_0}}\right)
    -\nu_A A_{i,t_0}+\nu_I I_{i,t_0}+\nu_G Gap_{i,t_0}.
\end{equation}
The calibration registry records the source of $q^{use,target}_{i,t_0}$: an OECD--Eurostat benchmark adoption average scaled by country AIPI relative to the benchmark AIPI, with a declared $\pm30\%$ support. This is an anchored diffusion convention: $\mu_{i,s,t_0}$ is a residual initial-condition term, not evidence that $A$, $I$, and $Gap$ have been structurally identified as predictors of the initial level. If no adoption data are available, low/central/high targets are calibrated from digital access, firm digitalization, and country preparedness. If the paper wants a predictive adoption model, it must use a separate structural-intercept robustness specification with $\mu_s$ or $\mu_s+\alpha_i$ and evaluate fit without target inversion. The module contains no endogenous imitation or network-diffusion term of the Bass type: with the anchored intercept held fixed, adoption paths are driven by scenario updates to preparedness, frictions, gaps, and ceilings. Adoption trajectories are therefore scenario paths, not endogenous-diffusion forecasts; if GenAI diffuses through imitation faster than preparedness improves, the baseline understates medium-term adoption.

\subsection{Productive-adoption lag}

The productive adoption path is
\begin{equation}
    q^{prod}_{i,s,t}=\min\left\{\bar q_{i,t},(1-\lambda_q)q^{prod}_{i,s,t-1}+\lambda_q q^{use}_{i,s,t}\right\},
    \qquad \lambda_q\in\{0.15,0.30,0.50\}\quad\text{in baseline, with }\lambda_q=1\text{ only as an instant-adoption stress test}.
\end{equation}
The default initial condition is $q^{prod}_{i,s,t_0}=q^{use,target,adj}_{i,t_0}$, with $q^{prod}_{i,s,t_0}=0$ used as a conservative robustness check. The value $\lambda_q=1$ is an immediate-pass-through stress test, not the baseline.

\subsection{Coefficient restrictions and robustness}

The baseline restrictions are
\begin{equation}
    \nu_A\geq0,\qquad \nu_I\geq0,\qquad \nu_G\geq0,
\end{equation}
where the signs are written so that informality and gaps lower adoption through minus signs in $\ell^q$. If the robustness equation includes $\nu_EE^{prod}$, then $\nu_E\geq0$ and the direct exposure term is removed from $S^{NDC}$ to preserve the no-double-counting convention.

\section{Functional Income Allocation Module}\label{app:labor_share_module}

This appendix formalizes the robustness checks behind the labor-share channel in Section \ref{sec:labor_share}. The module is not a wage-forecasting model. It is an accounting device for allocating AI-induced output between labor and capital/profit channels so that the fiscal conversion block can apply different taxability assumptions to each channel.

\subsection{Raw and adjusted labor share}

The preferred empirical input is an adjusted labor share. In economies with substantial self-employment or informal work, the observed compensation-of-employees share can understate the labor component of income because mixed income combines labor and capital returns. A defensible implementation should report:
\begin{equation}
    LS^{raw}_{i,t}=\frac{\text{Compensation of employees}_{i,t}}{Y^0_{i,t}},
\end{equation}
plus an adjusted measure $LS^{adj}_{i,t}$ when mixed-income imputations are available. The main results state the version used: the country-specific PWT labor share, recorded in the calibration registry; ranking stability under raw versus adjusted measures was tested in a labelled post-baseline robustness extension using the raw UN SNA compensation-of-employees share as the alternative input, with effective rates re-derived consistently: policy rankings and country-policy classifications are unchanged across the full grid, and the only cell-level effect is that Peru's most marginal stress-scenario pension crossing, in the status quo regime, falls just below the accounting threshold ($V^{gross}$ from $1.02$ to $0.98$), quantifying the fragility already flagged (amendment registry, reproducibility package).

\subsection{Automation, augmentation, and reskilling scenarios}

The latent factor-share equation separates three mechanisms:
\begin{equation}
    \Delta \ell_{i,s,t}
    =
    -\left(\delta_s-\zeta_s RST_{i,t}\right)E^{auto}_{i,t}q^{prod}_{i,s,t}
    +\psi_s E^{aug}_{i,t}q^{prod}_{i,s,t}.
\end{equation}
The baseline restrictions are
\begin{equation}
    \delta_s,\psi_s,\zeta_s\geq0,
    \qquad
    0\leq \lambda_{LS}\leq1.
\end{equation}
The empirical implementation should treat these as scenario parameters, not as known structural constants.

\begin{longtable}{p{0.25\textwidth}p{0.11\textwidth}p{0.11\textwidth}p{0.14\textwidth}p{0.27\textwidth}}
\caption{Labor-share scenario logic}\label{tab:labor_share_scenarios}\\
\toprule
Scenario & $\delta_s$ & $\psi_s$ & $\zeta_s$ & Interpretation \\
\midrule
\endfirsthead
\toprule
Scenario & $\delta_s$ & $\psi_s$ & $\zeta_s$ & Interpretation \\
\midrule
\endhead
Automation-heavy & High & Low & Low & AI mainly substitutes tasks \\
Balanced & Medium & Medium & Medium & Substitution and complementarity coexist \\
Augmentation-heavy & Low & High & Medium/high & AI mainly complements labor \\
High-reskilling & Medium & Medium & High & Institutions absorb displacement pressure \\
\bottomrule
\end{longtable}


\subsection{Exact identities and first-order approximation}

The exact labor-income and non-labor-residual changes are
\begin{equation}
    \frac{\Delta W^{AI}}{Y^0}=LS'(1+g^{AI})-LS,
\end{equation}
\begin{equation}
    \frac{\Delta NLS^{AI}}{Y^0}=NLS'(1+g^{AI})-NLS,
    \qquad NLS=1-LS,\quad NLS'=1-LS'.
\end{equation}
Only a share $\chi^{Kbase}$ of $\Delta NLS^{AI}$ is treated as taxable capital/profit base:
\begin{equation}
    \frac{\Delta \Pi^{AI,gross}}{Y^0}=\chi^{Kbase}\frac{\Delta NLS^{AI}}{Y^0},\qquad 0\leq\chi^{Kbase}\leq1.
\end{equation}
Since $\Delta LS^{AI}=LS'-LS$, expanding the labor equation gives
\begin{equation}
    \frac{\Delta W^{AI}}{Y^0}
    =LS\,g^{AI}+\Delta LS^{AI}+\Delta LS^{AI}g^{AI}.
\end{equation}
The first-order approximation is therefore
\begin{equation}
    \frac{\Delta W^{AI}}{Y^0}\approx LS\,g^{AI}+\Delta LS^{AI},
\end{equation}
when the interaction term $\Delta LS^{AI}g^{AI}$ is small. For the non-labor residual,
\begin{equation}
    \frac{\Delta NLS^{AI}}{Y^0}
    =NLS\,g^{AI}-\Delta LS^{AI}-\Delta LS^{AI}g^{AI},
\end{equation}
with first-order approximation
\begin{equation}
    \frac{\Delta NLS^{AI}}{Y^0}\approx NLS\,g^{AI}-\Delta LS^{AI}.
\end{equation}
The exact identities are the benchmark; the approximations are only interpretive shortcuts.

The parameter $\chi^{Kbase}$ converts the broad non-labor residual into the portion that is plausibly taxable as corporate profits, business income, capital income, or AI-related rents. Because it multiplies the AI-induced increment, $\chi^{Kbase}$ is a marginal-composition parameter: the AI increment is plausibly more concentrated in corporate and formal sectors than the economy-wide non-labor residual, so the taxable share of the increment may exceed the average taxable share of $NLS$, and calibrating $\chi^{Kbase}$ from average national-accounts composition is a conservative convention. Supports should therefore be chosen for the incremental base, mirroring the marginal-informality convention of Appendix \ref{app:informality_module}, with sectorally weighted variants reported where data allow. The labor and non-labor residual components exhaust the AI-induced output increment:
\begin{equation}
\begin{split}
    &\Bigl[LS'_{i,s,t}(1+g^{AI}_{i,s,t})-LS_{i,t}\Bigr]
    +\Bigl[NLS'_{i,s,t}(1+g^{AI}_{i,s,t})-NLS_{i,t}\Bigr] \\
    &=g^{AI}_{i,s,t}.
\end{split}
\end{equation}

\subsection{Bounds and functional-form robustness}

The baseline uses a logit transformation to keep $LS'$ in $(0,1)$. Define the share-point shock, measured in labor-share points, implied by the logit model as
\begin{equation}
    \Delta LS^{logit\text{-}eq}_{i,s,t}
    =\Lambda\left(\log\frac{LS_{i,t}}{1-LS_{i,t}}+\Delta\ell_{i,s,t}\right)-LS_{i,t}.
\end{equation}
Since the first term in the difference is $LS^{target}_{i,s,t}$, direct addition of this shock is algebraically identical to the baseline partial-adjustment rule:
\begin{equation}
    LS_{i,t}+\lambda_{LS}\Delta LS^{logit\text{-}eq}_{i,s,t}
    =LS_{i,t}+\lambda_{LS}(LS^{target}_{i,s,t}-LS_{i,t})
    =(1-\lambda_{LS})LS_{i,t}+\lambda_{LS}LS^{target}_{i,s,t}
    =LS'_{i,s,t}.
\end{equation}
Because $LS_{i,t},LS^{target}_{i,s,t}\in(0,1)$ and $\lambda_{LS}\in[0,1]$, this convex combination remains in $(0,1)$, so outer truncation never activates and is not an independent robustness check. The only functional-form robustness is a linear-share specification with separately calibrated linear-scale parameters $\delta^{lin}_s$, $\psi^{lin}_s$, and $\zeta^{lin}_s$; unlike the logit baseline, the linear specification is not self-bounding, so $LS'$ must there be explicitly truncated to $(0,1)$ and any activated truncation must be reported. Coefficients calibrated for log-odds space must not be reused as direct percentage-point labor-share shocks.

\subsection{Taxable labor-income channel under informality}

The labor-share module produces an income-flow change. The fiscal conversion block then determines how much of that flow is taxable. The following base-side representation is an alternative to the baseline rate-side convention in Section \ref{subsec:gross_effective_fiscal_conversion}; it is valid only if $\tau^{L,fisc,eff}$ does not apply the labor-informality penalty again. In high-informality economies, taxable labor income may be approximated by
\begin{equation}\label{eq:labor_informality_base_side_v25}
    \frac{\Delta W^{taxable}_{i,s,t}}{Y^0_{i,t}}
    =
    (1-I^{L,marg}_{i,s,t})^{\rho_L}
    \frac{\Delta W^{AI}_{i,s,t}}{Y^0_{i,t}},
\end{equation}
where $I^{L,marg}_{i,s,t}$ is marginal labor informality as defined in Section \ref{subsec:marginal_informality}; aggregate $I^L_{i,t}$ is used only when sectoral AI-shock weights are unavailable, and $\rho_L\geq0$ controls the strength of the penalty. This alternative must follow the assignment rule in Section \ref{subsec:gross_effective_fiscal_conversion} and Table \ref{tab:no_double_counting_informality}. A more detailed base-side implementation may split formal and informal labor income:
\begin{equation}
    \Delta W^{taxable}=\Delta W^{formal}+\omega^{inf}\Delta W^{informal},
    \qquad 0\leq \omega^{inf}<1.
\end{equation}
Here $\omega^{inf}$ is the taxable fraction of the informal labor-income increment, and the split is exhaustive: $\Delta W^{formal}+\Delta W^{informal}=\Delta W^{AI}$. This is reported as a fiscal-capture adjustment, not as a claim about the welfare effect of informality.

\subsection{Labor-share robustness table}

\begin{longtable}{p{0.24\textwidth}p{0.30\textwidth}p{0.34\textwidth}}
\caption{Labor-share robustness dimensions}\label{tab:labor_share_robustness}\\
\toprule
Block & Baseline & Robustness \\
\midrule
\endfirsthead
\toprule
Block & Baseline & Robustness \\
\midrule
\endhead
Labor-share measure & Adjusted $LS$ when available & Raw $LS$; alternative mixed-income imputations \\
Exposure & $E^{auto}$ and $E^{aug}$ separated & Single labor exposure as fallback \\
Functional form & Logit-bounded $LS^{target}$ with partial adjustment & Linear-share specification with separately calibrated $\delta^{lin}_s$, $\psi^{lin}_s$, and $\zeta^{lin}_s$ \\
Adjustment speed & Partial adjustment $\lambda_{LS}$ from the declared grid & Grid endpoints: immediate adjustment ($\lambda_{LS}=1$); slower adjustment (low $\lambda_{LS}$) \\
Scenarios & Balanced automation/augmentation & Automation-heavy; augmentation-heavy; high-reskilling \\
Parameters & Scenario grid for $\delta,\psi,\zeta$ & Monte Carlo draws from bounded distributions \\
Informality & Effective labor-tax penalty & Formal/informal labor split \\
\bottomrule
\end{longtable}

The core reporting rule is to show both distributional and income-flow objects: $\Delta LS^{AI}$, $\Delta W^{AI}/Y^0$, $\Delta NLS^{AI}/Y^0$, and $\Delta \Pi^{AI,gross}/Y^0$. This avoids the mistake of interpreting a lower labor share as automatically lower labor income in levels.


\section{Fiscal Conversion Module}\label{app:fiscal_conversion_module}

This appendix formalizes the fiscal conversion block in Section \ref{sec:fiscal_conversion}. The central convention is simple: effective tax rates already contain capturability adjustments. Capturability shares $\chi^j$ are used only to construct $\tau^{j,eff}$ from statutory or benchmark rates; they are not multiplied again in the main $MFC$ equation. Throughout this appendix, $\tau^{j,eff}$ abbreviates the fiscal-capture rates of the main text, $\tau^{L,fisc,eff}$, $\tau^{K,fisc,eff}$, $\tau^{C,eff}$, and $\tau^{AI,eff}$; disposable-income rates are always written with the $disp$ superscript.

\subsection{Nested specifications and no double counting}

The framework defines three nested fiscal-conversion specifications; the channel specification is the operative official mode:
\begin{longtable}{p{0.25\textwidth}p{0.39\textwidth}p{0.24\textwidth}}
\caption{Nested fiscal conversion specifications}\label{tab:fiscal_nested_specs}\\
\toprule
Specification & Formula & Use \\
\midrule
\endfirsthead
\toprule
Specification & Formula & Use \\
\midrule
\endhead
Average conversion & $fs^{avg}=\theta g^{AI}$ & Transparent benchmark \\
Buoyancy conversion & $fs^{buoy}=b\theta g^{AI}$ & Reduced-form macro benchmark \\
Channel gross conversion & $fs^{gross}=MFC^{gross}g^{AI}$ & Preferred gross revenue object \\
Effective fiscal space & $fs^{eff}=\widetilde{MFC}^{gross}g^{AI}-ac^{net,fix}-tr^{ann}-leak^{fix}$ & Feasibility benchmark \\
\bottomrule
\end{longtable}
The average and buoyancy formulas are never summed with the channel-based formula. They are alternative benchmarks for checking whether the preferred channel specification produces plausible magnitudes. Here $\theta$ denotes one of the observed revenue anchors $\theta^{tax}$, $\theta^{rev}$, or $\theta^{nonres}$ of Section \ref{sec:fiscal_conversion}, and $b$ a historical buoyancy estimate $Buoyancy^{hist}$ from Section \ref{sec:historical_plausibility}; the anchor used must be stated.

\subsection{Statutory, average, effective, and marginal rates}

The relevant tax object for the AI shock is the effective marginal fiscal capture rate applied to the incremental tax base:
\begin{equation}
    \tau^{j,eff}_{i,s,r,t},\qquad j\in\{L,K,C,AI\}.
\end{equation}
When an observed effective marginal rate is available, it enters directly. When only a statutory or benchmark rate is available, the effective rate is constructed as
\begin{equation}\label{eq:tau_eff_decomposition_v10}
    \tau^{j,eff}_{i,s,r,t}=\tau^{j,base}_{i,r,t}\chi^j_{i,s,r,t}.
\end{equation}
After this construction, the main channel formula uses $\tau^{j,eff}$ only:
\begin{equation}
    MFC^{gross}=\sum_{j\in\{L,K,C\}}\omega^j\tau^{j,eff}+\lambda^{AI}\tau^{AI,eff}.
\end{equation}
It does not use $\tau^{j,eff}\chi^j$. Under the subset convention of Equation \eqref{eq:omega_k_net_v16}, the capital weight in this sum is $\omega^{K,net}$ rather than $\omega^K$.

\subsection{Channel-revenue function and frozen-weight convention}

\begin{equation}\label{eq:channel_revenue_function_v24}
\begin{split}
    rev^{ch}_{i,s,r,t}(g)
    &=\sum_{j\in\{L,K,C\}}\tau^{j,fisc,eff}_{i,s,r,t}\frac{\Delta B^{j,tax}_{i,s,t}(g)}{Y^0_{i,t}}
      +\lambda^{AI}_{i,s,t}(g)\tau^{AI,eff}_{i,s,r,t}g,\\
    fs^{eff,ch}_{i,p,s,r,t}(g)
    &=rev^{ch}_{i,s,r,t}(g)
      -g\left(ac^{ME}_{i,s,r,t}+tr^{ME}_{i,s,r,t}+leak^{ME}_{i,s,r,t}\right)
      -ac^{net,fix}_{i,p,s,r,t}-tr^{ann}_{i,s,r,t}-leak^{fix}_{i,s,r,t}.
\end{split}
\end{equation}
The reduced-form and frozen-weight versions are the special cases in which $rev^{ch}(g)$ is linear in $g$. Under the subset convention of Equation \eqref{eq:omega_k_net_v16}, $\Delta B^{K,tax}_{i,s,t}(g)$ in Equation \eqref{eq:channel_revenue_function_v24} is the domestic post-shifting profit base net of the AI-rent component, mirroring $\omega^{K,net}$; otherwise the rent flow would enter both the $K$ term and the AI-rent term. This distinction protects the threshold formulas below from treating an endogenous fiscal-conversion object as if it were fixed.

\subsection{Consumption channel robustness}

The baseline consumption channel is induced by income flows:
\begin{equation}
\begin{split}
    \Delta W^{AI,disp}&=(1-\tau^{L,disp})\Delta W^{AI},\\
    \Delta\Pi^{AI,disp}&=(1-\tau^{K,disp})\Delta\Pi^{AI,dom},\\
    \Delta R^{AI,gross}&=(1-\chi^{Kbase})\Delta NLS^{AI},\\
    \Delta R^{AI,disp}&=\theta^{R,dom}(1-\tau^{R,disp})\Delta R^{AI,gross}.
\end{split}
\end{equation}
\begin{equation}
    \Delta C^{AI,taxable}=(1-exempt^C)(1-m^M)\left(mpc^W\Delta W^{AI,disp}+mpc^\Pi\Delta\Pi^{AI,disp}+mpc^R\Delta R^{AI,disp}\right).
\end{equation}
The term $\Delta R^{AI,disp}$ captures the residual non-labor income not treated as a directly taxable capital/profit base. It may still finance purchases in formal markets and therefore induce consumption-tax revenue. Equivalently, one may define $\Delta M^{AI}=m^M(mpc^W\Delta W^{AI,disp}+mpc^\Pi\Delta\Pi^{AI,disp}+mpc^R\Delta R^{AI,disp})$ and subtract $\Delta M^{AI}$ once. The coefficient $m^M$ and an already-defined import increment must not both be applied to the same base. Consumption exemptions are assigned either to the taxable-consumption base through $(1-exempt^C)$ or to the construction of $\tau^{C,eff}$ through $\chi^C$, but not both. Robustness checks vary $mpc^W$, $mpc^\Pi$, $mpc^R$, import leakages, and VAT exemptions.

\subsection{Capital, foreign capture, and profit shifting}

AI-induced profit income may be domestic, foreign-owned, shifted, or deductible through payments for cloud services, licenses, software, or intellectual property. The object entering this block is the gross AI-induced profit base already separated from the national-accounts residual:
\begin{equation}
    \Delta \Pi^{AI,dom}=\chi^{dom}(1-\psi^{shift})\Delta \Pi^{AI,gross}.
\end{equation}
When the tax base rather than disposable income is required, the implementation uses $\Delta \Pi^{AI,dom}$ or a country-specific taxable-profit adjustment explicitly labelled as such.
Where withholding taxes or digital services taxes apply, foreign-remitted income can be added as a separate channel. This channel is treated as a robustness extension unless country-specific rules are observed.

\subsection{AI-specific rents and fiscal regimes}

AI-specific rents are written as
\begin{equation}
    Rent^{AI}=\lambda^{AI}g^{AI}Y^0.
\end{equation}
If $\tau^{AI}$ is statutory, then
\begin{equation}\label{eq:ai_eff_decomp_v21}
    \tau^{AI,eff}_{i,s,r,t}=\tau^{AI}_{i,r,t}\chi^{AI}_{i,s,r,t},
\end{equation}
where the AI-rent capturability term is explicitly decomposed as
\begin{equation}\label{eq:chi_ai_decomp_v21}
    \chi^{AI}_{i,s,r,t}
    =\chi^{AI,dom}_{i,s,t}(1-\psi^{shift,tech}_{i,s,t})
    \chi^{AI,compliance}_{i,r,t}\chi^{AI,enforce}_{i,r,t}.
\end{equation}
The AI-rent channel is interpreted as incremental rent capture not already included in the capital/profit taxable-base channel. Each simulation must declare one convention. If AI rents are inside the capital/profit base, Equation \eqref{eq:omega_k_net_v16} uses a unit-compatible $\lambda^{AI,Kbase}$ and $\omega^{K,net}=\max\{0,\omega^K-\lambda^{AI,Kbase}\}$ or sets $\lambda^{AI}\tau^{AI,eff}=0$. Let $\Delta\Pi^{AI,rent,gross}$ denote the AI-rent component of $\Delta\Pi^{AI,gross}$. Unit compatibility with the domestic post-shifting base used by $\omega^K$ requires
\begin{equation}
\begin{split}
    \lambda^{AI,Kbase}_{i,s,t}
    &=\chi^{dom}_{i,s,t}(1-\psi^{shift}_{i,s,t})
    \frac{\Delta\Pi^{AI,rent,gross}_{i,s,t}}{\Delta\Pi^{AI,gross}_{i,s,t}}
    \frac{\Delta\Pi^{AI,gross}_{i,s,t}/Y^0_{i,t}}{g^{AI}_{i,s,t}}\\
    &=\chi^{dom}_{i,s,t}(1-\psi^{shift}_{i,s,t})
    \frac{\Delta\Pi^{AI,rent,gross}_{i,s,t}/Y^0_{i,t}}{g^{AI}_{i,s,t}}.
\end{split}
\end{equation}
If $\lambda^{AI}$ is measured on the total AI shock, it must be rescaled to this domestic post-shifting base before it is subtracted from $\omega^K$. If AI rents are outside the measured capital/profit base, $\omega^{K,net}=\omega^K$ and the AI-rent channel is additional. This rule prevents double counting between the $K$ channel and the AI-specific rent channel.
Here $\chi^{AI,dom}$ captures the domestically reachable share of AI rents, $\psi^{shift,tech}$ captures technology-specific profit shifting through intellectual property, cloud location, licensing, royalties, or platform routing, and $\chi^{AI,compliance}$ and $\chi^{AI,enforce}$ capture legal compliance and administrative enforcement. The technology-shifting term is not assumed to equal ordinary capital-profit shifting; for emerging-market applications it may be larger. Existing digital services taxes in the region are typically turnover taxes on gross revenues rather than profit taxes; they map into this block as a higher effective rate on the rent margin with its own distortion costs. If $\tau^{AI,eff}$ is used directly, none of the components in $\chi^{AI}$ is multiplied again. Positive AI-rent taxation is not assumed in the status-quo baseline unless it exists in the country-year being modeled. The AI-rent channel must represent incremental capture outside the capital/profit channel. If AI rents are modeled as part of taxable profits, set $\lambda^{AI}\tau^{AI,eff}=0$ for the separate AI-rent term or net the corresponding component out of $\omega^K$. Technology profit shifting must be placed either inside $\chi^{AI}$ or inside residual leakage, but not both.

\subsection{Fixed costs, transition annualization, and marginal-equivalent deductions}

The implementation separates administrative flow costs, annualized one-off transition costs, fixed leakage, and marginal-equivalent deductions. All objects in the fiscal-space equation are annual shares of no-AI counterfactual GDP. Administrative costs are net of savings from programs that the new policy replaces:
\begin{equation}\label{eq:admin_net_costs_v21}
\begin{split}
    ac^{new,fix}_{i,p,s,r,t}&=\frac{AC^{new,fix}_{i,p,s,r,t}}{Y^0_{i,t}},\qquad
    ac^{exist,sav}_{i,p,t}=\frac{AC^{exist,sav}_{i,p,t}}{Y^0_{i,t}},\\
    ac^{net,fix}_{i,p,s,r,t}&=ac^{new,fix}_{i,p,s,r,t}-ac^{exist,sav}_{i,p,t}.
\end{split}
\end{equation}
The administrative-savings term is included only when the gross or net policy-cost object does not already include administrative expenditure from the replaced program. If $C^{exist}$ already includes both transfer spending and administrative spending, $ac^{exist,sav}$ is set to zero to avoid double counting. A scenario-pairing rule also applies: administrative savings presuppose that the replaced programs are actually terminated, which is the same replacement interpretation that defines $C^{exist}$ on the cost side. Pure-layering implementations, in which the new instrument is added on top of existing programs, must set $ac^{exist,sav}=0$. If $V^{gross}$ is reported with $ac^{exist,sav}>0$, the combination must be labelled a conservative hybrid: the numerator credits replacement administrative savings while the gross-cost denominator does not credit the replaced transfer budget.
Per-beneficiary administrative costs are systematically ordered: means-tested instruments exceed categorical instruments, which exceed universal instruments. Pilot implementations must not apply a uniform administrative-cost ratio across instruments, since that would compound the ideal-targeting bias in the GMI comparison.

Transition expenditure is compared with recurring annual policy costs only after annualization. Let $tr^{oneoff}_{i,s,r,t}=TR^{oneoff}_{i,s,r,t}/Y^0_{i,t}$ denote an implementation stock cost and $tr^{rec}_{i,s,r,t}$ any recurring transition-management cost. For horizon $H^{tr}$ and discount or financing rate $r^{tr}$, define
\begin{equation}\label{eq:transition_annualization_v21}
    a(H^{tr},r^{tr})=\begin{cases}
    \dfrac{r^{tr}(1+r^{tr})^{H^{tr}}}{(1+r^{tr})^{H^{tr}}-1}, & r^{tr}>0,\\[6pt]
    \dfrac{1}{H^{tr}}, & r^{tr}=0,
    \end{cases}
\end{equation}
\begin{equation}\label{eq:annualized_transition_v21}
    tr^{ann}_{i,s,r,t}=tr^{rec}_{i,s,r,t}+a(H^{tr}_{i,s,r},r^{tr}_{i,s,r})tr^{oneoff}_{i,s,r,t},
    \qquad
    leak^{fix}_{i,s,r,t}=\frac{Leak^{fix}_{i,s,r,t}}{Y^0_{i,t}}.
\end{equation}
A transition component already reported as an annual flow enters $tr^{rec}$ and is not annualized again. A one-off component enters only through $tr^{oneoff}$ and the annuity factor.

Marginal-equivalent deductions are per unit of the AI-induced fiscal-base shock:
\begin{equation}\label{eq:marginal_equiv_costs_v17}
    ac^{ME}_{i,s,r,t}\geq0,
    \qquad tr^{ME}_{i,s,r,t}\geq0,
    \qquad leak^{ME}_{i,s,r,t}\geq0.
\end{equation}
They are not computed by dividing fixed or annualized costs by $g^{AI}$. A cost component must be assigned to one margin only. If a cost is an annual administrative flow, an annualized one-off transition cost, or a fixed leakage, it enters $ac^{net,fix}$, $tr^{ann}$, or $leak^{fix}$. If it scales with the AI-induced output increment, it enters $ac^{ME}$, $tr^{ME}$, or $leak^{ME}$.

\subsection{Timing, base erosion, and behavioral response}

Lagged collection is modeled on the revenue object net of marginal-equivalent deductions but before fixed costs:
\begin{equation}
    rev^{gross}_{i,s,r,t}=\widetilde{MFC}^{gross}_{i,s,r,t}g^{AI,level}_{i,s,t}.
\end{equation}
The superscript gross here means before fixed costs only: $rev^{gross}$ is already net of marginal-equivalent deductions and therefore differs from $fs^{gross}=MFC^{gross}g^{AI}$, which is gross of them. Marginal-equivalent deductions are lagged together with the revenue they net out; fixed costs are not lagged.
\begin{equation}
    rev^{lag,gross}_{i,s,r,t}=\sum_{\ell=0}^{L}w^{lag}_\ell rev^{gross}_{i,s,r,t-\ell},
    \qquad \sum_{\ell=0}^{L}w^{lag}_\ell=1,
\end{equation}
where $rev^{gross}$ is already normalized by $Y^0$. Net lag-adjusted fiscal space is then
\begin{equation}
    fs^{eff,lag}_{i,p,s,r,t}=rev^{lag,gross}_{i,s,r,t}-ac^{net,fix}_{i,p,s,r,t}-tr^{ann}_{i,s,r,t}-leak^{fix}_{i,s,r,t}.
\end{equation}
Behavioral taxable-base erosion can be modeled as
\begin{equation}
    B^{j,taxable}=B^j(1-\epsilon^{ero,j}\tau^j),
\end{equation}
where $\epsilon^{ero,j}$ captures avoidance, relocation, deductions, or reporting responses not already embedded in the baseline capital-base chain. The erosion elasticities $\epsilon^{ero,j}$ are calibrated structural parameters in the sense of Section \ref{sec:parameters}: they carry no preferred point value, are varied over a declared grid or bounded support recorded in the calibration audit trail, and must satisfy $\epsilon^{ero,j}\tau^j<1$ so that taxable bases remain positive.
This extension is the mandatory robustness companion for reform-regime results ($r\geq1$).

\subsection{Fiscal-conversion parameter grid}

\begin{longtable}{>{\raggedright\arraybackslash}p{0.20\textwidth}>{\raggedright\arraybackslash}p{0.35\textwidth}>{\raggedright\arraybackslash}p{0.31\textwidth}}
\caption{Fiscal conversion robustness dimensions}\label{tab:fiscal_conversion_robustness}\\
\toprule
Parameter & Meaning & Robustness \\
\midrule
\endfirsthead
\toprule
Parameter & Meaning & Robustness \\
\midrule
\endhead
$\tau^{L,fisc,eff}$ & Effective labor-income capture & Country-specific; low/medium/high \\
$\tau^{K,fisc,eff}$ & Effective profit capture & Baseline: no extra ordinary profit shifting after $\Delta\Pi^{AI,dom}$; rate-side shifting only if $\psi^{shift}=0$ in the base chain \\
$\tau^{C,eff}$ & Effective consumption-tax capture & With and without VAT exemptions; do not also apply $(1-exempt^C)$ in the consumption base if exemptions enter here \\
$\tau^{AI,eff}$ & Effective AI-specific rent capture & Zero in status quo; positive scenarios \\
$\omega^L$ & Tax-base exposure weight for labor income & Derived from labor-share block \\
$\omega^K$ & Tax-base exposure weight for profits/capital & Labor-share block plus domestic-ownership/profit-shifting adjustment; $\Delta\Pi^{AI,dom}/Y^0$ divided by $g^{AI}$ \\
$\omega^C$ & Induced taxable-consumption response & Low/medium/high MPC, residual-income consumption, and import leakage \\
$mpc^R,\theta^{R,dom},\tau^{R,disp}$ & Residual non-labor income routed into consumption & Conservative/central/high residual-consumption pass-through \\
$\lambda^{AI}$ & AI-rent share of the shock, incremental to the capital/profit channel & Zero/low/medium/high; set to zero if included in $\omega^K$ \\
$MFC^{gross}$ & Drawn/calibrated gross marginal fiscal capture & Channel model, positive-capture conditional PERT, or unconditional two-component mixture \\
$MFC^{eff}$ & Derived effective reporting object after fixed costs and marginal-equivalent deductions & Computed from $fs^{eff}/g^{AI,level}$; not independently drawn \\
$ac^{net,fix},tr^{ann},leak^{fix}$; $ac^{ME},tr^{ME},leak^{ME}$ & Fixed GDP-share costs and marginal-equivalent deductions; each component assigned to one margin only & Low/medium/high \\
$\Omega^{res}$ & Residual stress-test penalty not already included elsewhere & Baseline zero; stress-test only \\
$AC^{new,fix},TR^{oneoff},Leak^{fix}$ & Nominal fixed-cost and transition-stock objects & Used only before GDP normalization to $Y^0$ \\
\bottomrule
\end{longtable}

The central reporting object remains $V^{gross}=fs^{eff}/c^{gross}$, with $V^{net}=fs^{eff}/c^{net}$ reported only as a robustness object when $c^{net}>0$; if $c^{net}=0$, the cell is labelled as already covered by existing spending.

\subsection{Historical fiscal-capture implementation details}

The historical plausibility check is implemented as a separate data object in the reproducibility package. For each country-year, the replication file stores nominal GDP, tax revenue, total revenue, non-resource revenue where available, and government level. The preferred window is 2000--latest available year, with robustness windows 2010--latest and 2015--latest. Cleaning rules follow Section \ref{sec:historical_plausibility}. The historical output reports, for each country,
\begin{equation}
    \{P10,P25,P50,P75,P90\}\left(MFC^{hist,gross,+}_i\right),
    \qquad
    \{P10,P25,P50,P75,P90\}\left(Buoyancy^{hist}_i\right),
\end{equation}
and the share of years with negative gross capture.


\section{Informality and Fiscal-Capture Frictions}\label{app:informality_module}

This appendix formalizes how informality enters the fiscal-conversion block without becoming a generic penalty applied to the entire model. The key principle is that informality affects the capturability of specific tax bases. Labor informality, firm informality, consumption informality, exemptions, and profit shifting are therefore modeled separately. As in Appendix \ref{app:fiscal_conversion_module}, $\tau^{j,eff}$ abbreviates the fiscal-capture rates $\tau^{j,fisc,eff}$ of the main text.

\subsection{Informality by tax base}

The implementation distinguishes
\begin{equation}
    I^L_{i,t},\qquad I^F_{i,t},\qquad I^C_{i,t},
\end{equation}
where $I^L$ denotes labor informality, $I^F$ firm or business informality, and $I^C$ consumption or sales informality. In the baseline convention, profit shifting is represented by $\psi^{shift}_{i,s,t}$ only in the domestic profit-base chain, while consumption exemptions are represented by $exempt^C_{i,t}$.

\subsection{Taxable-share representation}

If a statutory or benchmark tax rate is used, capturability is used to construct the effective rate:
\begin{equation}\label{eq:tau_eff_informal_v10}
    \tau^{j,eff}_{i,s,r,t}=\tau^{j,base}_{i,r,t}\chi^j_{i,s,r,t},
    \qquad j\in\{L,K,C\}.
\end{equation}
The channel formula then uses $\tau^{j,eff}$ directly and does not multiply by $\chi^j$ again. Capturable shares may be constructed as
\begin{align}
    \chi^L_{i,s,r,t}&=(1-I^{L,marg}_{i,s,t})^{\rho_L},\\
    \chi^K_{i,s,r,t}&=(1-I^{F,marg}_{i,s,t})^{\rho_K},\\
    \chi^C_{i,s,r,t}&=(1-I^{C,marg}_{i,s,t})^{\rho_C}.
\end{align}
This is the baseline when profit shifting already enters $\Delta\Pi^{AI,dom}$ and $exempt^C$ already enters $\Delta C^{AI,taxable}$. If profit shifting is instead embedded in the effective capital-tax rate, use the alternative $\chi^K=(1-I^{F,marg})^{\rho_K}(1-\psi^{K,rate})^{\rho_\psi}$ and set $\psi^{shift}=0$ in the base-side capital chain. If consumption exemptions are instead embedded in the effective consumption-tax rate, use $\chi^C=(1-I^{C,marg})^{\rho_C}(1-exempt^C)^{\rho_E}$ and omit the $(1-exempt^C)$ multiplier from $\Delta C^{AI,taxable}$. If observed effective fiscal rates are used, additional informality, profit-shifting, or exemption penalties are reported only as robustness checks.

\subsection{Marginal informality of the AI shock}\label{subsec:marginal_informality}

If sectoral information is available, the marginal informality object is
\begin{equation}
\begin{aligned}
    I^{j,marg}_{i,s,t}&=\sum_k \omega^{AI}_{i,k,s,t}I^j_{i,k,t},
    &j&\in\{L,F,C\},\\
    \omega^{AI}_{i,k,s,t}&\geq0,
    &\sum_k\omega^{AI}_{i,k,s,t}&=1.
\end{aligned}
\end{equation}
The weights $\omega^{AI}_{i,k,s,t}$ are sectoral shares of the AI shock, so $I^{j,marg}_{i,s,t}$ is a weighted average in $[0,1]$. Where data allow, the weights are channel-matched: sectoral shares of the AI-induced labor-income increment for $j=L$, of the profit increment for $j=F$, and of induced consumption for $j=C$, with output-increment shares as the declared fallback. This prevents the model from mechanically applying aggregate national informality to shocks concentrated in more formal sectors such as finance, telecommunications, professional services, software, or formal manufacturing.

\subsection{No-double-counting rule}

\begin{longtable}{p{0.30\textwidth}p{0.23\textwidth}p{0.33\textwidth}}
\caption{No-double-counting rule for informality adjustments}\label{tab:no_double_counting_informality}\\
\toprule
Fiscal-rate source & Apply extra informality adjustment? & Reason \\
\midrule
\endfirsthead
\toprule
Fiscal-rate source & Apply extra informality adjustment? & Reason \\
\midrule
\endhead
Statutory or normative rate & Yes, inside $\tau^{eff}=\tau^{base}\chi$ & It does not embed compliance or informality \\
Observed effective rate from revenue/base & Not automatically & It already embeds average compliance, evasion, exemptions, and informality \\
Effective marginal rate for formal sector & Yes, with $I^{marg}$ if needed & The AI shock may occur partly outside the formal base \\
Aggregate tax-to-GDP ratio & No in the baseline & It already reflects the existing fiscal structure and informality \\
Labor channel, rate-side baseline & No base-side labor-informality penalty & Construct $\omega^L$ from gross $\Delta W^{AI}$ and include $\chi^L=(1-I^{L,marg})^{\rho_L}$ once in $\tau^{L,eff}$ \\
Labor channel, base-side alternative & No rate-side labor-informality penalty & Construct $\omega^L$ from $\Delta W^{taxable}$ and omit $\chi^L$ from $\tau^{L,eff}$ \\
Residual penalty $\Omega^{res}$ & Stress test only & Used only for frictions not included in channel-specific rates \\
\bottomrule
\end{longtable}
The official implementation derives effective fiscal rates from observed revenue composition, so no additional informality penalty is applied to them, consistent with the observed-effective-rate row of Table \ref{tab:no_double_counting_informality}.
Legal consumption exemptions and consumption informality overlap empirically: informal purchases concentrate in exempt staples. To avoid discounting the same untaxed consumption twice, $exempt^C$ must be measured on the formal consumption basket, or $I^C$ on the taxable basket, and the calibration registry states which convention is used. A parallel overlap applies to the capital channel: mixed income and informal-business income already excluded from the taxable base through $\chi^{Kbase}$ must not be penalized again through $I^{F,marg}$. Under the baseline, $\chi^{Kbase}$ is a composition filter selecting the profit-type components of the non-labor residual, while $I^{F,marg}$ is a compliance filter on firms inside that taxable-profit perimeter; the calibration registry states that $I^{F,marg}$ is measured on the perimeter-consistent firm population, or assigns the informal-business margin to one side only.

\subsection{Informality in adoption versus fiscal capture}

Informality can constrain both adoption and fiscal capture, but through different equations. In adoption, informality affects $q^{prod}$ and therefore $g^{AI}$. In fiscal conversion, informality affects $\tau^{j,eff}$ and therefore $fs^{eff}$ conditional on $g^{AI}$. The calibration registry declares whether an informality variable is being used on the adoption side, the fiscal side, or both; if both, the interpretation must be that they represent separate mechanisms rather than a repeated penalty on the same object.

\subsection{Formalization-enhancing AI as a scenario}

The baseline does not assume that AI automatically formalizes the economy. A reform scenario may allow digitalization, e-invoicing, payments traceability, or AI-assisted tax administration to reduce informality:
\begin{equation}
    I^{j,post}_{i,s,r,t}=I^j_{i,t}\left(1-\varphi^j_{r,t}q^{prod}_{i,s,t}\right).
\end{equation}
For each labor, firm, and consumption informality channel, the object is truncated to keep $I^{j,post}\in[0,1]$. This extension is inactive in the status quo and activated only in improved-compliance or combined-reform scenarios.
The endogenous formalization term applies only to fiscal capturability within the period: $I^{j,post}$ enters $\chi^j$ and $\tau^{j,eff}$, while the adoption-side informality in $\ell^q$ and $\bar q$ is held at observed values to avoid a simultaneous loop between adoption and the adoption ceiling. A dynamic extension may feed $I^{j,post}$ back into adoption with a one-period lag; it must then be labelled as a dynamic-feedback robustness scenario.

The symmetric adverse scenario is displacement-driven informalization: when automation pressure is strong and reinstatement weak, displaced formal workers may move into informality. A labelled robustness scenario allows $I^{L,post}_{i,s,r,t}=I^L_{i,t}\left(1+\varphi^{disp}_{r,t}(\delta_s-\zeta_s RST_{i,t})E^{disp}_{i,t}q^{prod}_{i,s,t}\right)$, truncated to keep $I^{L,post}\in[0,1]$, with the same fiscal-side-only timing convention as the formalization scenario. The baseline holds informality fixed; of the two endogenous deviations, only reporting the favorable one would bias automation-heavy scenarios optimistically, so implementations that activate formalization scenarios must also report this adverse counterpart.

\subsection{Elasticity and robustness grid}

The baseline elasticities $\rho_L,\rho_K,\rho_C,\rho_E,\rho_\psi$ are varied over the grid $\{0.5,1,1.5,2\}$. The profit-shifting elasticity $\rho_\psi$ is active only in the alternative rate-side convention for capital taxation. The exemption elasticity $\rho_E$ is active only in the alternative convention where consumption exemptions are embedded in $\chi^C$ rather than in $\Delta C^{AI,taxable}$. They are not treated as known structural constants.


\section{Threshold Inversion Module}\label{app:threshold_module}

This appendix formalizes additional threshold inversions used as stress tests. All thresholds are algebraic inversions of the same feasibility condition, not independent assumptions and not forecasts.

\subsection{Notation and no-double-counting rule}

Let
\begin{equation}
    R^{req}_{i,p,t}(\xi)=(1+\xi)c^{gross}_{i,p,t},
    \qquad
    R^{req,DC}_{i,p,t}(\xi)=(1+\xi)c^{gross}_{i,p,t}+s^{PB}_{i,t}.
\end{equation}
If the threshold uses gross fiscal conversion, the required gross revenue includes only fixed-cost subtractions. Marginal-equivalent deductions enter the denominator through $\widetilde{MFC}^{gross}$ and must not be added again to the required-revenue term:
\begin{equation}
    R^{req,gross}_{i,p,s,r,t}(\xi)=R^{req}_{i,p,t}(\xi)+ac^{net,fix}_{i,p,s,r,t}+tr^{ann}_{i,s,r,t}+leak^{fix}_{i,s,r,t},
\end{equation}
\begin{equation}
    R^{req,gross,DC}_{i,p,s,r,t}(\xi)=R^{req,DC}_{i,p,t}(\xi)+ac^{net,fix}_{i,p,s,r,t}+tr^{ann}_{i,s,r,t}+leak^{fix}_{i,s,r,t}.
\end{equation}
The implementation must not add $ac^{net,fix}$, $tr^{ann}$, or $leak^{fix}$ to a threshold that already uses $MFC^{eff}$. It must also not add $ac^{ME}$, $tr^{ME}$, or $leak^{ME}$ to $R^{req,gross}$, because those terms are already embedded in $\widetilde{MFC}^{gross}=MFC^{gross}-ac^{ME}-tr^{ME}-leak^{ME}$. If $\Omega^{res}$ is activated, it enters only as a stress-test penalty: either use $fs^{eff,res}=fs^{eff}-\Omega^{res}$ or add $\Omega^{res}$ to the required-revenue object, but not both.

The closed-form threshold formulas in this appendix are valid for the historical reduced-form mode or for channel-based calculations in which $\omega^j$, $\lambda^{AI}$, and $\widetilde{MFC}^{gross}$ are frozen at the evaluated scenario. If the channel bases $\Delta B^{j,tax}$ are recomputed as $g^{AI}$, $T$, or $q^{prod}$ changes, the implementation solves the threshold directly from $fs^{eff,ch}(\cdot)$ in Equation \eqref{eq:channel_revenue_function_v24}.

\subsection{Minimum AI-rent capture threshold}

Suppose the gross fiscal conversion block is written with the AI-rent term separated from the labor, capital, and consumption channels, $MFC^{base,gross}_{i,s,r,t}=MFC^{gross}_{i,s,r,t}-\lambda^{AI}_{i,s,t}\tau^{AI,eff}_{i,s,r,t}$ in the notation of Equation \eqref{eq:mfc_gross_channel_v10}, so that
\begin{equation}
    fs^{eff}_{i,p,s,r,t}=g^{AI}_{i,s,t}\left[\widetilde{MFC}^{base,gross}_{i,s,r,t}+\tau^{AI,eff}_{i,s,r,t}\lambda^{AI}_{i,s,t}\right]-ac^{net,fix}_{i,p,s,r,t}-tr^{ann}_{i,s,r,t}-leak^{fix}_{i,s,r,t},
\end{equation}
where
\begin{equation}
    \widetilde{MFC}^{base,gross}_{i,s,r,t}=MFC^{base,gross}_{i,s,r,t}-ac^{ME}_{i,s,r,t}-tr^{ME}_{i,s,r,t}-leak^{ME}_{i,s,r,t}.
\end{equation}
Solving for the minimum effective AI-rent tax gives
\begin{equation}\label{eq:tau_ai_eff_star_v10}
    \tau^{AI,eff,*}_{i,p,s,r,t}(\xi)=
    \max\left\{0,
    \frac{\frac{R^{req,gross}_{i,p,s,r,t}(\xi)}{g^{AI}_{i,s,t}}-\widetilde{MFC}^{base,gross}_{i,s,r,t}}
    {\lambda^{AI}_{i,s,t}}
    \right\},
\end{equation}
provided $g^{AI}_{i,s,t}>0$ and $\lambda^{AI}_{i,s,t}>0$. If the policy variable is a statutory rate and $\tau^{AI,eff}=\tau^{AI}\chi^{AI}$, then for $\chi^{AI}_{i,s,r,t}>0$,
\begin{equation}\label{eq:tau_ai_stat_star_v10}
    \tau^{AI,*}_{i,p,s,r,t}(\xi)=\frac{\tau^{AI,eff,*}_{i,p,s,r,t}(\xi)}{\chi^{AI}_{i,s,r,t}}.
\end{equation}
If $\chi^{AI}_{i,s,r,t}=0$ and $\tau^{AI,eff,*}>0$, the cell is AI-rent-capture-impossible. If $\tau^{AI,eff,*}>\bar\tau^{AI,eff}$ or $\tau^{AI,*}>\bar\tau^{AI}$, the policy-scenario pair is infeasible through the AI-rent-tax channel.

\subsection{Minimum productive-adoption threshold}

To avoid circularity and relative-sample inflation, the adoption threshold uses a fixed external reference score $S^{ref}_{s,t}=S^{frontier}_{s,t}$ that does not depend on the country-policy cell being inverted or on the internal pilot sample. Under the baseline no-double-counting score,
\begin{equation}
    S^{NDC}_{i,s,t}=\mathbf{1}\{q^{prod}_{i,s,t}>0,\ E^{prod}_{i,t}>0\}(\tilde E^{prod}_{i,t})^\beta(\tilde q^{prod,num}_{i,s,t})^\gamma,
\end{equation}
The closed-form translation and adoption thresholds, including the conversion back from the floor-regularized scale and lower-support truncation, are stated in Proposition \ref{prop:q_threshold_v10}. If $\phi^{Y,nom}_s\leq0$ or $\widetilde{MFC}^{gross}\leq0$ with positive required revenue, no finite adoption threshold exists and the cell is labelled non-positive technology/fiscal support. If $T^{req,static}>1$, even the saturated value $T=1$ cannot cross the threshold and the cell is classified as technology/fiscal-support insufficient. Diagnostics apply to the truncated original-scale object $q^{prod,*}$; if $q^{prod,*}>\bar q_i$, feasibility cannot be obtained by adoption alone under the maintained exposure, AI-productivity, and fiscal-conversion assumptions. The unique operational definition of floor-driven is the one in Section \ref{subsec:impossibility_score}, Equation \eqref{eq:floor_driven_operational_v25}: $q^{prod,*}=0$, $\epsilon>0$, and the crossing fails when $\epsilon=0$. Such a cell is not counted as adoption-feasible. A substantive adoption crossing must satisfy $q^{prod,*}>0$ and $q^{prod,*}\leq\bar q_i$.

\subsection{Minimum base-shock threshold}

Under the one-period static approximation and the frozen-weight convention, the required nominal GDP-equivalent base shock is
\begin{equation}\label{eq:phi_star_static_app_v25}
    \phi^{*,static}_{i,p,s,r,t}(\xi)
    =\frac{R^{req,gross}_{i,p,s,r,t}(\xi)}
    {T_{i,s,t}\widetilde{MFC}^{gross}_{i,s,r,t}}.
\end{equation}
This threshold is valid when $T_{i,s,t}>0$ and $\widetilde{MFC}^{gross}_{i,s,r,t}>0$, under the same existence and frozen-weight rules as $T^*$. For accumulated horizons, $\phi^*$ is solved numerically from Equation \eqref{eq:accumulated_threshold_condition_v22}. If $\phi^{*,static}_{i,p,s,r,t}(\xi)>\phi^{Y,nom,max}_s$, the cell is technology-scenario-insufficient.

\subsection{Existence conditions for adoption thresholds}\label{subsec:adoption_threshold_existence}

An interior adoption threshold exists only if $R>0$, $\phi^{Y,nom}_s>0$, $\widetilde{MFC}^{gross}>0$, $\gamma>0$, $E^{prod}>0$, $\tilde E^{prod}>0$, and $S^{ref}=S^{frontier}>0$. This list is identical to the condition $Valid^{T*}_{i,p,s,r,t}(\xi)=1$ in Proposition \ref{prop:q_threshold_v10}. A growth threshold similarly requires $\widetilde{MFC}^{gross}>0$, and a gross-capture threshold requires $g^{AI}>0$. If $\widetilde{MFC}^{gross}\leq0$ while required revenue is positive, the draw is kept as an economic failure draw but no finite positive growth or adoption threshold is reported. In simulations, any use of $\epsilon>0$ is numerical regularization conditional on positive adoption; substantive feasibility at zero adoption or zero productive exposure is ruled out by construction through the gate $q^{prod}\leq0\ \text{or}\ E^{prod}\leq0\Rightarrow S=T=g^{AI}=0$.

\subsection{Preparedness and exposure thresholds}

Preparedness and exposure thresholds are reported as stress tests only. If $q^{prod}$ is simulated from preparedness, a preparedness threshold can be computed from the adoption equation. If preparedness is also in $S^{full}$, then the specification must be labeled full-score robustness rather than no-double-counting baseline.

\subsection{Impossibility diagnostics}

The threshold module reports realized coverage, support headroom, and absolute gaps separately. Realized coverage is
\begin{align}\label{eq:distance_ratios_main_v19}
    D^{g,real}&=\frac{g^{AI}}{g^{AI,*}},
    &D^{MFC,gross,real}&=\frac{MFC^{gross}}{MFC^{gross,*}},\\
    D^{q,real}&=\frac{q^{prod}}{q^{prod,*}}\quad\text{if }q^{prod,*}>0,
    &D^{AI,eff,real}&=\frac{\tau^{AI,eff}}{\tau^{AI,eff,*}}\quad\text{if }\tau^{AI,eff,*}>0.
\end{align}
Support headroom is
\begin{align}\label{eq:headroom_ratios_main_v19}
    H^g&=\frac{g^{AI,max}}{g^{AI,*}},
    &H^q&=\frac{\bar q_i}{q^{prod,*}}\quad\text{if }q^{prod,*}>0,\\
    H^{AI,eff}&=\frac{\bar\tau^{AI,eff}}{\tau^{AI,eff,*}}\quad\text{if }\tau^{AI,eff,*}>0.
\end{align}
Values above one mean that the realized value or support ceiling covers the requirement within its family. The effective AI-rent bound combines the statutory cap and capturability, $\bar\tau^{AI,eff}_{i,s,r,t}=\bar\tau^{AI}_{i,r,t}\chi^{AI}_{i,s,r,t}$. Realized and headroom families are never pooled into one cross-margin ranking. A derived effective-conversion distance must be separately labelled $D^{MFC,eff,real}=MFC^{eff}/MFC^{eff,*}$ and used only for derived reporting; gross realized capture remains the baseline plausibility metric because historical capture is gross.
If $q^{prod,*}=0$, neither $D^{q,real}$ nor $H^q$ is reported. Adoption-margin already satisfied requires feasibility without reliance on the numerical floor and without violating the zero-adoption/zero-exposure gate; otherwise the cell is floor-driven or non-substantive and not adoption-feasible. If $0<q^{prod,*}\leq\bar q_i$, support is sufficient; if $q^{prod,*}>\bar q_i$, the ceiling fails. The absolute excess $Gap^q=q^{prod,*}-\bar q_i$ is reported as a gap, never a distance ratio.


For accumulated-horizon implementations, the scenario-supported maximum AI level gain is
\begin{equation}
    g^{AI,max}_{s,H}=\prod_{\tau=1}^{H}\left(1+\phi^{Y,nom,max}_{s,\tau}D_{s,\tau}T^{max}\right)-1.
\end{equation}
The static object $g^{AI,max}_{s,t}=\phi^{Y,nom,max}_sD_{s,t}T^{max}$ is used only for one-period static runs. The consolidated classification of impossibility diagnostics is in Table \ref{tab:impossibility_diagnostics_v10}; this module reports the ratios and gaps defined above.

\subsection{Comparative statics}

Under the frozen-weight convention, each required threshold is increasing in policy cost and decreasing in the realized values of the other margins: fiscal capture for $g^{AI,*}$; exposure and adoption for thresholds that operate through $T$; and the productivity scenario for $MFC^{gross,*}$. A threshold does not depend on its own realized value. These comparative statics are algebraic and should be reported as threshold movements rather than causal estimates.

\subsection{Reporting protocol}

Every threshold table should state whether it uses gross or effective fiscal conversion, gross or net policy cost, debt-consistent or non-debt requirements, and whether the adoption score is no-double-counting or full-score robustness.


\section{Uncertainty Propagation and Robustness Module}\label{app:uncertainty_module}

This appendix formalizes the robustness and Monte Carlo protocol. The Monte Carlo exercise propagates parameter uncertainty through the fiscal-feasibility framework; it is not a forecast of AI development.


\subsection{Monte Carlo inputs}

The uncertain-input vector is mode-specific. In the channel-based fiscal mode, draw channel rates and construct $MFC^{gross}$ from them. In the historical reduced-form mode, draw $MFC^{gross}$ directly from the historical gross-capture distribution and do not also use channel rates to construct it. The positive-capture PERT version is a conditional-feasibility benchmark; the two-component mixture is the unconditional fiscal-risk benchmark.

The implementation must also declare its adoption mode. In observed-adoption mode, $q^{prod}$ is a primitive drawn from observed adoption or measurement-error distributions. In the baseline anchored simulated-adoption mode, $A,I,Gap,\lambda_q,\nu_A,\nu_I,\nu_G$, and $q^{use,target}_{i,t_0}$ are primitives, while $\mu_{i,s,t_0}$ is derived by the target inversion. In structural-intercept robustness mode, $\mu_s$ or $\mu_s+\alpha_i$ is drawn directly and the target inversion is not used. In all simulated modes, $q^{use}$ and $q^{prod}$ are derived endogenously.

Channel-based fiscal mode:
\begin{equation}
\begin{split}
    \Theta^{(m)}_{channel}=\{&\phi^{raw}_{s,u},\kappa^{u\to Y},\pi^{AI,Y},A,E^{prod},I,Gap,\lambda_q,\nu_A,\nu_I,\nu_G,\omega_I,\omega_G,q^{use,target},\\
    &\tau^{L,fisc,eff},\tau^{K,fisc,eff},\tau^{C,eff},\lambda^{AI},\tau^{AI},\chi^{AI,dom},\psi^{shift,tech},\chi^{AI,compliance},\chi^{AI,enforce},\\
    &ac^{net,fix},tr^{ann},leak^{fix},ac^{ME},tr^{ME},leak^{ME},\\
    &\rho_L,\rho_K,\rho_C,\rho_E,\rho_\psi,\alpha,\beta,\gamma,
    LS,E^{auto},E^{aug},RST,\delta,\psi,\zeta,\lambda_{LS},\chi^{Kbase},\\
    &mpc^W,mpc^{\Pi},mpc^R,m^M,exempt^C,\chi^{dom},\psi^{shift},\theta^{R,dom},\tau^{L,disp},\tau^{K,disp},\tau^{R,disp},\Omega^{res},s^{PB},c^{gross},c^{net}\}^{(m)}.
\end{split}
\end{equation}
Historical reduced-form fiscal mode:
\begin{equation}
\begin{split}
    \Theta^{(m)}_{hist}=\{&\phi^{raw}_{s,u},\kappa^{u\to Y},\pi^{AI,Y},A,E^{prod},I,Gap,\lambda_q,\nu_A,\nu_I,\nu_G,q^{use,target},\omega_I,\omega_G,
    MFC^{gross},ac^{net,fix},tr^{ann},leak^{fix},\\
    &ac^{ME},tr^{ME},leak^{ME},
    \alpha,\beta,\gamma,\Omega^{res},s^{PB},c^{gross},c^{net}\}^{(m)}.
\end{split}
\end{equation}
Nominal $AC$, $TR$, and $Leak$ objects are reporting objects derived from $Y^0$ when needed. They are not drawn independently from normalized fixed-cost or marginal-equivalent objects in the same simulation. Observed variables may be fixed in the baseline, but the replication table must label them as fixed and give source-year values. If they are stochastic, the table must give distribution, support, source, and scenario cell.

\begin{footnotesize}
\begin{longtable}{>{\raggedright\arraybackslash}p{0.18\textwidth}>{\raggedright\arraybackslash}p{0.27\textwidth}>{\raggedright\arraybackslash}p{0.28\textwidth}>{\raggedright\arraybackslash}p{0.17\textwidth}}
\caption{Monte Carlo input protocol}\label{tab:mc_inputs}\\
\toprule
Parameter block & Meaning & Baseline distribution & Robustness \\
\midrule
\endfirsthead
\toprule
Parameter block & Meaning & Baseline distribution & Robustness \\
\midrule
\endhead
AI shock block & Raw AI shock, output bridge, nominal bridge, and accumulated level object & Bounded PERT on raw support; bridge fixed or PERT & Triangular or truncated-uniform shock supports \\
Adoption block & $A$, $E^{prod}$, $I$, $Gap$, $q^{prod}$, $\lambda_q$, $\omega_I$, $\omega_G$, and anchored target moment & Observed inputs fixed or beta/PERT measurement error; simulated adoption derives $q^{use}$ and $q^{prod}$ from the anchored diffusion module & Slow/fast diffusion and structural-intercept robustness \\
Fiscal-rate block & Labor, capital, consumption, and AI-rent fiscal rates & Observed effective rates where available; otherwise PERT or fixed scenario values & Low/high fiscal capture \\
AI-rent capture block & $\lambda^{AI}$, statutory $\tau^{AI}$, and the components of $\chi^{AI}$: domestic reach, technology profit shifting, compliance, and enforcement & Zero AI-rent tax in status quo; positive scenarios use bounded shares & No AI-rent tax; high technology-shifting stress \\
Cost and leakage block & Net administrative cost, annualized transition cost, fixed leakage, and marginal-equivalent deductions & Bounded PERT or fixed source values; each cost assigned to exactly one margin & High-cost/high-leakage stress \\
Tax-base elasticities & Informality elasticities; $\rho_E$ only under the labelled exemptions-in-rate convention; $\rho_\psi$ only under labelled rate-side profit shifting & Bounded PERT; optional terms active only in their labelled conventions & Discrete grid \\
Behavioral erosion & $\epsilon^{ero,j}$ taxable-base erosion under reform regimes & Fixed at zero in the status-quo baseline & Mandatory declared grid or bounded PERT companion for $r\geq1$ \\
Residual penalty & $\Omega^{res}$ for residual stress-test frictions not counted elsewhere & Fixed at zero in baseline & Positive stress-test draws only \\
Translation elasticities & $\alpha$, $\beta$, and $\gamma$ & Discrete grid & Factorial grid \\
Labor-share block & Automation, augmentation, reskilling, and labor-share damping & Bounded PERT & Truncated-uniform robustness \\
Policy-cost block & Gross cost, net cost, and debt-consistent guardrail & Fixed baseline or scenario objects & Net-cost and debt-guardrail robustness \\
\bottomrule
\end{longtable}
\end{footnotesize}

\subsection{Scenario-consistent dependence}

The baseline does not require all parameters to be independent. Correlated latent shocks can encode $Corr(A,q)>0$, $Corr(I,q)<0$, $Corr(I,MFC)<0$, $Corr(A,MFC)>0$, and additionally $Corr(I,\chi^{Kbase})<0$ (informality raises the mixed-income share of the non-labor residual) and $Corr(A,E^{prod})>0$ (more prepared economies are also more exposed). Because the signs cluster advantages, scenario-consistent dependence widens the cross-country dispersion of $V$ relative to independence: strong countries cross earlier and weak countries later. These are robustness assumptions, not estimated causal correlations.

\subsection{Monte Carlo outputs}

For each country-policy-scenario-regime cell, report
\begin{equation}
    \{P5,P25,P50,P75,P95\}(V^{gross})
\end{equation}
and, when relevant,
\begin{equation}
    \{P5,P25,P50,P75,P95\}(V^{net})\quad\text{only for cells with }c^{net}>0.
\end{equation}
The crossing-probability grid is
\begin{equation}
    \widehat{\pi}(\xi)=\frac{1}{M}\sum_{m=1}^{M}\mathbf{1}\{V^{gross,(m)}\geq1+\xi\},
    \qquad \xi\in\{0,0.05,0.10,0.25,0.50\}.
\end{equation}
The basic-valid and support-valid versions are
\begin{equation}
    \widehat{\pi}^{basic}(\xi)=
    \frac{\sum_m Valid^{basic,(m)}\mathbf{1}\{V^{gross,(m)}\geq1+\xi\}}
    {\sum_m Valid^{basic,(m)}}
    \quad\text{if }\sum_m Valid^{basic,(m)}>0,
\end{equation}
\begin{equation}
    \widehat{\pi}^{support}(\xi)=
    \frac{\sum_m Valid^{support,(m)}\mathbf{1}\{V^{gross,(m)}\geq1+\xi\}}
    {\sum_m Valid^{support,(m)}}
    \quad\text{if }\sum_m Valid^{support,(m)}>0.
\end{equation}
Report $\pi^{threshold|basic}$ and $IS^{support|threshold}$ to separate non-computable thresholds from support violations, and retain $\pi^{support|basic}$ as the aggregate support-valid share. The debt-consistent version additionally requires $fs^{eff,(m)}\geq(1+\xi)c^{gross,(m)}+s^{PB,(m)}$.

\begin{longtable}{p{0.18\textwidth}p{0.36\textwidth}p{0.34\textwidth}}
\caption{Required Monte Carlo output layers}\label{tab:mc_output_layers}\\
\toprule
Layer & Output & Purpose \\
\midrule
\endfirsthead
\toprule
Layer & Output & Purpose \\
\midrule
\endhead
Distribution & $P5$, $P25$, $P50$, $P75$, $P95$ of $V^{gross}$ & Shows downside, median, and upside feasibility \\
Crossing probability & $\widehat{\pi}(\xi)$, $\widehat{\pi}^{basic}(\xi)$, and $\widehat{\pi}^{support}(\xi)$ & Shows whether feasibility survives buffers with all-draw, basic-valid, and support-valid denominators \\
Debt-consistent probability & $\widehat{\pi}^{DC}(\xi)$ & Separates accounting coverage from fiscal sustainability \\
Failure severity & $E[1+\xi-V\mid V<1+\xi]$ & Distinguishes slight failure from severe failure \\
Validity / support chain & $\pi^{basic}$, $\pi^{threshold|basic}$, and $IS^{support|threshold}$ & Separates denominator-invalid draws, non-computable thresholds, and bounded-support violations \\
Driver ranking & Spearman/tornado/Sobol outputs & Identifies what moves $V$ \\
\bottomrule
\end{longtable}

\subsection{Global sensitivity analysis}

A parsimonious first implementation reports Spearman rank correlations
\begin{equation}
    \rho^S_j=Corr_S(\theta_j,V^{gross}).
\end{equation}
A tornado chart varies one input at a time between low and high values. Where inputs are independent, first-order and total-order Sobol indices may be reported \citep{morris1991factorial,saltelli2008global,sobol2001global}:
\begin{equation}
    S_j=\frac{Var_{\theta_j}(E[V\mid\theta_j])}{Var(V)},
    \qquad
    S^T_j=1-\frac{Var_{\theta_{-j}}(E[V\mid\theta_{-j}])}{Var(V)}.
\end{equation}
If the Monte Carlo imposes dependence, Sobol indices are reported only in an independent diagnostic run or with a dependence-aware method.

\subsection{Deterministic stress tests}

The deterministic stress menu spans low AI productivity, slow adoption, weak fiscal capture, high informality, high administrative cost, high leakage, no AI-rent tax, no formalization effect, and no labor augmentation; the executed stress and ablation outputs are recorded in the reproducibility package.

\subsection{Ablation, placebo, and negative-control outputs}

The replication package includes temporal placebo, sectoral negative-control, leave-one-source-out, and channel-ablation outputs.

\subsection{Numerical convergence and reproducibility}

The replication package reports the random seed, draw count, convergence checks for crossing probabilities (its Monte Carlo convergence output), support-valid draw share, code version, and data version. Any table using $V^{net}$ must also report $V^{gross}$.

\section{Calibration Matrix and Assumption Ledger}\label{app:calibration_ledger}

This appendix consolidates the minimum parameter-reporting standard and the ledger of maintained assumptions with their bias directions.

\subsection{Master parameter-assignment and calibration matrix}

The following matrix is the minimum reporting standard for the empirical application. It is designed to make each value reproducible and to prevent the model from hiding fragile assumptions inside a single composite index.

\scriptsize
\begin{longtable}{p{0.14\textwidth}p{0.22\textwidth}p{0.25\textwidth}p{0.20\textwidth}p{0.11\textwidth}}
\caption{Master parameter-assignment and calibration matrix}\label{tab:master_calibration_matrix}\\
\toprule
Input & Definition and unit & Baseline assignment & Uncertainty treatment & Primary source discipline \\
\midrule
\endfirsthead
\toprule
Input & Definition and unit & Baseline assignment & Uncertainty treatment & Primary source discipline \\
\midrule
\endhead
$Y_{i,t}$, $Pop_{i,t}$, $N_{i,p,t}$ & GDP, population, and eligible population for policy $p$ & Observed country-year values; keep nominal and PPP versions separate & Fixed in baseline; measurement-error check if data are interpolated & World Bank PIP/WDI; national accounts \\
$c^{gross,0}_{i,p,t}$, $c^{net,0}_{i,p,t}$ & Annual gross or net policy cost as share of no-AI counterfactual GDP $Y^0_{i,t}$ & Construct from benefit rule $B_{i,p,t}$ and eligible population; subtract existing spending only in the net-cost specification; post-AI denominators are robustness only & Low/central/high benefit rules; no mixing gross and net cost & ASPIRE, PIP, national budgets, \citet{gentilini2020ubi} \\
$\eta_{MUT},\eta_{PBI},\eta_{UBI},\chi_{GMI},B^{PEN},\theta^{target}$ & Benefit-rule parameters, pension benefit level, and GMI targeting-and-behavior loading & Indicative support $[0.25,1.00]$ for poverty-line fractions and gap coverage, including $B^{PEN}/PL$; $\theta^{target}\in[1.0,1.5]$ for labelled GMI robustness; finalized in the official calibration registry (reproducibility package, appendix\_i files) & Low/central/high grid or bounded PERT within declared support & PIP, ASPIRE, and national benefit rules \\
$\phi^{raw}_{s,u}$, $\phi^{Y,nom}_s$, $g^{AI,level}$ & Raw source-unit AI shock, nominal GDP-equivalent annual bridge, and accumulated level shock used against recurring costs & Draw raw support from Table \ref{tab:phi_numerical_closure}; convert through $\kappa^{u\to Y}$ and $\pi^{AI,Y}$; accumulate before feasibility comparison & Bounded PERT on raw support; bridge uncertainty reported separately; triangular/uniform only robustness & \citet{acemoglu2024simple}, \citet{filippucci2024miracle}, \citet{oecd2025g7} \\
$A_{i,t}$ & AI preparedness, normalized $[0,1]$ & IMF AIPI or documented reconstruction from its four dimensions: digital infrastructure, human capital and labour-market policies, innovation/economic integration, regulation/ethics & Fixed observed index plus normalization robustness; avoid using as causal estimate & \citet{imf2024aipi,cazzaniga2024genai} \\
$Gap_{i,t}$ & Digital and skills gap, normalized to $[0,1]$ & Construct only from indicators excluded from AIPI & Fixed observed index plus normalization robustness & ITU and World Bank digital indicators \\
$E^{prod}_{i,t}$, $\varpi$ & Productive AI exposure and occupational-versus-sectoral combination weight & Prefer occupation/sector exposure mapped to employment or value added; for LAC, use productivity-enhancing exposure as baseline rather than total exposure & Beta/PERT distribution around low-central-high exposure; $\varpi\in\{0.25,0.50,0.75\}$ or a declared weight; broad exposure is robustness, not baseline & \citet{worldbankilo2024lacgenai,cazzaniga2024genai} \\
$E^{disp}_{i,t}$ & Displacement or automation exposure & Separate from $E^{prod}$; used in pressure and labour-share modules, not directly as productivity opportunity & Scenario range; never add to $E^{prod}$ & \citet{worldbankilo2024lacgenai,cazzaniga2024genai} \\
$q^{use}_{i,s,t}$, $q^{prod}_{i,s,t}$, $\omega_I$, $\omega_G$ & AI use, productive adoption, and adoption-ceiling parameters; adoption shares lie in $[0,1]$ & If observed, use firm/worker adoption data; if simulated, use logistic module with country ceiling $\bar q_i$ & Logit-normal or bounded PERT for adoption; grid or bounded PERT with declared support for $\omega_I$ and $\omega_G$; enforce $0\le q^{prod}\le\bar q_i\le1$; impose $q^{prod}\le q^{use}$ only under non-decreasing-use robustness & Diffusion literature \citep{bass1969newproduct,rogers2003diffusion}; micro-productivity evidence \citep{noy2023experimental,brynjolfsson2023generative} \\
$\lambda_q$ & Productive-adoption lag parameter & Baseline grid $\{0.15,0.30,0.50\}$; $1.00$ only as instant-adoption stress & Discrete grid & Diffusion theory and adoption-lag logic \\
$\nu_A,\nu_I,\nu_G,q^{use,target}$ & Logistic adoption coefficients and anchored adoption moment & Calibrate the anchored simulated-adoption mode in Appendix \ref{app:adoption_module} & Grid or bounded PERT with declared support & Appendix \ref{app:adoption_module}; diffusion literature; national ICT/firm surveys or World Bank technology adoption surveys for the target moment \\
$S^{frontier}_{s,t}$ & External translation benchmark & Document the frontier-economy set and benchmark year & Fixed within each threshold inversion; alternative external sets as robustness & External frontier sample; Appendix \ref{app:ai_shock_module} \\
$\alpha,\beta,\gamma$ & Elasticities in the translation factor & $\beta=\gamma=1$ baseline; $\alpha=0$ under the no-double-counting baseline ($\alpha=1$ only in full-score robustness; see Appendix \ref{app:ai_shock_module} elasticity grid) & Grid $\{0.5,1,1.5,2\}$ for active elasticities and global sensitivity & Model discipline; no direct causal claim \\
$RST_{i,t}$, $\delta_s,\psi_s,\zeta_s,\lambda_{LS}$ & Reskilling/task-reinstatement capacity in $[0,1]$ plus labor-share automation, augmentation, mitigation, and partial-adjustment parameters & Indicative provisional grids $\delta_s,\psi_s,\zeta_s\in\{0.25,0.50,1.00\}$ and $\lambda_{LS}\in\{0.25,0.50,1.00\}$; enforce $\zeta_sRST\leq\delta_s$ outside augmentation-dominant robustness; finalized in the official calibration registry (reproducibility package, appendix\_i files) & Low/central/high grid for scenario parameters; $RST\in[0,1]$ from declared proxy; bounded PERT only after pilot supports are declared & \citet{acemoglu2018race,acemoglu2019automation}; active labour-market policy spending, training participation (ILOSTAT), or residualized human-capital component of AIPI; calibrated scenario discipline \\
$MFC^{gross}_{i,s,r,t}$ & Gross marginal fiscal capture of AI-induced output before administrative costs, transition costs, and leakage & Construct from channel model or draw from historical gross-capture distribution, but not both in the same run & Channel-mode derived object or historical-mode stochastic object; report mode explicitly & UNU-WIDER GRD; \citet{brollo2024fiscal}; IMF fiscal-space guidance \\
$MFC^{eff}_{i,p,s,r,t}$ & Derived reporting object after subtracting $ac$, $tr$, and $leak$ from fiscal space & Compute only as $MFC^{eff}=fs^{eff}/g^{AI}$ when $g^{AI}>0$; not drawn as an independent primitive & Report for interpretation; do not use for gross historical plausibility checks & Derived from $MFC^{gross}$, $g^{AI}$, $ac$, $tr$, and $leak$ \\
$\tau^{L,fisc,eff},\tau^{K,fisc,eff},\tau^{C,eff},\tau^{AI}$; $\tau^{L,disp},\tau^{K,disp}$ & Fiscal-capture rates for revenue and disposable-income rates for induced consumption & Infer fiscal rates from revenue composition where possible; disposable rates may coincide only in simplified calibration & Grid by regime; cap AI-rent tax at plausible maximum $\bar\tau^{AI}$ & GRD revenue composition; \citet{brollo2024fiscal} \\
$\chi^{Kbase},\chi^{dom},\psi^{shift}$; $\psi^{K,rate}$ & Capital-base chain parameters, each in $[0,1]$; $\psi^{K,rate}$ is the rate-side shifting parameter, active only under the labelled rate-side convention with $\psi^{shift}=0$ & Construct the taxable domestic post-shifting capital base without double counting & Bounded PERT on low/central/high supports & National accounts and profit-shifting literature \\
$mpc^W,mpc^{\Pi},mpc^R,mpc^{ben},m^M,exempt^C,\theta^{R,dom},\tau^{R,disp}$ & Induced-consumption shares and rates, each in $[0,1]$; $mpc^{ben}$ active only in the labelled first-round recycling variant & Calibrate disposable-income pass-through, domestic routing, imports, and exemptions consistently & Bounded PERT or declared grid & Household surveys, national accounts, and VAT rules \\
$\lambda^{AI},\chi^{AI,dom},\psi^{shift,tech},\chi^{AI,compliance},\chi^{AI,enforce}$ & AI-rent share and capturability components in $[0,1]$ & Zero AI-rent channel in the status quo through $\tau^{AI}=0$ unless such an instrument already exists; bounded shares in positive scenarios & Bounded scenario grid or PERT & Fiscal law, international taxation, and Appendix \ref{app:fiscal_conversion_module} \\
$I^L_{i,t}, I^F_{i,t}, I^C_{i,t}$ & Labour, firm, and consumption-base informality & Use source-specific measures where available; do not collapse all frictions into one informal-economy scalar unless data force it & Fixed observed value; marginal informality $I^{j,marg}$ preferred when sectoral AI-shock weights exist, aggregate $I^j$ as declared fallback & ILOSTAT; \citet{worldbank2021informality,medina2018shadow} \\
$\rho_L,\rho_K,\rho_C,\rho_E,\rho_\psi$ & Elasticities for capturability, exemptions, and rate-side profit shifting & No preferred point value without estimation; $\rho_E$ and $\rho_\psi$ active only in their labelled conventions & Grid $\{0.5,1,1.5,2\}$; report sensitivity and Spearman/Sobol diagnostics & Calibrated structural parameters \\
$\epsilon^{ero,j}$ & Behavioral taxable-base erosion elasticities by channel & Zero only in the mechanical upper-bound reading of reform regimes; mandatory companion for $r\geq1$ results, with $\epsilon^{ero,j}\tau^j<1$ & Declared grid or bounded PERT; recorded in the calibration audit trail & Tax avoidance and evasion elasticity literature; Appendix \ref{app:fiscal_conversion_module} \\
$ac^{net,fix},tr^{ann},leak^{fix}$ & Fixed administrative, transition, and residual leakage costs as GDP-normalized ratios & Subtract from fiscal space after the AI shock is converted into revenue & Bounded PERT supports; triangular only as robustness & ASPIRE expenditure concepts; IMF fiscal policy; GRD \\
$ac^{ME},tr^{ME},leak^{ME}$ & Marginal-equivalent deductions per unit of AI-induced output & Enter only through $\widetilde{MFC}^{gross}=MFC^{gross}-ac^{ME}-tr^{ME}-leak^{ME}$ & Bounded PERT or zero in baseline if costs are fixed & IMF fiscal policy; implementation stress tests \\
$\Omega^{res},s^{PB}$ & Residual stress penalty and primary-balance correction & $\Omega^{res}=0$ in the baseline; derive $s^{PB}$ from baseline and debt-stabilizing primary balances & Positive stress grid for $\Omega^{res}$; WEO/DSA scenarios for $s^{PB}$ & IMF WEO and debt-sustainability analysis \\
$\xi$, $\bar\pi$ & Prudential buffer and probability label & $\xi\in\{0,0.05,0.10,0.25,0.50\}$ and $\bar\pi\in\{0.50,0.75,0.90\}$ & Fixed reporting grid, not estimated parameter & Robust Decision Making logic \citep{lempert2003shaping,rand2013rdm} \\
\bottomrule
\end{longtable}
\normalsize

\subsection{Data construction rules}

Observed variables should be constructed as follows:
\begin{itemize}
    \item Fiscal capacity: tax revenue, non-tax revenue, grants, social contributions, and revenue composition from the UNU-WIDER Government Revenue Dataset \citep{unu2025grd}; use central/general government consistently.
    \item Informality: informal employment as a share of total employment from ILOSTAT where available, supplemented by shadow-economy or informality studies only when official labour measures are missing \citep{ilostat2026informal,medina2018shadow,worldbank2021informality}.
    \item Social protection: coverage, adequacy, expenditure, beneficiary information, and program categories from ASPIRE where available \citep{worldbank2026aspire}.
    \item Poverty lines and poverty gaps: World Bank Poverty and Inequality Platform or national poverty lines, with PPP and local-currency concepts kept separate \citep{worldbank2026pip}.
    \item AI preparedness: IMF AI Preparedness Index and its four components; do not reinterpret the index as invention leadership \citep{cazzaniga2024genai}.
    \item AI exposure: occupation- or sector-level exposure harmonized to national employment or value-added shares. Productive exposure and displacement exposure must be reported as separate variables.
    \item Digital and skills gap: construct $Gap_{i,t}$ from ITU and World Bank digital indicators excluded from AIPI, and document the normalization to $[0,1]$.
    \item Primary-balance correction: obtain $PB^0_{i,t}$ and $PB^{stab}_{i,t}$ for $s^{PB}_{i,t}$ from the IMF WEO and debt-sustainability analysis.
    \item External translation frontier: document the benchmark economies and year used to construct $S^{frontier}_{s,t}$, which remains fixed within each threshold inversion, and state that the benchmark set matches the population underlying the scenario shock $\phi_s$ (population consistency).
\end{itemize}

\subsection{Maintained economic assumptions and bias directions}

The accounting frontier rests on maintained assumptions stated throughout the paper. They are consolidated here with their bias direction: conservative assumptions understate feasibility, optimistic ones overstate it.

\begin{longtable}{p{0.43\textwidth}p{0.22\textwidth}p{0.25\textwidth}}
\caption{Maintained economic assumptions and bias directions}\label{tab:maintained_assumptions_bias}\\
\toprule
Assumption & Stated in & Direction \\
\midrule
\endfirsthead
\toprule
Assumption & Stated in & Direction \\
\midrule
\endhead
One-directional chain; no feedback from protection to productivity & Section \ref{sec:conceptual} & Ambiguous \\
Endpoint level gain persists as long as the policy; composition may degrade through rent dissipation & Section \ref{subsec:phi_ladder_main} & Optimistic for long-run capture \\
Non-market and zero-price AI gains outside the fiscal base; framework measures taxable income, not welfare & Sections \ref{subsec:phi_ladder_main} and \ref{sec:claims} & Conservative \\
\(\kappa^{TFP\to Y}\leq1\) excludes induced capital deepening & Section \ref{subsec:phi_ladder_main} & Conservative \\
No general equilibrium: no wage spillovers, first-round consumption only & Sections \ref{sec:conceptual} and \ref{sec:fiscal_conversion} & Conservative \\
\(S^{frontier}\) benchmark matches the population of \(\phi_s\) & Section \ref{subsec:ndc_translation_main} & Conservative wedge if mismatched \\
No endogenous imitation diffusion; adoption paths are scenario paths & Appendix \ref{app:adoption_module} & Conservative in the medium term \\
Informality fixed in baseline; formalization and informalization only as labelled scenarios, reported jointly & Appendix \ref{app:informality_module} & Ambiguous \\
One adoption stock gates automation and augmentation; \(E^{auto}\) anchored to full-automation risk is a lower bound on displacement & Section \ref{sec:labor_share} & Optimistic for the labor tax base \\
\(\chi^{Kbase}\) calibrated on average non-labor-residual composition and applied to the AI increment & Section \ref{sec:labor_share} & Conservative \\
LAC anchoring uses productivity-enhancing exposure only; automation-driven output gains excluded from \(E^{prod}\) & Section \ref{subsec:prod_disp_exposure_main} & Conservative \\
Statutory-accrual incidence; no incidence shifting & Section \ref{sec:fiscal_conversion} & Ambiguous \\
Reform-regime gains are mechanical upper bounds; behavioral-erosion companion mandatory for \(r\geq1\) & Section \ref{sec:fiscal_conversion} & Optimistic for \(r\geq1\) \\
Historical capture prior is composition-weighted and policy-inclusive (buoyancy, not elasticity) & Section \ref{sec:historical_plausibility} & Ambiguous; generous for $r=0$ plausibility labels \\
Recipient behavior fixed; GMI priced at ideal targeting absent \(\theta^{target}\); no transfer recycling & Section \ref{sec:policy_costs} & GMI-optimistic; recycling omission conservative \\
Feasibility is an endpoint condition; crossing year is not a launch year & Section \ref{sec:thresholds} & Naive launch reading optimistic \\
\(PB^{stab}\) static at baseline growth & Section \ref{sec:thresholds} & Conservative \\
\bottomrule
\end{longtable}

\begin{thebibliography}{99}

\bibitem[Acemoglu(2024)]{acemoglu2024simple}
Acemoglu, D. (2024).
\newblock The simple macroeconomics of AI.
\newblock \emph{NBER Working Paper No. 32487}. DOI: 10.3386/w32487.

\bibitem[Acemoglu and Restrepo(2018)]{acemoglu2018race}
Acemoglu, D. and Restrepo, P. (2018).
\newblock The race between man and machine: Implications of technology for growth, factor shares, and employment.
\newblock \emph{American Economic Review}, 108(6), 1488--1542. DOI: 10.1257/aer.20160696.

\bibitem[Acemoglu and Restrepo(2019)]{acemoglu2019automation}
Acemoglu, D. and Restrepo, P. (2019).
\newblock Automation and new tasks: How technology displaces and reinstates labor.
\newblock \emph{Journal of Economic Perspectives}, 33(2), 3--30. DOI: 10.1257/jep.33.2.3.


\bibitem[Athey and Imbens(2017)]{athey2017state}
Athey, S. and Imbens, G. W. (2017).
\newblock The state of applied econometrics: Causality and policy evaluation.
\newblock \emph{Journal of Economic Perspectives}, 31(2), 3--32. DOI: 10.1257/jep.31.2.3.

\bibitem[Bass(1969)]{bass1969newproduct}
Bass, F. M. (1969).
\newblock A new product growth for model consumer durables.
\newblock \emph{Management Science}, 15(5), 215--227. DOI: 10.1287/mnsc.15.5.215.

\bibitem[Bertrand, Duflo and Mullainathan(2004)]{bertrand2004trust}
Bertrand, M., Duflo, E., and Mullainathan, S. (2004).
\newblock How much should we trust differences-in-differences estimates?
\newblock \emph{Quarterly Journal of Economics}, 119(1), 249--275. DOI: 10.1162/003355304772839588.

\bibitem[Brollo et al.(2024)]{brollo2024fiscal}
Brollo, F., Dabla-Norris, E., de Mooij, R., Garcia-Macia, D., Hanappi, T., Liu, L., and Nguyen, A. D. M. (2024).
\newblock Broadening the gains from generative AI: The role of fiscal policies.
\newblock \emph{IMF Staff Discussion Note 2024/002}. DOI: 10.5089/9798400277177.006.

\bibitem[Browne and Immervoll(2017)]{oecd2017basicincome}
Browne, J. and Immervoll, H. (2017).
\newblock Mechanics of replacing benefit systems with a basic income: Comparative results from a microsimulation approach.
\newblock \emph{OECD Social, Employment and Migration Working Papers No. 201}. DOI: 10.1787/ec38a279-en.

\bibitem[Brynjolfsson, Li and Raymond(2023)]{brynjolfsson2023generative}
Brynjolfsson, E., Li, D., and Raymond, L. R. (2023).
\newblock Generative AI at work.
\newblock \emph{NBER Working Paper No. 31161}. DOI: 10.3386/w31161.

\bibitem[Cazzaniga et al.(2024)]{cazzaniga2024genai}
Cazzaniga, M., Jaumotte, F., Li, L., Melina, G., Panton, A. J., Pizzinelli, C., Rockall, E. J., and Tavares, M. M. (2024).
\newblock Gen-AI: Artificial intelligence and the future of work.
\newblock \emph{IMF Staff Discussion Note 2024/001}. DOI: 10.5089/9798400262548.006.



\bibitem[DANE(2026)]{dane2026informalidad}
Departamento Administrativo Nacional de Estadistica. (2026).
\newblock Empleo informal y seguridad social: informacion enero--marzo 2026.
\newblock Updated May 13, 2026.

\bibitem[Filippucci et al.(2024)]{filippucci2024miracle}
Filippucci, F., Gal, P., and Schief, M. (2024).
\newblock Miracle or myth? Assessing the macroeconomic productivity gains from artificial intelligence.
\newblock \emph{OECD Artificial Intelligence Papers No. 29}. DOI: 10.1787/b524a072-en.


\bibitem[Filippucci et al.(2025)]{oecd2025g7}
Filippucci, F., Gal, P., Laengle, K., and Schief, M. (2025).
\newblock Macroeconomic productivity gains from artificial intelligence in G7 economies.
\newblock \emph{OECD Artificial Intelligence Papers No. 41}. DOI: 10.1787/a5319ab5-en.


\bibitem[Gaspar et al.(2016)]{gaspar2016tax}
Gaspar, V., Jaramillo, L., and Wingender, P. (2016).
\newblock Tax capacity and growth: Is there a tipping point?
\newblock \emph{IMF Working Paper 16/234}. DOI: 10.5089/9781475558142.001.

\bibitem[Gentilini et al.(2020)]{gentilini2020ubi}
Gentilini, U., Grosh, M., Rigolini, J., and Yemtsov, R. (2020).
\newblock \emph{Exploring Universal Basic Income: A Guide to Navigating Concepts, Evidence, and Practices}.
\newblock World Bank. DOI: 10.1596/978-1-4648-1458-7.


\bibitem[Gollin(2002)]{gollin2002getting}
Gollin, D. (2002).
\newblock Getting income shares right.
\newblock \emph{Journal of Political Economy}, 110(2), 458--474. DOI: 10.1086/338747.

\bibitem[ILOSTAT(2026)]{ilostat2026informal}
ILOSTAT. (2026).
\newblock Informal employment rate.
\newblock International Labour Organization statistical snapshot and database, accessed June 2026.

\bibitem[INEGI(2024)]{inegi2024enoe}
Instituto Nacional de Estadistica y Geografia. (2024).
\newblock Encuesta Nacional de Ocupacion y Empleo, segundo trimestre de 2024.
\newblock Press release 504/24, September 2, 2024.


\bibitem[INE Chile(2025)]{inechile2025ene}
Instituto Nacional de Estadisticas de Chile. (2025).
\newblock Separata tecnica anual de la Encuesta Nacional de Empleo 2024.
\newblock Press release dated March 20, 2025.

\bibitem[International Monetary Fund(2016)]{imf2016fiscalspace}
International Monetary Fund. (2016).
\newblock Assessing fiscal space: An initial consistent set of considerations.
\newblock \emph{IMF Policy Paper}, June 23, 2016. DOI: 10.5089/9781498345583.007.

\bibitem[International Monetary Fund(2024)]{imf2024aipi}
International Monetary Fund. (2024).
\newblock AI Preparedness Index (AIPI).
\newblock IMF DataMapper dataset and documentation.

\bibitem[Karabarbounis and Neiman(2014)]{karabarbounis2014global}
Karabarbounis, L. and Neiman, B. (2014).
\newblock The global decline of the labor share.
\newblock \emph{Quarterly Journal of Economics}, 129(1), 61--103. DOI: 10.1093/qje/qjt032.

\bibitem[Lempert, Popper and Bankes(2003)]{lempert2003shaping}
Lempert, R. J., Popper, S. W., and Bankes, S. C. (2003).
\newblock \emph{Shaping the Next One Hundred Years: New Methods for Quantitative, Long-Term Policy Analysis}.
\newblock RAND Corporation, MR-1626-RPC.



\bibitem[Lipsitch, Tchetgen Tchetgen and Cohen(2010)]{lipsitch2010negative}
Lipsitch, M., Tchetgen Tchetgen, E., and Cohen, T. (2010).
\newblock Negative controls: A tool for detecting confounding and bias in observational studies.
\newblock \emph{Epidemiology}, 21(3), 383--388. DOI: 10.1097/EDE.0b013e3181d61eeb.

\bibitem[McKay, Beckman and Conover(1979)]{mckay1979lhs}
McKay, M. D., Beckman, R. J., and Conover, W. J. (1979).
\newblock A comparison of three methods for selecting values of input variables in the analysis of output from a computer code.
\newblock \emph{Technometrics}, 21(2), 239--245. DOI: 10.2307/1268522.

\bibitem[Medina and Schneider(2018)]{medina2018shadow}
Medina, L. and Schneider, F. (2018).
\newblock Shadow economies around the world: What did we learn over the last 20 years?
\newblock \emph{IMF Working Paper 18/17}. DOI: 10.5089/9781484338636.001.

\bibitem[Morris(1991)]{morris1991factorial}
Morris, M. D. (1991).
\newblock Factorial sampling plans for preliminary computational experiments.
\newblock \emph{Technometrics}, 33(2), 161--174. DOI: 10.1080/00401706.1991.10484804.

\bibitem[Noy and Zhang(2023)]{noy2023experimental}
Noy, S. and Zhang, W. (2023).
\newblock Experimental evidence on the productivity effects of generative artificial intelligence.
\newblock \emph{Science}, 381(6654), 187--192. DOI: 10.1126/science.adh2586.

\bibitem[OECD(2024)]{oecd2024impactai}
OECD. (2024).
\newblock The impact of artificial intelligence on productivity, distribution and growth.
\newblock \emph{OECD Artificial Intelligence Papers No. 15}. DOI: 10.1787/8d900037-en.

\bibitem[OECD(2025)]{oecd2025peru}
OECD. (2025).
\newblock \emph{OECD Economic Surveys: Peru 2025}.
\newblock OECD Publishing.

\bibitem[OECD et al.(2025)]{oecd2025revenue}
OECD, ECLAC, CIAT, and IDB. (2025).
\newblock \emph{Revenue Statistics in Latin America and the Caribbean 2025}.
\newblock OECD Publishing. DOI: 10.1787/7594fbdd-en.

\bibitem[RAND Corporation(2013)]{rand2013rdm}
RAND Corporation. (2013).
\newblock Making good decisions without predictions: Robust Decision Making for planning under deep uncertainty.
\newblock RAND Research Highlights, RB-9701.


\bibitem[Rogers(2003)]{rogers2003diffusion}
Rogers, E. M. (2003).
\newblock \emph{Diffusion of Innovations}, 5th ed.
\newblock Free Press.



\bibitem[Saltelli et al.(2008)]{saltelli2008global}
Saltelli, A., Ratto, M., Andres, T., Campolongo, F., Cariboni, J., Gatelli, D., Saisana, M., and Tarantola, S. (2008).
\newblock \emph{Global Sensitivity Analysis: The Primer}.
\newblock Wiley.

\bibitem[Sobol(2001)]{sobol2001global}
Sobol, I. M. (2001).
\newblock Global sensitivity indices for nonlinear mathematical models and their Monte Carlo estimates.
\newblock \emph{Mathematics and Computers in Simulation}, 55(1--3), 271--280. DOI: 10.1016/S0378-4754(00)00270-6.

\bibitem[UNU-WIDER(2025)]{unu2025grd}
UNU-WIDER. (2025).
\newblock Government Revenue Dataset, Version 2025.
\newblock DOI: 10.35188/UNU-WIDER/GRD-2025.

\bibitem[World Bank(2021)]{worldbank2021informality}
World Bank. (2021).
\newblock \emph{The Long Shadow of Informality: Challenges and Policies}.
\newblock World Bank, Washington, DC. DOI: 10.1596/978-1-4648-1753-3.


\bibitem[World Bank(2026a)]{worldbank2026aspire}
World Bank. (2026a).
\newblock ASPIRE: The Atlas of Social Protection Indicators of Resilience and Equity.
\newblock Social protection and labor indicators, accessed June 2026.

\bibitem[World Bank(2026b)]{worldbank2026pip}
World Bank. (2026b).
\newblock Poverty and Inequality Platform.
\newblock Poverty, inequality, and poverty-line data, accessed June 2026.

\bibitem[World Bank(2026c)]{worldbank2026wdi}
World Bank. (2026c).
\newblock World Development Indicators.
\newblock Country data for GDP, GDP per capita, population, and macroeconomic indicators, accessed June 2026.

\bibitem[World Bank(2026d)]{worldbank2026humancapital}
World Bank. (2026d).
\newblock Human Capital country indicators: Peru.
\newblock Adult informal employment and related labor indicators, accessed June 2026.

\bibitem[World Bank and ILO(2024)]{worldbankilo2024lacgenai}
World Bank and International Labour Organization. (2024).
\newblock Generative AI and jobs in Latin America and the Caribbean: Is the digital divide a buffer or bottleneck?
\newblock World Bank Group and International Labour Organization.

\end{thebibliography}

\end{document}
